% Vector Product and its Applications — Further Mathematics Upper Sixth — Cours signé OnBuch+
% Official MINESEC syllabus. Compiler avec : tectonic FMA14-vector-product-and-its-applications.tex
\newcommand{\DOCMATIERE}{Further Mathematics}
\newcommand{\DOCNIVEAU}{Upper Sixth}
\newcommand{\DOCMODULE}{Upper Sixth — Further mathematics}
\newcommand{\DOCLECON}{14}
\newcommand{\DOCTITRE}{Vector Product and its Applications}
% OnBuch+ — préambule « cours signés » ENRICHI (compilé avec tectonic / XeLaTeX)
% Système visuel « Pop » d'OnBuch+ : fond crème (jamais blanc), bordures encre
% épaisses, ombres « dures » décalées, titres Archivo Black en capitales.
% Version MATHÉMATIQUES : graphes (pgfplots/tikz), boîtes pédagogiques
% (définition, propriété, méthode, attention, le savais-tu, exercice résolu,
% exercice, TP, bilan), page de garde et signature OnBuch+ ancrée en bas.
%
% Macros attendues dans main.tex AVANT \input{preamble} :
%   \DOCMATIERE  \DOCNIVEAU  \DOCMODULE  \DOCLECON (numéro)  \DOCTITRE
\documentclass[11pt]{article}

\usepackage{fontspec}
\usepackage[english]{babel}
\usepackage[a4paper,margin=2cm,top=3.4cm,bottom=2.8cm,headheight=1.4cm,footskip=1.3cm]{geometry}
\usepackage{amsmath,amssymb,amsfonts}
\usepackage{mathtools} % \xmapsto, \coloneqq... souvent utilisés par les modèles
\usepackage{cancel} % simplifications barrées (bilans de forces)
% \abs{z} (module, valeur absolue) et \norm{u} (norme), utilisés par les modèles
\providecommand{\abs}{}\let\abs\relax\DeclarePairedDelimiter{\abs}{\lvert}{\rvert}
\providecommand{\norm}{}\let\norm\relax\DeclarePairedDelimiter{\norm}{\lVert}{\rVert}
\usepackage[table]{xcolor} % [table] : fournit aussi \rowcolors (alternance de lignes)
\usepackage{pagecolor}
\usepackage{tikz}
\usetikzlibrary{arrows.meta,positioning,calc,shapes.geometric,shapes.misc,shapes.arrows,decorations.pathmorphing,decorations.pathreplacing,decorations.markings,patterns,shadows,mindmap,trees,fit,backgrounds,matrix,intersections,angles,quotes}
\usepackage{pgfplots}
\usepackage[version=4]{mhchem} % équations nucléaires \ce{^{235}_{92}U}
\usepackage[european,siunitx]{circuitikz} % schémas électriques
\pgfplotsset{compat=1.18}
\usepgfplotslibrary{fillbetween,groupplots} % aires entre courbes (calcul intégral)
\usepackage{enumitem}
\usepackage{fancyhdr}
% [most] charge toutes les bibliothèques tcolorbox (listings, minted, raster,
% documentation...) et gonfle inutilement la mémoire TeX sur un document long
% (~300 boîtes) : on ne charge que ce qui est réellement utilisé.
\usepackage[skins,breakable]{tcolorbox}
\usepackage{graphicx}
\usepackage{colortbl}
\usepackage{booktabs}
\usepackage{tabularx}
\usepackage{diagbox} % case d'en-tête diagonale des tableaux à double entrée
% Colonnes extensibles centrée / gauche / droite (C, L, R), souvent utilisées par les modèles
\newcolumntype{C}{>{\centering\arraybackslash}X}
\newcolumntype{L}{>{\raggedright\arraybackslash}X}
\newcolumntype{R}{>{\raggedleft\arraybackslash}X}
\usepackage{array}
\usepackage{multicol}
\usepackage{siunitx}
% parse-numbers=false : accepte aussi \SI{2,5 \times 10^{-3}}{...}
\sisetup{locale=UK,output-decimal-marker={.},group-minimum-digits=4,parse-numbers=false}
\DeclareSIUnit\ohm{\ensuremath{\symup{\Omega}}} % U+2126 absent de la police
\DeclareSIUnit\division{div} % réglages d'oscilloscope (ms/div, V/div)
\DeclareSIUnit\tr{tr} % tours (vitesse de rotation, ex. \SI{1500}{\tr\per\minute})
\usepackage{qrcode}
% \degree : commande fréquemment utilisée par les modèles pour un angle ou une
% température en dehors de \SI{}{} (non fournie par les paquets chargés).
\providecommand{\degree}{^{\circ}}
% \qed (fin de démonstration), utilisé par les modèles sans amsthm
\providecommand{\qed}{\hfill$\square$}
% \barre : double barre des valeurs interdites dans un tableau de variations (tkz-tab)
\providecommand{\barre}{\ensuremath{\|}}
% environnement proof (amsthm non chargé) utilisé par les modèles
\ifdefined\proof\else\newenvironment{proof}[1][Proof]{\par\noindent\textit{#1.}\ }{\qed\par}\fi
\usepackage{lastpage}
\usepackage{hyperref}
\hypersetup{hidelinks}

% ============================================================
% Palette « Pop » OnBuch+ — mêmes hex que la classe P (plus_ui.dart)
% ============================================================
\definecolor{poporange}{HTML}{F59321}
\definecolor{popdark}{HTML}{E07A0C}
\definecolor{popcream}{HTML}{FFF6E8}
\definecolor{popink}{HTML}{1C1714}
\definecolor{poppurple}{HTML}{7A5AE0}
\definecolor{popgreen}{HTML}{2E9E6A}
\definecolor{poppink}{HTML}{E0457B}
\definecolor{popblue}{HTML}{2F7DE1}
\definecolor{popgold}{HTML}{C98A12}
\definecolor{popmuted}{HTML}{8A7F76}
\definecolor{popline}{HTML}{EADFD2}
\definecolor{popgray}{HTML}{8A7F76}
\colorlet{popgrey}{popgray}
% Alias tolérants (le modèle nomme parfois les couleurs différemment)
\colorlet{popwhite}{white}
\colorlet{popblack}{popink}
\colorlet{popred}{poppink}
\colorlet{popyellow}{popgold}
% Teintes claires pour fonds de tableaux / zones
\colorlet{popblueL}{popblue!10!white}
\colorlet{poporangeL}{poporange!14!white}
\colorlet{popgreenL}{popgreen!12!white}
\colorlet{poppurpleL}{poppurple!12!white}
\colorlet{poppinkL}{poppink!12!white}
\colorlet{popgoldL}{popgold!14!white}

\pagecolor{popcream}

% ============================================================
% Typographie
% ============================================================
\defaultfontfeatures{Ligatures=TeX}
\setmainfont{PlusJakartaSans-Medium.ttf}[Path=../../../fonts/,BoldFont=PlusJakartaSans-Bold.ttf]
% Maths en sans-serif (Fira Math) avec chiffres et lettres droites de Plus
% Jakarta, pour que les formules se fondent dans le texte.
\usepackage[math-style=ISO,bold-style=ISO]{unicode-math}
\setmathfont{FiraMath-Regular.otf}[Path=../../../fonts/]
\setmathfont{PlusJakartaSans-Medium.ttf}[Path=../../../fonts/,range={up/num,up/{Latin,latin},\mathup}]
% (icomma retiré : décimales avec point en anglais)
% Caractères mathématiques Unicode absents des polices de texte (Plus Jakarta,
% Archivo Black) : rendus par leur équivalent mathématique, en texte comme en
% formule — sinon ils s'affichent en carré vide (ex. « Arithmétique dans ℤ »).
\usepackage{newunicodechar}
\newunicodechar{ℝ}{\ensuremath{\mathbb{R}}}
\newunicodechar{ℤ}{\ensuremath{\mathbb{Z}}}
\newunicodechar{ℂ}{\ensuremath{\mathbb{C}}}
\newunicodechar{ℕ}{\ensuremath{\mathbb{N}}}
\newunicodechar{ℚ}{\ensuremath{\mathbb{Q}}}
\newunicodechar{𝔻}{\ensuremath{\mathbb{D}}}
\newunicodechar{α}{\ensuremath{\alpha}}
\newunicodechar{β}{\ensuremath{\beta}}
\newunicodechar{γ}{\ensuremath{\gamma}}
\newunicodechar{δ}{\ensuremath{\delta}}
\newunicodechar{ε}{\ensuremath{\varepsilon}}
\newunicodechar{θ}{\ensuremath{\theta}}
\newunicodechar{λ}{\ensuremath{\lambda}}
\newunicodechar{μ}{\ensuremath{\mu}}
\newunicodechar{σ}{\ensuremath{\sigma}}
\newunicodechar{τ}{\ensuremath{\tau}}
\newunicodechar{φ}{\ensuremath{\varphi}}
\newunicodechar{ψ}{\ensuremath{\psi}}
\newunicodechar{ω}{\ensuremath{\omega}}
\newunicodechar{Σ}{\ensuremath{\Sigma}}
% majuscules produites par \MakeUppercase dans les titres (α-aminés -> Α)
\newunicodechar{Α}{\ensuremath{\alpha}}
\newunicodechar{Β}{\ensuremath{\beta}}
\newunicodechar{Γ}{\ensuremath{\Gamma}}
\newunicodechar{Δ}{\ensuremath{\Delta}}
\newunicodechar{Ε}{\ensuremath{\varepsilon}}
\newunicodechar{Θ}{\ensuremath{\Theta}}
\newunicodechar{Λ}{\ensuremath{\Lambda}}
\newunicodechar{Μ}{\ensuremath{\mu}}
\newunicodechar{Τ}{\ensuremath{\tau}}
\newunicodechar{Φ}{\ensuremath{\Phi}}
\newunicodechar{Ψ}{\ensuremath{\Psi}}
\newunicodechar{Ω}{\ensuremath{\Omega}}
\newunicodechar{↦}{\ensuremath{\mapsto}}
\newunicodechar{∈}{\ensuremath{\in}}
\newunicodechar{∉}{\ensuremath{\notin}}
\newunicodechar{∧}{\ensuremath{\wedge}}
\newunicodechar{⇔}{\ensuremath{\Leftrightarrow}}
\newunicodechar{⟺}{\ensuremath{\Longleftrightarrow}}
\newunicodechar{⇒}{\ensuremath{\Rightarrow}}
\newunicodechar{⟹}{\ensuremath{\Longrightarrow}}
\newunicodechar{∘}{\ensuremath{\circ}}
\newunicodechar{∪}{\ensuremath{\cup}}
\newunicodechar{∩}{\ensuremath{\cap}}
\newunicodechar{≡}{\ensuremath{\equiv}}
\newunicodechar{⊂}{\ensuremath{\subset}}
\newunicodechar{⊥}{\ensuremath{\perp}}
\newunicodechar{∃}{\ensuremath{\exists}}
\newunicodechar{∀}{\ensuremath{\forall}}
\newunicodechar{∛}{\ensuremath{\sqrt[3]{\,}}}
\newunicodechar{∜}{\ensuremath{\sqrt[4]{\,}}}
\newunicodechar{⃗}{\ensuremath{{}^{\to}}}
\newunicodechar{ⁿ}{\ifmmode^{n}\else\textsuperscript{\ensuremath{n}}\fi}
\newunicodechar{ˣ}{\ifmmode^{x}\else\textsuperscript{\ensuremath{x}}\fi}
\newunicodechar{ᵇ}{\ifmmode^{b}\else\textsuperscript{\ensuremath{b}}\fi}
\newunicodechar{ʳ}{\ifmmode^{r}\else\textsuperscript{\ensuremath{r}}\fi}
\newunicodechar{ᵘ}{\ifmmode^{u}\else\textsuperscript{\ensuremath{u}}\fi}
\newunicodechar{ʸ}{\ifmmode^{y}\else\textsuperscript{\ensuremath{y}}\fi}
\newunicodechar{ᵃ}{\ifmmode^{a}\else\textsuperscript{\ensuremath{a}}\fi}
\newunicodechar{ᵉ}{\ifmmode^{e}\else\textsuperscript{\ensuremath{e}}\fi}
\newunicodechar{ᵏ}{\ifmmode^{k}\else\textsuperscript{\ensuremath{k}}\fi}
\newunicodechar{ᵖ}{\ifmmode^{p}\else\textsuperscript{\ensuremath{p}}\fi}
\newunicodechar{ᴺ}{\ifmmode^{N}\else\textsuperscript{\ensuremath{N}}\fi}
\newunicodechar{ᶜ}{\ifmmode^{c}\else\textsuperscript{\ensuremath{c}}\fi}
\newunicodechar{ʰ}{\ifmmode^{h}\else\textsuperscript{\ensuremath{h}}\fi}
\newunicodechar{ᵠ}{\ifmmode^{q}\else\textsuperscript{\ensuremath{q}}\fi}
\newunicodechar{ₙ}{\ifmmode_{n}\else\textsubscript{\ensuremath{n}}\fi}
\newunicodechar{ₐ}{\ifmmode_{a}\else\textsubscript{\ensuremath{a}}\fi}
\newunicodechar{ᵢ}{\ifmmode_{i}\else\textsubscript{\ensuremath{i}}\fi}
\newunicodechar{ᵣ}{\ifmmode_{r}\else\textsubscript{\ensuremath{r}}\fi}
\newunicodechar{ₖ}{\ifmmode_{k}\else\textsubscript{\ensuremath{k}}\fi}
\newunicodechar{ᵦ}{\ifmmode_{\beta}\else\textsubscript{\ensuremath{\beta}}\fi}
\newunicodechar{ₓ}{\ifmmode_{x}\else\textsubscript{\ensuremath{x}}\fi}
\newfontfamily\popdisplay{ArchivoBlack-Regular.ttf}[Path=../../../fonts/]
\newfontfamily\poplogo{SpaceGrotesk-Bold.ttf}[Path=../../../fonts/]
\newfontfamily\popsemi{PlusJakartaSans-SemiBold.ttf}[Path=../../../fonts/]
\newfontfamily\popxbold{PlusJakartaSans-ExtraBold.ttf}[Path=../../../fonts/]
\newfontfamily\popxboldls{PlusJakartaSans-ExtraBold.ttf}[Path=../../../fonts/,LetterSpace=3.0]

\setlist[enumerate]{leftmargin=1.7em,itemsep=5pt,topsep=4pt}
\setlist[itemize]{leftmargin=1.7em,itemsep=4pt,topsep=4pt,label={\color{poporange}\textbullet}}
\setlength{\parindent}{0pt}
\setlength{\parskip}{5pt}
\renewcommand{\arraystretch}{1.25}

% pgfplots : style Pop par défaut
\pgfplotsset{
  every axis/.append style={axis line style={popink,line width=1pt},tick style={popink},
    grid style={popline,line width=0.5pt},label style={font=\small},tick label style={font=\footnotesize}},
  popcurve/.style={poporange,line width=1.8pt,no markers,smooth},
}
% Même style disponible dans un \draw TikZ simple (hors pgfplots)
\tikzset{popcurve/.style={poporange,line width=1.8pt}}
\tikzset{popgrid/.style={popline,very thin}}
\colorlet{popcurve}{poporange} % le nom du style est parfois utilisé comme couleur

% ============================================================
% Pill / squiggle
% ============================================================
\newcommand{\poppill}[3]{%
  \tikz[baseline=-0.55ex]\node[draw=popink,line width=1.2pt,rounded corners=2.6pt,%
    fill=#2,inner xsep=6.5pt,inner ysep=3pt]{%
    \popxboldls\fontsize{7}{7}\selectfont\color{#3}\MakeUppercase{#1}};%
}
\newcommand{\popsquiggle}[1]{%
  \tikz[baseline=-0.4ex]{
    \draw[#1,line width=2.6pt,line cap=round]
      plot [smooth] coordinates {(0,0.09) (0.34,0.01) (0.68,0.21) (1.02,0.01) (1.36,0.10)};
  }%
}
% Mot-clé surligné (marqueur orange sous le texte)
\newcommand{\cle}[1]{{\popxbold\color{popdark}#1}}

% ============================================================
% En-tête / pied
% ============================================================
\pagestyle{fancy}
\fancyhf{}
\renewcommand{\headrulewidth}{0pt}
\renewcommand{\footrulewidth}{0pt}
\lhead{\raisebox{-0.15ex}{\poplogo\fontsize{13}{13}\selectfont\color{popink}OnBuch\color{poporange}+}}
\rhead{\poppill{\DOCMATIERE\ \textbullet\ \DOCNIVEAU\ \textbullet\ Chapter \DOCLECON}{white}{popink}}
\lfoot{\popxbold\fontsize{7.6}{7.6}\selectfont\color{popmuted}\MakeUppercase{Signed course OnBuch+}\ \textbullet\ {\popsemi\DOCTITRE}}
\rfoot{\tikz[baseline=-0.6ex]\node[rounded corners=8.5pt,fill=popink,minimum height=17pt,minimum width=17pt,inner xsep=5pt,inner sep=0]{\popxbold\fontsize{8}{8}\selectfont\color{popcream}\thepage\,/\,\pageref*{LastPage}};}

% ============================================================
% Titres
% ============================================================
\newcommand{\chaptitre}[2]{%
  \begin{tcolorbox}[enhanced,colback=poporange,colframe=popink,boxrule=2.6pt,arc=11pt,
      drop shadow=popink,left=15pt,right=15pt,top=12pt,bottom=11pt,before skip=3pt,after skip=18pt]
    \poppill{#2}{popink}{popcream}\par\vspace{9pt}
    {\raggedright\hyphenpenalty=10000\exhyphenpenalty=10000\popdisplay\fontsize{22}{24}\selectfont\color{popcream}\MakeUppercase{#1}\par}
  \end{tcolorbox}
}
% Section numérotée : pastille orange avec le numéro + titre Archivo
\newcounter{popsec}
\newcommand{\coursec}[1]{%
  \stepcounter{popsec}\setcounter{popsub}{0}%
  \par\vspace{16pt}\noindent
  \begin{minipage}{\linewidth}
  \tikz[baseline=-0.7ex]\node[draw=popink,line width=1.6pt,fill=poporange,rounded corners=4pt,minimum size=22pt,inner sep=0]{\popdisplay\fontsize{12}{12}\selectfont\color{popcream}\thepopsec};%
  \hspace{8pt}{\popdisplay\fontsize{15}{17}\selectfont\color{popink}\MakeUppercase{#1}}\par
  \vspace{4pt}\noindent{\color{poporange}\rule{36pt}{3.6pt}}
  \end{minipage}\par\nopagebreak\vspace{8pt}
}
\newcounter{popsub}
\newcommand{\courssub}[1]{%
  \stepcounter{popsub}%
  \par\vspace{9pt}\noindent{\popxbold\fontsize{11.5}{13}\selectfont\color{popink}{\color{poporange}\thepopsec.\thepopsub}\hspace{6pt}#1}\par\nopagebreak\vspace{4pt}
}

% ============================================================
% Boîtes pédagogiques — même recette : fond blanc, bordure épaisse colorée,
% ombre dure encre, badge plein. Titre optionnel [#1].
% ============================================================
\tcbset{popbase/.style={breakable,enhanced,colback=white,boxrule=2.3pt,arc=9pt,
  shadow={3pt}{-3pt}{0pt}{popink},left=14pt,right=14pt,top=11pt,bottom=11pt,before skip=11pt,after skip=15pt,
  fontupper=\popsemi\fontsize{10.6}{14.5}\selectfont\color{popink},
  fontlower=\popsemi\fontsize{10.6}{14.5}\selectfont\color{popink}}}
% Coche dessinée (le glyphe ✓ n'existe pas dans les polices du document)
\newcommand{\popcheck}[1]{\tikz[baseline=-0.5ex]\draw[#1,line width=1.6pt,line cap=round,line join=round](-0.13,0.0)--(-0.04,-0.09)--(0.13,0.1);}
\providecommand{\coche}{\popcheck{popgreen}}
\providecommand{\resultat}[1]{\par\smallskip\noindent\fcolorbox{popgreen}{popgreenL}{\parbox{\dimexpr\linewidth-2\fboxsep-2\fboxrule\relax}{\textbf{Result.} #1}}\par\smallskip}
\newcommand{\popboxhead}[3]{\poppill{#1}{#2}{white}\ifx\relax#3\relax\else\hspace{6pt}{\popxbold\fontsize{10.6}{12}\selectfont\color{popink}#3}\fi\par\vspace{7pt}}

\newtcolorbox{aretenir}[1][]{popbase,colframe=popgreen,before upper={\popboxhead{Key points}{popgreen}{#1}}}
\newtcolorbox{exemplebox}[1][]{popbase,colframe=poporange,before upper={\popboxhead{Exemple}{poporange}{#1}}}
\newtcolorbox{definition}[1][]{popbase,colframe=popblue,before upper={\popboxhead{Definition}{popblue}{#1}}}
\newtcolorbox{propriete}[1][]{popbase,colframe=poppurple,before upper={\popboxhead{Property}{poppurple}{#1}}}
\newtcolorbox{methode}[1][]{popbase,colframe=popgold,before upper={\popboxhead{Method}{popgold}{#1}}}
\newtcolorbox{attention}[1][]{popbase,colframe=poppink,before upper={\popboxhead{Warning}{poppink}{#1}}}
\newtcolorbox{experience}[1][]{popbase,colframe=popblue,colback=popblueL,before upper={\popboxhead{Experiment}{popblue}{#1}}}
% « Le savais-tu ? » : carte violette pleine (contexte, culture, Cameroun)
\newtcolorbox{savaistu}[1][]{popbase,colback=poppurple,colframe=popink,
  fontupper=\popsemi\fontsize{10.4}{14.2}\selectfont\color{white},
  before upper={\poppill{Did you know?}{white}{poppurple}\ifx\relax#1\relax\else\hspace{6pt}{\popxbold\color{white}#1}\fi\par\vspace{7pt}}}
% Exercice résolu : énoncé en haut, \tcblower puis la solution sur fond crème
\newcounter{popexo}\newcounter{popexr}
\newtcolorbox{exoresolu}[1][]{popbase,colframe=poporange,colbacklower=popcream,
  segmentation style={popink,line width=1.2pt,dashed},
  before upper={\stepcounter{popexr}\popboxhead{Worked example \thepopexr}{poporange}{#1}},
  before lower={\poppill{Solution}{popink}{popcream}\par\vspace{6pt}}}
% Exercice (non résolu en place) — la correction est dans la partie Corrigés
\newtcolorbox{exercice}[1][]{popbase,colframe=popink,
  before upper={\stepcounter{popexo}\popboxhead{Exercise \thepopexo}{popink}{#1}}}
% Corrigé
\newtcolorbox{corrige}[1][]{popbase,colframe=popgreen,colback=popgreenL,
  before upper={\popboxhead{Solution}{popgreen}{#1}}}
% Objectifs / prérequis / bilan
\newtcolorbox{objectifs}{popbase,colframe=popink,colback=poporangeL,
  before upper={\popboxhead{Chapter objectives}{popink}{}}}
\newtcolorbox{prerequis}{popbase,colframe=popink,colback=popblueL,
  before upper={\popboxhead{Prerequisites}{popblue}{}}}
\newtcolorbox{bilan}{popbase,colframe=popink,colback=white,boxrule=2.8pt,
  before upper={\popboxhead{Fiche bilan}{poporange}{L'essentiel à connaître pour l'examen}}}
% Figure encadrée avec légende
\newtcolorbox{popfigure}[1][]{enhanced,breakable=false,colback=white,colframe=popline,boxrule=1.4pt,arc=8pt,
  left=10pt,right=10pt,top=10pt,bottom=8pt,before skip=10pt,after skip=12pt,halign=center,
  lower separated=false,halign lower=center,
  fontlower=\popsemi\fontsize{9}{11.5}\selectfont\color{popmuted},
  #1}
\newcommand{\legende}[1]{\par\vspace{6pt}{\popsemi\fontsize{9}{11.5}\selectfont\color{popmuted}#1}}

% ============================================================
% Signature OnBuch+
% ============================================================
\newcommand{\poplogoinline}[1]{{\poplogo\fontsize{#1}{#1}\selectfont\color{popink}OnBuch\color{poporange}+}}
\newcommand{\signatureonbuch}{%
  \begin{tcolorbox}[enhanced,colback=popink,colframe=popink,boxrule=2.2pt,arc=10pt,
      drop shadow=poporange,left=14pt,right=14pt,top=10pt,bottom=10pt,before skip=0pt,after skip=0pt]
    \tikz[baseline=-0.7ex]\node[circle,fill=popgreen,draw=popcream,line width=1.3pt,minimum size=22pt,inner sep=0]{\popcheck{white}};%
    \hspace{9pt}\begin{minipage}[c]{0.62\linewidth}
      {\popxbold\fontsize{12}{13}\selectfont\color{popcream}Signed\ \textbullet\ {\poplogo OnBuch\color{poporange}+}}\par\vspace{2pt}
      {\raggedright\popsemi\fontsize{8}{10.5}\selectfont\color{popline}\MakeUppercase{OnBuch+ teaching team}\\\MakeUppercase{Follows the official MINESEC syllabus}\par}
    \end{minipage}\hfill
    {\poplogo\fontsize{22}{22}\selectfont\color{popcream}OnBuch\color{poporange}+}
  \end{tcolorbox}%
}

% ============================================================
% Page de garde
% #1 titre · #2 matière · #3 niveau · #4 module · #5 n° leçon · #6 durée ·
% #7 description / objectifs (texte ou itemize)
% ============================================================
\newcommand{\pagegarde}[7]{%
  \thispagestyle{empty}
  \newgeometry{a4paper,margin=1.7cm,top=1.6cm,bottom=1.4cm}%
  \begin{tikzpicture}[remember picture,overlay]
    \fill[poporange!18] ([xshift=-3.2cm,yshift=-3.6cm]current page.north east) circle (4.6cm);
    \fill[poppurple!14] ([xshift=2.4cm,yshift=7.6cm]current page.south west) circle (3.8cm);
  \end{tikzpicture}%
  {\poplogo\fontsize{30}{30}\selectfont\color{popink}OnBuch\color{poporange}+}\hfill
  \raisebox{0.6ex}{\poppill{Signed course \textbullet\ Premium}{popink}{popcream}}\par
  \vspace{1pt}\popsquiggle{poppurple}\par
  \vspace{8pt}
  \begin{tcolorbox}[enhanced,colback=poporange,colframe=popink,boxrule=2.8pt,arc=13pt,
      drop shadow=popink,left=18pt,right=18pt,top=12pt,bottom=13pt,after skip=10pt]
    \poppill{#2 \textbullet\ #3}{popink}{popcream}\hspace{5pt}\poppill{#4}{white}{popink}\par\vspace{8pt}
    {\popdisplay\fontsize{14}{14}\selectfont\color{popink}CHAPTER #5}\par\vspace{4pt}
    {\raggedright\hyphenpenalty=10000\exhyphenpenalty=10000\popdisplay\fontsize{25}{27}\selectfont\color{popcream}\MakeUppercase{#1}\par}
    \vspace{8pt}
    {\popxbold\fontsize{9.5}{9.5}\selectfont\color{popink}\MakeUppercase{Suggested time: #6}}
  \end{tcolorbox}
  \begin{tcolorbox}[enhanced,colback=white,colframe=popink,boxrule=2.2pt,arc=10pt,
      drop shadow=popink,left=16pt,right=16pt,top=10pt,bottom=10pt,after skip=8pt]
    \poppill{In this chapter}{poppurple}{white}\par\vspace{5pt}
    \popsemi\fontsize{9.3}{12.6}\selectfont\color{popink}#7
  \end{tcolorbox}
  \noindent\begin{minipage}[t]{0.487\linewidth}
    \begin{tcolorbox}[enhanced,colback=popgreen,colframe=popink,boxrule=2.2pt,arc=10pt,
        drop shadow=popink,left=11pt,right=11pt,top=10pt,bottom=10pt,height=3.1cm,valign=center]
      \tcbox[colback=white,colframe=popink,boxrule=1.3pt,arc=5pt,boxsep=0pt,nobeforeafter,
        left=3pt,right=3pt,top=3pt,bottom=3pt,valign=center]{\qrcode[height=1.9cm]{https://play.google.com/store/apps/details?id=com.luvvix.onbuch&hl=en}}%
      \hspace{8pt}\begin{minipage}[c]{0.5\linewidth}
        {\popdisplay\fontsize{11}{12.5}\selectfont\color{white}\MakeUppercase{Download OnBuch}}\par\vspace{4pt}
        {\raggedright\popsemi\fontsize{8.4}{11}\selectfont\color{white}Free \textbullet\ Google Play\\AI tutor, past papers, results\par}
      \end{minipage}
    \end{tcolorbox}
  \end{minipage}\hfill
  \begin{minipage}[t]{0.487\linewidth}
    \begin{tcolorbox}[enhanced,colback=poppurple,colframe=popink,boxrule=2.2pt,arc=10pt,
        drop shadow=popink,left=11pt,right=11pt,top=10pt,bottom=10pt,height=3.1cm,valign=center]
      \tikz[baseline=-0.7ex]\node[circle,fill=white,draw=popink,line width=1.3pt,minimum size=1.6cm,inner sep=0]{\popdisplay\fontsize{24}{24}\selectfont\color{poppurple}+};%
      \hspace{8pt}\begin{minipage}[c]{0.6\linewidth}
        {\popdisplay\fontsize{11}{12.5}\selectfont\color{white}\MakeUppercase{Go OnBuch+}}\par\vspace{4pt}
        {\raggedright\popsemi\fontsize{8.4}{11}\selectfont\color{white}All signed courses, unlimited AI tutor, PDF sheets — OnBuch+ tab in the app\par}
      \end{minipage}
    \end{tcolorbox}
  \end{minipage}
  \vspace{8pt}
  \popcontenu
  \vfill
  \signatureonbuch
  \restoregeometry
  \newpage
}

% Rangée « Ce cours contient » (page de garde)
\newcommand{\poptile}[3]{% #1 couleur #2 gros texte #3 légende
  \begin{tcolorbox}[enhanced,colback=#1,colframe=popink,boxrule=2pt,arc=9pt,shadow={2.5pt}{-2.5pt}{0pt}{popink},
      left=6pt,right=6pt,top=9pt,bottom=9pt,halign=center,valign=center,height=1.75cm,width=\linewidth,nobeforeafter]
    {\popdisplay\fontsize{12.5}{13}\selectfont\color{white}\MakeUppercase{#2}}\par\vspace{3pt}
    {\centering\popsemi\fontsize{7.8}{9.5}\selectfont\color{white}#3\par}
  \end{tcolorbox}}
\newcommand{\popcontenu}{%
  \par\noindent{\popxboldls\fontsize{7.5}{7.5}\selectfont\color{popmuted}THIS SIGNED COURSE CONTAINS}\par\vspace{7pt}
  \noindent\begin{minipage}[t]{0.235\linewidth}\poptile{popblue}{Course}{complete, illustrated, step by step}\end{minipage}\hfill
  \begin{minipage}[t]{0.235\linewidth}\poptile{poporange}{Methods}{and worked examples}\end{minipage}\hfill
  \begin{minipage}[t]{0.235\linewidth}\poptile{poppink}{Exercises}{exam style + solutions}\end{minipage}\hfill
  \begin{minipage}[t]{0.235\linewidth}\poptile{popgreen}{Summary sheet}{the essentials to remember}\end{minipage}\par}

% Sommaire
\newtcolorbox{sommaire}{enhanced,colback=white,colframe=popink,boxrule=2.2pt,arc=10pt,
  drop shadow=popink,left=18pt,right=18pt,top=13pt,bottom=13pt,before skip=6pt,after skip=20pt,
  before upper={\poppill{Contents}{poppurple}{white}\par\vspace{9pt}\popsemi\fontsize{10.6}{14.5}\selectfont\color{popink}}}

% Signature de fin de cours — poussée tout en bas de la dernière page
\newcommand{\findecours}{%
  \par\vfill\nopagebreak
  \begin{center}\popsquiggle{poporange}\end{center}\vspace{4pt}
  \signatureonbuch
}
\AtBeginDocument{\def\llbracket{\mathopen{[\mkern-3.5mu[}}\def\rrbracket{\mathclose{]\mkern-3.5mu]}}} % intervalles d'entiers (glyphe ⟦ absent de la police)
% Glyphes absents de Fira Math : redéfinis à partir de glyphes disponibles
\AtBeginDocument{%
  \def\setminus{\mathbin{\backslash}}\let\smallsetminus\setminus
  \def\vdots{\mathord{\vbox{\baselineskip3.2pt\lineskiplimit0pt\kern4pt\hbox{.}\hbox{.}\hbox{.}}}}%
  \def\ddots{\mathinner{\mkern1mu\raise6.5pt\hbox{.}\mkern2mu\raise3.6pt\hbox{.}\mkern2mu\raise0.7pt\hbox{.}\mkern1mu}}%
  \def\triangle{\mathord{\scalebox{0.9}{$\symup{\Delta}$}}}%
  \def\nabla{\mathord{\raisebox{\depth}{\scalebox{1}[-1]{$\symup{\Delta}$}}}}%
  \def\checkmark{\popcheck{popdark}}%
  \def\star{\mathbin{\ast}}\def\bigstar{\mathord{\ast}}%
  \def\diamond{\mathbin{\scalebox{0.6}{$\Diamond$}}}%
  \def\bigcup{\mathop{\mathchoice{\raisebox{-0.3em}{\scalebox{1.7}{$\cup$}}}{\scalebox{1.25}{$\cup$}}{\cup}{\cup}}\displaylimits}%
  \def\bigcap{\mathop{\mathchoice{\raisebox{-0.3em}{\scalebox{1.7}{$\cap$}}}{\scalebox{1.25}{$\cap$}}{\cap}{\cap}}\displaylimits}%
  \def\lesseqqgtr{\mathrel{\lessgtr}}\def\bigoplus{\mathop{\oplus}\displaylimits}%
}
\providecommand{\ssi}{if and only if} % abréviation fréquente
\AtBeginDocument{\providecommand{\wideparen}{\overparen}} % arc de cercle

% Fonctions usuelles que les modèles utilisent sans que amsmath les fournisse
\providecommand{\arsinh}{\operatorname{arsinh}}\providecommand{\arcosh}{\operatorname{arcosh}}
\providecommand{\artanh}{\operatorname{artanh}}\providecommand{\arsech}{\operatorname{arsech}}
\providecommand{\arcsch}{\operatorname{arcsch}}\providecommand{\arcoth}{\operatorname{arcoth}}
\providecommand{\arcsinh}{\operatorname{arsinh}}\providecommand{\arccosh}{\operatorname{arcosh}}
\providecommand{\arctanh}{\operatorname{artanh}}\providecommand{\sech}{\operatorname{sech}}
\providecommand{\csch}{\operatorname{csch}}\providecommand{\arcsec}{\operatorname{arcsec}}
\providecommand{\arccsc}{\operatorname{arccsc}}\providecommand{\arccot}{\operatorname{arccot}}
\providecommand{\sgn}{\operatorname{sgn}}\providecommand{\lcm}{\operatorname{lcm}}
\providecommand{\Var}{\operatorname{Var}}\providecommand{\Cov}{\operatorname{Cov}}

%% FIGFIX-BEGIN v5 — correctifs globaux des figures (docs/onbuch-generation/tools/figfix)
\usepackage{adjustbox}\usepackage{etoolbox}\tcbuselibrary{raster}
% toute figure TikZ/pgfplots plus large que la ligne est réduite à la largeur disponible
\BeforeBeginEnvironment{tikzpicture}{\begin{adjustbox}{max width=\linewidth}}
\AfterEndEnvironment{tikzpicture}{\end{adjustbox}}
% légendes pgfplots lisibles par-dessus les courbes
\pgfplotsset{every axis/.append style={legend style={fill=white,fill opacity=0.92,text opacity=1,draw=black!15,inner sep=3pt}}}
% étiquettes d'arêtes (syntaxe edge["..."]) avec fond blanc
\tikzset{every edge quotes/.append style={fill=white,inner sep=1.2pt}}
%% FIGFIX-END
\begin{document}

\pagegarde{Vector Product and its Applications}{Further Mathematics}{Upper Sixth}{Upper Sixth — Further mathematics}{14}{3 periods}{This chapter introduces the vector (cross) product of two vectors in space, its algebraic and geometric properties, and its spectacular applications: areas, volumes via the scalar triple product, and the vector and Cartesian equations of a plane. It is a cornerstone chapter of the GCE Advanced Level Further Mathematics paper, where spatial geometry questions are guaranteed.
\vspace{4pt}
\begin{multicols}{2}\small
\begin{itemize}
\item By the end of this chapter I can define the vector product of two vectors and compute it both from the determinant formula and from the geometric definition.
\item By the end of this chapter I can state and use the algebraic properties of the vector product (anticommutativity, distributivity, behaviour with scalar multiples).
\item By the end of this chapter I can interpret the magnitude of a vector product as the area of a parallelogram or triangle.
\item By the end of this chapter I can compute the scalar triple product and use it to find volumes of parallelepipeds and tetrahedra.
\item By the end of this chapter I can use the scalar triple product to test whether three vectors are coplanar or four points are coplanar.
\item By the end of this chapter I can determine a normal vector to a plane defined by three points or by two directions.
\end{itemize}\end{multicols}}

\begin{sommaire}
\begin{multicols}{2}
\begin{enumerate}
\item Section 1: The Vector Product — Definition and Computation
\item Section 2: Algebraic Properties and Geometric Applications (Areas)
\item Section 3: The Scalar Triple Product and Volumes
\item Section 4: Equations of a Plane — Vector and Cartesian Forms
\item Integration activity
\item Methods and worked examples
\item Exercises
\item Solutions to the exercises
\item Summary sheet
\end{enumerate}
\end{multicols}
\end{sommaire}

\begin{objectifs}
\begin{itemize}
\item By the end of this chapter I can define the vector product of two vectors and compute it both from the determinant formula and from the geometric definition.
\item By the end of this chapter I can state and use the algebraic properties of the vector product (anticommutativity, distributivity, behaviour with scalar multiples).
\item By the end of this chapter I can interpret the magnitude of a vector product as the area of a parallelogram or triangle.
\item By the end of this chapter I can compute the scalar triple product and use it to find volumes of parallelepipeds and tetrahedra.
\item By the end of this chapter I can use the scalar triple product to test whether three vectors are coplanar or four points are coplanar.
\item By the end of this chapter I can determine a normal vector to a plane defined by three points or by two directions.
\item By the end of this chapter I can write the vector equation and the Cartesian equation of a plane and pass from one form to the other.
\item By the end of this chapter I can find the perpendicular distance from a point to a plane and the angle between two planes (GCE extension).
\item By the end of this chapter I can solve GCE-style problems combining vector products, triple products and plane equations.
\end{itemize}
\end{objectifs}

\begin{prerequis}
\begin{itemize}
\item Vectors in space: components, magnitude, addition and scalar multiplication (Lower Sixth).
\item The scalar (dot) product and its use for angles and perpendicularity (Lower Sixth).
\item Position vectors, vector equation of a line in space (Lower Sixth).
\item Determinants of 2x2 and 3x3 matrices (Further Pure, Lower Sixth).
\item Basic trigonometry: sine of an angle, area of a triangle (1/2)ab sin C.
\end{itemize}
\end{prerequis}

\newpage

\coursec{Section 1: The Vector Product — Definition and Computation}

At a construction site in Bonabéri, a mason must check that a metal gate is perfectly rectangular: he measures the diagonals, but an engineer instead computes a cross product of two edge vectors to test perpendicularity and find the exact area in one step. Meanwhile, a carpenter in Bamenda wants the volume of a wedge‑shaped roof beam, and a surveyor near Mount Cameroon needs the equation of an inclined plane to position a water tank. All three problems, seemingly different, are solved by one powerful tool: the \cle{vector product} (also called the \cle{cross product}). How can a single operation between two vectors give directions, areas, volumes and even the equations of planes in space? This section introduces the vector product, its geometric meaning, its algebraic properties and the practical determinant method for computation.

\courssub{Motivation and Geometric Definition}

In three‑dimensional space, given two non‑zero vectors $\vec{a}$ and $\vec{b}$, there are infinitely many vectors perpendicular to both. To single out one, we need both a magnitude and an orientation. The vector product provides exactly that: a vector perpendicular to the plane containing $\vec{a}$ and $\vec{b}$, whose length equals the area of the parallelogram spanned by $\vec{a}$ and $\vec{b}$, and whose direction follows the \cle{right‑hand rule}.

\begin{definition}[Vector Product — Geometric Definition]
Let $\vec{a}$ and $\vec{b}$ be two vectors in $\mathbb{R}^3$ and let $\theta\in[0,\pi]$ be the angle between them. The \textbf{vector product} (or \textbf{cross product}) $\vec{a}\times\vec{b}$ is the vector defined by
\[
\vec{a}\times\vec{b}=|\vec{a}|\,|\vec{b}|\,\sin\theta\;\vec{n},
\]
where $\vec{n}$ is the unit vector perpendicular to both $\vec{a}$ and $\vec{b}$ such that $(\vec{a},\vec{b},\vec{n})$ forms a right‑handed system (right‑hand rule: curl the fingers of the right hand from $\vec{a}$ towards $\vec{b}$; the extended thumb points in the direction of $\vec{n}$).
\end{definition}

\begin{popfigure}
\begin{tikzpicture}[scale=1.2,>=stealth]
% coordinate axes
\draw[->,popmuted] (0,0,0) -- (2.5,0,0) node[right] {$x$};
\draw[->,popmuted] (0,0,0) -- (0,2.5,0) node[above] {$y$};
\draw[->,popmuted] (0,0,0) -- (0,0,2.5) node[above left] {$z$};
% vectors a and b
\draw[->,popblue,very thick] (0,0,0) -- (2,0.5,0) node[midway,below right] {$\vec{a}$};
\draw[->,poporange,very thick] (0,0,0) -- (0.5,2,0) node[midway,above left] {$\vec{b}$};
% plane shading
\fill[popblueL,opacity=0.2] (0,0,0) -- (2,0.5,0) -- (2.5,2.5,0) -- (0.5,2,0) -- cycle;
% cross product vector
\draw[->,popgreen,very thick] (0,0,0) -- (0,0,2) node[midway,right] {$\vec{a}\times\vec{b}$};
% angle arc
\draw[popdark,thick] (0.5,0,0) arc (0:30:0.5);
\node[popdark] at (0.7,0.2,0) {$\theta$};
% right-hand rule arrow
\draw[poppurple,thick,->] (1.5,1.2,0.3) arc (0:90:0.5 and 0.3);
\node[poppurple,align=center] at (1.8,1.5,0.5) {Right-hand\\rule};
\end{tikzpicture}
\legende{Geometric definition of $\vec{a}\times\vec{b}$: magnitude $|\vec{a}||\vec{b}|\sin\theta$, direction by right‑hand rule.}
\end{popfigure}

Immediate consequences of the definition:
\begin{itemize}
\item If $\vec{a}=\vec{0}$ or $\vec{b}=\vec{0}$, then $\vec{a}\times\vec{b}=\vec{0}$.
\item $\vec{a}\times\vec{a}=\vec{0}$ because $\theta=0$ and $\sin0=0$.
\item For non‑zero vectors, $\vec{a}\times\vec{b}=\vec{0}$ \textbf{if and only if} $\vec{a}$ and $\vec{b}$ are parallel (collinear), i.e. $\theta=0$ or $\pi$.
\end{itemize}

\begin{propriete}[Zero Cross Product]
For non‑zero vectors $\vec{a},\vec{b}\in\mathbb{R}^3$,
\[
\vec{a}\times\vec{b}=\vec{0}\quad\Longleftrightarrow\quad \vec{a}\text{ and }\vec{b}\text{ are collinear}.
\]
\end{propriete}

\courssub{Algebraic Properties of the Vector Product}

The vector product is \textbf{not commutative}; reversing the order changes the sign because the right‑hand rule flips the orientation.

\begin{propriete}[Anti‑commutativity]
\[
\vec{a}\times\vec{b}=-(\vec{b}\times\vec{a}).
\]
\end{propriete}

The standard orthonormal basis $\{\vec{i},\vec{j},\vec{k}\}$ obeys a cyclic rule that is easy to remember.

\begin{propriete}[Cross Products of Basis Vectors]
\[
\vec{i}\times\vec{j}=\vec{k},\qquad
\vec{j}\times\vec{k}=\vec{i},\qquad
\vec{k}\times\vec{i}=\vec{j},
\]
and the reverse products give the opposite vectors:
\[
\vec{j}\times\vec{i}=-\vec{k},\qquad
\vec{k}\times\vec{j}=-\vec{i},\qquad
\vec{i}\times\vec{k}=-\vec{j}.
\]

\begin{popfigure}
\begin{tikzpicture}[scale=1.5,>=stealth]
% axes
\draw[->,popmuted] (0,0,0) -- (1.5,0,0) node[below right] {$\vec{i}$};
\draw[->,popmuted] (0,0,0) -- (0,1.5,0) node[above left] {$\vec{j}$};
\draw[->,popmuted] (0,0,0) -- (0,0,1.5) node[above] {$\vec{k}$};
% cyclic arrows
\draw[popblue,thick,->] (0.8,0.2,0) arc (0:90:0.8 and 0.2);
\node[popblue] at (1.0,0.6,0) {$\vec{i}\times\vec{j}=\vec{k}$};
\draw[poporange,thick,->] (0.2,0.8,0) arc (90:180:0.2 and 0.8);
\node[poporange] at (-0.2,1.0,0) {$\vec{j}\times\vec{k}=\vec{i}$};
\draw[popgreen,thick,->] (0,0.2,0.8) arc (90:0:0.2 and 0.8);
\node[popgreen] at (0.2,0,1.0) {$\vec{k}\times\vec{i}=\vec{j}$};
\end{tikzpicture}
\legende{Cyclic rule for the basis vectors: $\vec{i}\to\vec{j}\to\vec{k}\to\vec{i}$.}
\end{popfigure}
\end{propriete}

The vector product is \textbf{distributive} over addition:
\[
\vec{a}\times(\vec{b}+\vec{c})=\vec{a}\times\vec{b}+\vec{a}\times\vec{c},
\qquad
(\vec{a}+\vec{b})\times\vec{c}=\vec{a}\times\vec{c}+\vec{b}\times\vec{c}.
\]
It is \textbf{not associative}: $(\vec{a}\times\vec{b})\times\vec{c}\neq\vec{a}\times(\vec{b}\times\vec{c})$ in general.

\courssub{Component Formula and Determinant Method}

If $\vec{a}=(a_1,a_2,a_3)$ and $\vec{b}=(b_1,b_2,b_3)$ in the orthonormal basis, expanding $(\vec{a}\times\vec{b})$ using the basis rules yields the component formula.

\begin{definition}[Vector Product — Component Form]
\[
\vec{a}\times\vec{b}=
\begin{pmatrix}
a_2b_3-a_3b_2\\
a_3b_1-a_1b_3\\
a_1b_2-a_2b_1
\end{pmatrix}.
\]
This is conveniently remembered as the expansion of the symbolic $3\times3$ determinant
\[
\vec{a}\times\vec{b}=
\begin{vmatrix}
\vec{i} & \vec{j} & \vec{k}\\
a_1   & a_2   & a_3\\
b_1   & b_2   & b_3
\end{vmatrix}.
\]
\end{definition}

\begin{methode}[Computing $\vec{a}\times\vec{b}$ with the Determinant Scheme]
\begin{enumerate}
\item Write the determinant with first row $\vec{i},\vec{j},\vec{k}$, second row components of $\vec{a}$, third row components of $\vec{b}$.
\item Expand along the first row (cofactor expansion):
\begin{align*}
\vec{a}\times\vec{b}&=
\vec{i}\begin{vmatrix}a_2&a_3\\b_2&b_3\end{vmatrix}
-\vec{j}\begin{vmatrix}a_1&a_3\\b_1&b_3\end{vmatrix}
+\vec{k}\begin{vmatrix}a_1&a_2\\b_1&b_2\end{vmatrix}\\
&=(a_2b_3-a_3b_2)\vec{i}+(a_3b_1-a_1b_3)\vec{j}+(a_1b_2-a_2b_1)\vec{k}.
\end{align*}
\item \textbf{Check perpendicularity} (optional but recommended): compute $\vec{a}\cdot(\vec{a}\times\vec{b})$ and $\vec{b}\cdot(\vec{a}\times\vec{b})$; both must be $0$.
\end{enumerate}
\end{methode}

\courssub{Worked Examples}

\begin{exoresolu}[Computing a Cross Product]
Let $\vec{a}=2\vec{i}-\vec{j}+3\vec{k}$ and $\vec{b}=\vec{i}+4\vec{j}-2\vec{k}$. Compute $\vec{a}\times\vec{b}$ and verify that it is perpendicular to both $\vec{a}$ and $\vec{b}$.
\tcblower
\textbf{Solution.}
\[
\vec{a}\times\vec{b}=
\begin{vmatrix}
\vec{i} & \vec{j} & \vec{k}\\
2 & -1 & 3\\
1 & 4 & -2
\end{vmatrix}
=\vec{i}\begin{vmatrix}-1&3\\4&-2\end{vmatrix}
-\vec{j}\begin{vmatrix}2&3\\1&-2\end{vmatrix}
+\vec{k}\begin{vmatrix}2&-1\\1&4\end{vmatrix}.
\]
Calculate the $2\times2$ determinants:
\begin{align*}
\begin{vmatrix}-1&3\\4&-2\end{vmatrix}&=(-1)(-2)-3\cdot4=2-12=-10,\\
\begin{vmatrix}2&3\\1&-2\end{vmatrix}&=2(-2)-3\cdot1=-4-3=-7,\\
\begin{vmatrix}2&-1\\1&4\end{vmatrix}&=2\cdot4-(-1)\cdot1=8+1=9.
\end{align*}
Hence
\[
\vec{a}\times\vec{b}=-10\vec{i}-(-7)\vec{j}+9\vec{k}=-10\vec{i}+7\vec{j}+9\vec{k}.
\]
\textbf{Verification:}
\[
\vec{a}\cdot(\vec{a}\times\vec{b})=(2)(-10)+(-1)(7)+(3)(9)=-20-7+27=0,
\]
\[
\vec{b}\cdot(\vec{a}\times\vec{b})=(1)(-10)+(4)(7)+(-2)(9)=-10+28-18=0.
\]
Both dot products are zero, confirming perpendicularity.
\end{exoresolu}

\begin{exemplebox}[Cross Product of Parallel Vectors]
Let $\vec{u}=3\vec{i}-6\vec{j}+9\vec{k}$ and $\vec{v}=-\vec{i}+2\vec{j}-3\vec{k}$. Notice $\vec{u}=-3\vec{v}$, so they are collinear. Compute $\vec{u}\times\vec{v}$.
\tcblower
Because $\vec{u}$ and $\vec{v}$ are parallel, the angle $\theta=0$ (or $\pi$), $\sin\theta=0$, and the cross product must be the zero vector. Using the determinant:
\[
\vec{u}\times\vec{v}=
\begin{vmatrix}
\vec{i} & \vec{j} & \vec{k}\\
3 & -6 & 9\\
-1 & 2 & -3
\end{vmatrix}
=\vec{i}(18-18)-\vec{j}(-9+9)+\vec{k}(6-6)=\vec{0}.
\]
This illustrates the property $\vec{a}\times\vec{b}=\vec{0}\iff\vec{a}\parallel\vec{b}$ (for non‑zero vectors).
\end{exemplebox}

\courssub{Geometric Interpretation: Magnitude and Angle}

The magnitude of the cross product gives the area of the parallelogram spanned by $\vec{a}$ and $\vec{b}$ and allows us to find $\sin\theta$.

\begin{propriete}[Magnitude of the Cross Product]
\[
|\vec{a}\times\vec{b}|=|\vec{a}|\,|\vec{b}|\,\sin\theta.
\]
Consequently, for non‑zero vectors,
\[
\sin\theta=\frac{|\vec{a}\times\vec{b}|}{|\vec{a}|\,|\vec{b}|}.
\]
\end{propriete}

\begin{exoresolu}[Finding the Sine of the Angle]
Given $\vec{p}=2\vec{i}+\vec{j}-\vec{k}$ and $\vec{q}=\vec{i}-2\vec{j}+2\vec{k}$, find $\sin\theta$ where $\theta$ is the angle between $\vec{p}$ and $\vec{q}$.
\tcblower
First compute $\vec{p}\times\vec{q}$:
\[
\vec{p}\times\vec{q}=
\begin{vmatrix}
\vec{i} & \vec{j} & \vec{k}\\
2 & 1 & -1\\
1 & -2 & 2
\end{vmatrix}
=\vec{i}(2-2)-\vec{j}(4+1)+\vec{k}(-4-1)=0\vec{i}-5\vec{j}-5\vec{k}.
\]
Magnitudes:
\[
|\vec{p}\times\vec{q}|=\sqrt{0^2+(-5)^2+(-5)^2}=\sqrt{50}=5\sqrt{2},
\]
\[
|\vec{p}|=\sqrt{2^2+1^2+(-1)^2}=\sqrt{6},\qquad
|\vec{q}|=\sqrt{1^2+(-2)^2+2^2}=\sqrt{9}=3.
\]
Thus
\[
\sin\theta=\frac{5\sqrt{2}}{\sqrt{6}\cdot3}=\frac{5\sqrt{2}}{3\sqrt{6}}=\frac{5}{3\sqrt{3}}=\frac{5\sqrt{3}}{9}.
\]
\end{exoresolu}

\courssub{Checking Perpendicularity — A Quick Test}

The cross product is the standard tool to obtain a vector perpendicular to two given vectors. After computing $\vec{a}\times\vec{b}$, a quick dot‑product check guarantees no algebraic slip.

\begin{methode}[Verifying Perpendicularity]
After obtaining $\vec{c}=\vec{a}\times\vec{b}$, compute
\[
\vec{a}\cdot\vec{c}\quad\text{and}\quad\vec{b}\cdot\vec{c}.
\]
If both are zero, $\vec{c}$ is perpendicular to both $\vec{a}$ and $\vec{b}$. If not, re‑check the determinant expansion.
\end{methode}

\begin{exemplebox}[Perpendicular Vector to a Plane]
Find a vector perpendicular to the plane containing points $A(1,2,-1)$, $B(3,0,2)$ and $C(-1,1,3)$.
\tcblower
Two vectors in the plane are $\vec{AB}=B-A=(2,-2,3)$ and $\vec{AC}=C-A=(-2,-1,4)$. Their cross product is perpendicular to the plane:
\[
\vec{AB}\times\vec{AC}=
\begin{vmatrix}
\vec{i} & \vec{j} & \vec{k}\\
2 & -2 & 3\\
-2 & -1 & 4
\end{vmatrix}
=\vec{i}(-8+3)-\vec{j}(8+6)+\vec{k}(-2-4)=-5\vec{i}-14\vec{j}-6\vec{k}.
\]
Check: $(2,-2,3)\cdot(-5,-14,-6)=-10+28-18=0$; $(-2,-1,4)\cdot(-5,-14,-6)=10+14-24=0$. Hence $(-5,-14,-6)$ is a normal vector to the plane.
\end{exemplebox}

\courssub{Common Mistakes and Warnings}

\begin{attention}[Cross Product is Not Commutative]
$\vec{a}\times\vec{b}=-(\vec{b}\times\vec{a})$. Swapping the two rows of the determinant changes the sign of every component. Always keep the order as given in the problem.
\end{attention}

\begin{attention}[Do Not Confuse Dot and Cross Products]
\begin{itemize}
\item Dot product $\vec{a}\cdot\vec{b}$ is a \textbf{scalar} (number): $a_1b_1+a_2b_2+a_3b_3$.
\item Cross product $\vec{a}\times\vec{b}$ is a \textbf{vector}: $(a_2b_3-a_3b_2,\;a_3b_1-a_1b_3,\;a_1b_2-a_2b_1)$.
\end{itemize}
Writing $\vec{a}\cdot\vec{b}=\vec{c}$ or $\vec{a}\times\vec{b}=k$ (scalar) is a serious error.
\end{attention}

\begin{attention}[Determinant Sign Errors]
When expanding the $3\times3$ determinant, the $\vec{j}$ component comes with a \textbf{minus} sign:
\[
-\vec{j}\begin{vmatrix}a_1&a_3\\b_1&b_3\end{vmatrix}.
\]
Forgetting this minus sign is the most frequent computational mistake.
\end{attention}

\courssub{Summary of Key Points}

\begin{aretenir}[Vector Product — Essentials]
\begin{itemize}
\item \textbf{Geometric definition:} $\vec{a}\times\vec{b}=|\vec{a}||\vec{b}|\sin\theta\,\vec{n}$, $\vec{n}$ by right‑hand rule.
\item \textbf{Component formula:} $\vec{a}\times\vec{b}=(a_2b_3-a_3b_2,\;a_3b_1-a_1b_3,\;a_1b_2-a_2b_1)$.
\item \textbf{Determinant mnemonic:} $\begin{vmatrix}\vec{i}&\vec{j}&\vec{k}\\a_1&a_2&a_3\\b_1&b_2&b_3\end{vmatrix}$.
\item \textbf{Anti‑commutative:} $\vec{a}\times\vec{b}=-(\vec{b}\times\vec{a})$.
\item \textbf{Zero cross product $\iff$ collinear} (for non‑zero vectors).
\item \textbf{Magnitude:} $|\vec{a}\times\vec{b}|=|\vec{a}||\vec{b}|\sin\theta$ = area of parallelogram.
\item \textbf{Perpendicularity check:} $\vec{a}\cdot(\vec{a}\times\vec{b})=\vec{b}\cdot(\vec{a}\times\vec{b})=0$.
\end{itemize}
\end{aretenir}

The vector product is now at our disposal. In the next sections we will see how it yields areas of triangles and parallelograms, volumes of parallelepipeds and tetrahedra via the scalar triple product, and how it provides the normal vector needed to write the vector and Cartesian equations of a plane.

\coursec{Section 2: Algebraic Properties and Geometric Applications (Areas)}

In Section 1 we learned how to compute the vector product $\vec{a} \times \vec{b}$ using a determinant and we saw its geometric definition: a vector perpendicular to both $\vec{a}$ and $\vec{b}$ with magnitude $|\vec{a}||\vec{b}|\sin\theta$. Now we turn to the algebraic rules that govern this operation. Understanding these properties is essential for manipulating vector expressions in mechanics, geometry, and further mathematics. We will also discover the most immediate geometric application: calculating areas of parallelograms and triangles in three-dimensional space.

\courssub{Algebraic Properties of the Vector Product}

The vector product behaves differently from the scalar product and from ordinary multiplication of numbers. The three fundamental algebraic properties are anticommutativity, distributivity over addition, and compatibility with scalar multiplication.

\begin{propriete}[Algebraic Properties of the Cross Product]
Let $\vec{a}, \vec{b}, \vec{c} \in \mathbb{R}^3$ and $k \in \mathbb{R}$.
\begin{enumerate}
    \item \textbf{Anticommutativity:} $\vec{b} \times \vec{a} = -(\vec{a} \times \vec{b})$.
    \item \textbf{Distributivity over addition:} $\vec{a} \times (\vec{b} + \vec{c}) = \vec{a} \times \vec{b} + \vec{a} \times \vec{c}$.
    \item \textbf{Compatibility with scalar multiplication:} $(k\vec{a}) \times \vec{b} = k(\vec{a} \times \vec{b}) = \vec{a} \times (k\vec{b})$.
\end{enumerate}
\end{propriete}

\begin{exemplebox}[Proof of Anticommutativity from the Determinant]
Recall the determinant definition:
\[
\vec{a} \times \vec{b} = \begin{vmatrix}
\vec{i} & \vec{j} & \vec{k} \\
a_1 & a_2 & a_3 \\
b_1 & b_2 & b_3
\end{vmatrix}.
\]
Swapping the second and third rows (which correspond to the components of $\vec{a}$ and $\vec{b}$) changes the sign of the determinant. Hence
\[
\vec{b} \times \vec{a} = \begin{vmatrix}
\vec{i} & \vec{j} & \vec{k} \\
b_1 & b_2 & b_3 \\
a_1 & a_2 & a_3
\end{vmatrix} = -\begin{vmatrix}
\vec{i} & \vec{j} & \vec{k} \\
a_1 & a_2 & a_3 \\
b_1 & b_2 & b_3
\end{vmatrix} = -(\vec{a} \times \vec{b}).
\]
This also follows immediately from the geometric definition: $\vec{a} \times \vec{b}$ and $\vec{b} \times \vec{a}$ have the same magnitude $|\vec{a}||\vec{b}|\sin\theta$ but opposite directions (right-hand rule vs left-hand rule).
\end{exemplebox}

\begin{attention}[Non-Associativity -- A Critical Trap]
The vector product is \textbf{not associative}. In general,
\[
(\vec{a} \times \vec{b}) \times \vec{c} \neq \vec{a} \times (\vec{b} \times \vec{c}).
\]
\textbf{Counter-example:} Take the standard basis vectors $\vec{i}, \vec{j}, \vec{k}$.
\[
\vec{i} \times (\vec{i} \times \vec{j}) = \vec{i} \times \vec{k} = -\vec{j},
\]
but
\[
(\vec{i} \times \vec{i}) \times \vec{j} = \vec{0} \times \vec{j} = \vec{0}.
\]
Since $-\vec{j} \neq \vec{0}$, associativity fails. \textbf{Never} write $\vec{a} \times \vec{b} \times \vec{c}$ without parentheses; the order of operations matters.
\end{attention}

\begin{exoresolu}[Using Algebraic Properties to Simplify]
Simplify the expression $(\vec{a} + 2\vec{b}) \times (3\vec{a} - \vec{b})$.
\tcblower
\textbf{Solution:} We expand using distributivity and scalar compatibility, treating the cross product like a bilinear operation but respecting the order.
\begin{align*}
(\vec{a} + 2\vec{b}) \times (3\vec{a} - \vec{b})
&= \vec{a} \times (3\vec{a} - \vec{b}) + 2\vec{b} \times (3\vec{a} - \vec{b}) \\
&= 3(\vec{a} \times \vec{a}) - \vec{a} \times \vec{b} + 6(\vec{b} \times \vec{a}) - 2(\vec{b} \times \vec{b}).
\end{align*}
Since $\vec{v} \times \vec{v} = \vec{0}$ for any vector and $\vec{b} \times \vec{a} = -\vec{a} \times \vec{b}$,
\begin{align*}
&= 3\vec{0} - \vec{a} \times \vec{b} + 6(-\vec{a} \times \vec{b}) - 2\vec{0} \\
&= -7(\vec{a} \times \vec{b}).
\end{align*}

\end{exoresolu}

\courssub{Geometric Interpretation: Area of a Parallelogram}

The magnitude of the cross product has a beautiful and extremely useful geometric meaning.

\begin{propriete}[Area of a Parallelogram]
Let $\vec{a}$ and $\vec{b}$ be two non-zero, non-parallel vectors in $\mathbb{R}^3$. The area $A$ of the parallelogram having $\vec{a}$ and $\vec{b}$ as adjacent sides is
\[
A = |\vec{a} \times \vec{b}|.
\]
\end{propriete}

\begin{exemplebox}[Proof]
The area of a parallelogram is base $\times$ height. Take $|\vec{a}|$ as the base. The height is the perpendicular distance from the tip of $\vec{b}$ to the line along $\vec{a}$, which is $|\vec{b}|\sin\theta$, where $\theta$ is the angle between $\vec{a}$ and $\vec{b}$. Hence
\[
\text{Area} = |\vec{a}| \cdot (|\vec{b}|\sin\theta) = |\vec{a}||\vec{b}|\sin\theta = |\vec{a} \times \vec{b}|.
\]
\end{exemplebox}

\begin{popfigure}
\begin{tikzpicture}[scale=1.2, >=Stealth]
% Coordinate axes
\draw[->, thin, gray] (-0.5,0,0) -- (3,0,0) node[anchor=north west] {$x$};
\draw[->, thin, gray] (0,-0.5,0) -- (0,3,0) node[anchor=south east] {$y$};
\draw[->, thin, gray] (0,0,-0.5) -- (0,0,3) node[anchor=south] {$z$};

% Vectors a and b
\draw[popblue, very thick, ->] (0,0,0) -- (2.5,0.5,0) node[midway, below right] {$\vec{a}$};
\draw[poporange, very thick, ->] (0,0,0) -- (0.8,2,0.5) node[midway, left] {$\vec{b}$};

% Parallelogram completion
\draw[popblue, very thick, ->] (0.8,2,0.5) -- (3.3,2.5,0.5);
\draw[poporange, very thick, ->] (2.5,0.5,0) -- (3.3,2.5,0.5);

% Height line (perpendicular from b to a)
\draw[dashed, popdark, thick] (0.8,2,0.5) -- (1.5,0.3,0) node[midway, right, align=center] {$h = |\vec{b}|\sin\theta$};

% Right angle mark
\draw[popdark] (1.4,0.25,0) -- (1.5,0.25,0.1) -- (1.5,0.35,0.1);

% Cross product vector (normal)
\draw[popgreen, very thick, ->] (0,0,0) -- (0,0,2) node[midway, left] {$\vec{a}\times\vec{b}$};

% Angle theta
\draw[popmuted, thick] (0.5,0.1,0) arc (11:60:1.5);
\node[popmuted] at (0.8,0.5,0) {$\theta$};

% Area label
\node[popblue, align=center] at (2,1.5,1) {Area $= |\vec{a}\times\vec{b}|$\\$= |\vec{a}|\,h$};
\end{tikzpicture}
\legende{The parallelogram spanned by $\vec{a}$ and $\vec{b}$. The height is $h = |\vec{b}|\sin\theta$, so area $= |\vec{a}|\,h = |\vec{a}\times\vec{b}|$. The cross product vector is perpendicular to the plane.}
\end{popfigure}

\begin{exoresolu}[Area of a Parallelogram from Vertices]
Find the area of the parallelogram with vertices $A(1,2,-1)$, $B(3,1,0)$, $C(4,3,2)$, $D(2,4,1)$.
\tcblower
\textbf{Solution:} Two adjacent sides are $\vec{AB}$ and $\vec{AD}$.
\[
\vec{AB} = \begin{pmatrix}2\\-1\\1\end{pmatrix},\quad
\vec{AD} = \begin{pmatrix}1\\2\\2\end{pmatrix}.
\]
\[
\vec{AB} \times \vec{AD} = \begin{vmatrix}
\vec{i} & \vec{j} & \vec{k} \\
2 & -1 & 1 \\
1 & 2 & 2
\end{vmatrix}
= \vec{i}(-2-2) - \vec{j}(4-1) + \vec{k}(4+1)
= \begin{pmatrix}-4\\-3\\5\end{pmatrix}.
\]
Area $= |\vec{AB} \times \vec{AD}| = \sqrt{(-4)^2 + (-3)^2 + 5^2} = \sqrt{16+9+25} = \sqrt{50} = 5\sqrt{2}$ square units.

\end{exoresolu}

\courssub{Area of a Triangle in Space}

A triangle is exactly half of a parallelogram. This gives a very efficient formula for the area of a triangle given the coordinates of its vertices in 3D.

\begin{methode}[Area of Triangle $ABC$ in Space]
Given three points $A, B, C$ with position vectors $\vec{a}, \vec{b}, \vec{c}$:
\begin{enumerate}
    \item Compute two side vectors, e.g. $\vec{AB} = \vec{b} - \vec{a}$ and $\vec{AC} = \vec{c} - \vec{a}$.
    \item Compute the cross product $\vec{AB} \times \vec{AC}$.
    \item The area is $\displaystyle \frac{1}{2}|\vec{AB} \times \vec{AC}|$.
\end{enumerate}
\end{methode}

\begin{popfigure}
\begin{tikzpicture}[scale=1.3, >=Stealth]
% Triangle in 3D
\coordinate (A) at (0,0,0);
\coordinate (B) at (3,0.5,0);
\coordinate (C) at (1,2.5,0.5);

\draw[popblue, very thick, ->] (A) -- (B) node[midway, below] {$\vec{AB}$};
\draw[poporange, very thick, ->] (A) -- (C) node[midway, left] {$\vec{AC}$};
\draw[popmuted, thick] (B) -- (C) node[midway, right] {$\vec{BC}$};

% Labels
\node[anchor=north east] at (A) {$A$};
\node[anchor=north west] at (B) {$B$};
\node[anchor=south] at (C) {$C$};

% Cross product as normal vector
\draw[popgreen, very thick, ->] (A) -- (0,0,2) node[midway, left] {$\vec{AB}\times\vec{AC}$};

% Right angle mark at base of normal
\draw[popgreen] (0.1,0,0) -- (0.1,0,0.1) -- (0,0,0.1);

% Area formula label
\node[popdark, align=center, fill=white, fill opacity=0.8, text opacity=1, rounded corners] at (2,1,1.5) {Area $\triangle ABC = \frac{1}{2}|\vec{AB}\times\vec{AC}|$};
\end{tikzpicture}
\legende{Triangle $ABC$ in space. The vector $\vec{AB} \times \vec{AC}$ is perpendicular to the plane of the triangle; its magnitude is twice the area of the triangle.}
\end{popfigure}

\begin{exoresolu}[Area of a Triangle with Given Vertices]
Find the area of triangle $ABC$ where $A(2,-1,3)$, $B(4,0,1)$, $C(1,2,-1)$.
\tcblower
\textbf{Solution:}
\[
\vec{AB} = \begin{pmatrix}2\\1\\-2\end{pmatrix},\quad
\vec{AC} = \begin{pmatrix}-1\\3\\-4\end{pmatrix}.
\]
\[
\vec{AB} \times \vec{AC} = \begin{vmatrix}
\vec{i} & \vec{j} & \vec{k} \\
2 & 1 & -2 \\
-1 & 3 & -4
\end{vmatrix}
= \vec{i}(-4+6) - \vec{j}(-8-2) + \vec{k}(6+1)
= \begin{pmatrix}2\\10\\7\end{pmatrix}.
\]
\[
|\vec{AB} \times \vec{AC}| = \sqrt{2^2 + 10^2 + 7^2} = \sqrt{4+100+49} = \sqrt{153} = 3\sqrt{17}.
\]
Area $= \frac{1}{2} \times 3\sqrt{17} = \frac{3\sqrt{17}}{2}$ square units.

\end{exoresolu}

\courssub{Unit Vector Perpendicular to Two Given Vectors}

Since $\vec{a} \times \vec{b}$ is perpendicular to both $\vec{a}$ and $\vec{b}$, normalising it gives a unit normal vector. This is fundamental for defining plane orientations, flux integrals, and surface normals.

\begin{methode}[Finding a Unit Vector Perpendicular to $\vec{a}$ and $\vec{b}$]
\begin{enumerate}
    \item Compute $\vec{n} = \vec{a} \times \vec{b}$.
    \item Compute its magnitude $|\vec{n}|$.
    \item The required unit vectors are $\pm \dfrac{\vec{n}}{|\vec{n}|}$ (two opposite directions).
\end{enumerate}
If $\vec{a}$ and $\vec{b}$ are parallel, $\vec{a} \times \vec{b} = \vec{0}$ and there is no unique perpendicular direction (infinitely many).
\end{methode}

\begin{exoresolu}[Unit Normal Vector]
Find a unit vector perpendicular to both $\vec{a} = 2\vec{i} - \vec{j} + 3\vec{k}$ and $\vec{b} = \vec{i} + 2\vec{j} - \vec{k}$.
\tcblower
\textbf{Solution:}
\[
\vec{a} \times \vec{b} = \begin{vmatrix}
\vec{i} & \vec{j} & \vec{k} \\
2 & -1 & 3 \\
1 & 2 & -1
\end{vmatrix}
= \vec{i}(1-6) - \vec{j}(-2-3) + \vec{k}(4+1)
= \begin{pmatrix}-5\\5\\5\end{pmatrix}.
\]
Magnitude: $|\vec{a} \times \vec{b}| = \sqrt{25+25+25} = \sqrt{75} = 5\sqrt{3}$.
Unit vector: $\displaystyle \hat{n} = \frac{1}{5\sqrt{3}}\begin{pmatrix}-5\\5\\5\end{pmatrix} = \frac{1}{\sqrt{3}}\begin{pmatrix}-1\\1\\1\end{pmatrix}$.
The opposite vector $-\hat{n} = \frac{1}{\sqrt{3}}\begin{pmatrix}1\\-1\\-1\end{pmatrix}$ is also a valid answer.

\end{exoresolu}

\courssub{Distance from a Point to a Line in Space (GCE Extension)}

Although not explicitly listed in the core syllabus, this application appears frequently in GCE Advanced Level papers as a "useful extension". It combines the cross product magnitude (area) with the parallelogram area formula.

\begin{propriete}[Distance from Point $P$ to Line $\ell$]
Let $\ell$ be a line with direction vector $\vec{u}$ passing through point $A$. The perpendicular distance $d$ from a point $P$ to $\ell$ is
\[
d = \frac{|\vec{AP} \times \vec{u}|}{|\vec{u}|}.
\]
\end{propriete}

\begin{exemplebox}[Derivation]
The area of the parallelogram formed by $\vec{AP}$ and $\vec{u}$ is $|\vec{AP} \times \vec{u}|$. This area also equals base $\times$ height $= |\vec{u}| \cdot d$. Equating the two expressions gives the formula.
\end{exemplebox}

\begin{exoresolu}[Distance from Point to Line]
Find the perpendicular distance from the point $P(1,2,3)$ to the line $\ell: \frac{x-2}{1} = \frac{y+1}{2} = \frac{z}{-1}$.
\tcblower
\textbf{Solution:} The line $\ell$ passes through $A(2,-1,0)$ and has direction vector $\vec{u} = \begin{pmatrix}1\\2\\-1\end{pmatrix}$.
\[
\vec{AP} = \begin{pmatrix}1-2\\2-(-1)\\3-0\end{pmatrix} = \begin{pmatrix}-1\\3\\3\end{pmatrix}.
\]
\[
\vec{AP} \times \vec{u} = \begin{vmatrix}
\vec{i} & \vec{j} & \vec{k} \\
-1 & 3 & 3 \\
1 & 2 & -1
\end{vmatrix}
= \vec{i}(-3-6) - \vec{j}(1-3) + \vec{k}(-2-3)
= \begin{pmatrix}-9\\2\\-5\end{pmatrix}.
\]
\[
|\vec{AP} \times \vec{u}| = \sqrt{81+4+25} = \sqrt{110},\qquad |\vec{u}| = \sqrt{1+4+1} = \sqrt{6}.
\]
\[
d = \frac{\sqrt{110}}{\sqrt{6}} = \sqrt{\frac{110}{6}} = \sqrt{\frac{55}{3}} = \frac{\sqrt{165}}{3}\ \text{units}.
\]

\end{exoresolu}

\begin{aretenir}[Key Points of Section 2]
\begin{itemize}
    \item \textbf{Anticommutativity:} $\vec{b} \times \vec{a} = -(\vec{a} \times \vec{b})$. Order matters!
    \item \textbf{Distributivity:} $\vec{a} \times (\vec{b}+\vec{c}) = \vec{a}\times\vec{b} + \vec{a}\times\vec{c}$.
    \item \textbf{Scalar multiplication:} $(k\vec{a})\times\vec{b} = k(\vec{a}\times\vec{b})$.
    \item \textbf{Non-associativity:} $(\vec{a}\times\vec{b})\times\vec{c} \neq \vec{a}\times(\vec{b}\times\vec{c})$ in general.
    \item \textbf{Area of parallelogram:} $|\vec{a} \times \vec{b}|$.
    \item \textbf{Area of triangle $ABC$:} $\frac{1}{2}|\vec{AB} \times \vec{AC}|$.
    \item \textbf{Unit normal to $\vec{a},\vec{b}$:} $\pm \frac{\vec{a}\times\vec{b}}{|\vec{a}\times\vec{b}|}$.
    \item \textbf{Distance point-to-line:} $d = \frac{|\vec{AP}\times\vec{u}|}{|\vec{u}|}$.
\end{itemize}
\end{aretenir}

\begin{savaistu}[The Vector Product in Physics and Engineering]
The cross product is ubiquitous in physics. Torque (moment of a force) is $\vec{\tau} = \vec{r} \times \vec{F}$. The magnetic force on a moving charge is $\vec{F} = q\vec{v} \times \vec{B}$. Angular momentum is $\vec{L} = \vec{r} \times \vec{p}$. In all these cases, the magnitude gives the "strength" of the effect (proportional to $\sin\theta$) and the direction gives the axis of rotation or the direction of the force, following the right-hand rule. The area interpretation is used in computer graphics to compute surface normals for lighting and shading calculations (e.g., in video games or CAD software used in Douala's engineering firms).
\end{savaistu}

\begin{exercice}[Practice Exercises]
\begin{enumerate}
    \item Given $\vec{a} = \begin{pmatrix}1\\2\\-1\end{pmatrix}$, $\vec{b} = \begin{pmatrix}3\\0\\1\end{pmatrix}$, $\vec{c} = \begin{pmatrix}-1\\1\\2\end{pmatrix}$, verify the distributive law $\vec{a} \times (\vec{b}+\vec{c}) = \vec{a}\times\vec{b} + \vec{a}\times\vec{c}$.
    \item Show that $(\vec{i}+\vec{j}) \times (\vec{j}+\vec{k}) \times (\vec{k}+\vec{i})$ depends on the placement of parentheses by evaluating both $(\vec{u}\times\vec{v})\times\vec{w}$ and $\vec{u}\times(\vec{v}\times\vec{w})$ where $\vec{u}=\vec{i}+\vec{j}$, $\vec{v}=\vec{j}+\vec{k}$, $\vec{w}=\vec{k}+\vec{i}$.
    \item Find the area of the triangle with vertices $A(1,1,0)$, $B(2,-1,2)$, $C(0,3,1)$.
    \item Find a unit vector perpendicular to the plane containing points $P(1,0,2)$, $Q(2,1,0)$, $R(0,2,1)$.
    \item The line $\ell$ passes through $A(1,2,-1)$ and has direction vector $\vec{d} = 2\vec{i} - \vec{j} + 2\vec{k}$. Find the shortest distance from the origin $O$ to $\ell$.
\end{enumerate}
\end{exercice}

\begin{corrige}[Exercise 1]
$\vec{b}+\vec{c} = \begin{pmatrix}2\\1\\3\end{pmatrix}$.
$\vec{a}\times(\vec{b}+\vec{c}) = \begin{vmatrix}\vec{i}&\vec{j}&\vec{k}\\1&2&-1\\2&1&3\end{vmatrix} = \vec{i}(6+1) - \vec{j}(3+2) + \vec{k}(1-4) = \begin{pmatrix}7\\-5\\-3\end{pmatrix}$.
$\vec{a}\times\vec{b} = \begin{vmatrix}\vec{i}&\vec{j}&\vec{k}\\1&2&-1\\3&0&1\end{vmatrix} = \vec{i}(2-0) - \vec{j}(1+3) + \vec{k}(0-6) = \begin{pmatrix}2\\-4\\-6\end{pmatrix}$.
$\vec{a}\times\vec{c} = \begin{vmatrix}\vec{i}&\vec{j}&\vec{k}\\1&2&-1\\-1&1&2\end{vmatrix} = \vec{i}(4+1) - \vec{j}(2-1) + \vec{k}(1+2) = \begin{pmatrix}5\\-1\\3\end{pmatrix}$.
Sum $= \begin{pmatrix}7\\-5\\-3\end{pmatrix}$. Verified.
\end{corrige}

\begin{corrige}[Exercise 2]
$\vec{u}\times\vec{v} = (\vec{i}+\vec{j})\times(\vec{j}+\vec{k}) = \vec{i}\times\vec{j} + \vec{i}\times\vec{k} + \vec{j}\times\vec{j} + \vec{j}\times\vec{k} = \vec{k} - \vec{j} + \vec{0} + \vec{i} = \vec{i} - \vec{j} + \vec{k}$.
$(\vec{u}\times\vec{v})\times\vec{w} = (\vec{i}-\vec{j}+\vec{k})\times(\vec{k}+\vec{i}) = \vec{i}\times\vec{k} + \vec{i}\times\vec{i} - \vec{j}\times\vec{k} - \vec{j}\times\vec{i} + \vec{k}\times\vec{k} + \vec{k}\times\vec{i} = -\vec{j} + \vec{0} - \vec{i} + \vec{k} + \vec{0} + \vec{j} = -\vec{i} + \vec{k}$.
$\vec{v}\times\vec{w} = (\vec{j}+\vec{k})\times(\vec{k}+\vec{i}) = \vec{j}\times\vec{k} + \vec{j}\times\vec{i} + \vec{k}\times\vec{k} + \vec{k}\times\vec{i} = \vec{i} - \vec{k} + \vec{0} + \vec{j} = \vec{i}+\vec{j}-\vec{k}$.
$\vec{u}\times(\vec{v}\times\vec{w}) = (\vec{i}+\vec{j})\times(\vec{i}+\vec{j}-\vec{k}) = \vec{i}\times\vec{i} + \vec{i}\times\vec{j} - \vec{i}\times\vec{k} + \vec{j}\times\vec{i} + \vec{j}\times\vec{j} - \vec{j}\times\vec{k} = \vec{0} + \vec{k} + \vec{j} - \vec{k} + \vec{0} - \vec{i} = -\vec{i} + \vec{j}$.
Since $-\vec{i}+\vec{k} \neq -\vec{i}+\vec{j}$, the results differ. Non-associativity confirmed.
\end{corrige}

\begin{corrige}[Exercise 3]
$\vec{AB} = \begin{pmatrix}1\\-2\\2\end{pmatrix}$, $\vec{AC} = \begin{pmatrix}-1\\2\\1\end{pmatrix}$.
$\vec{AB}\times\vec{AC} = \begin{vmatrix}\vec{i}&\vec{j}&\vec{k}\\1&-2&2\\-1&2&1\end{vmatrix} = \vec{i}(-2-4) - \vec{j}(1+2) + \vec{k}(2-2) = \begin{pmatrix}-6\\-3\\0\end{pmatrix}$.
$|\vec{AB}\times\vec{AC}| = \sqrt{36+9} = \sqrt{45} = 3\sqrt{5}$.
Area $= \frac{1}{2} \times 3\sqrt{5} = \frac{3\sqrt{5}}{2}$ square units.
\end{corrige}

\begin{corrige}[Exercise 4]
$\vec{PQ} = \begin{pmatrix}1\\1\\-2\end{pmatrix}$, $\vec{PR} = \begin{pmatrix}-1\\2\\-1\end{pmatrix}$.
$\vec{PQ}\times\vec{PR} = \begin{vmatrix}\vec{i}&\vec{j}&\vec{k}\\1&1&-2\\-1&2&-1\end{vmatrix} = \vec{i}(-1+4) - \vec{j}(-1-2) + \vec{k}(2+1) = \begin{pmatrix}3\\3\\3\end{pmatrix}$.
$|\vec{PQ}\times\vec{PR}| = \sqrt{27} = 3\sqrt{3}$.
Unit normal $\hat{n} = \pm \frac{1}{3\sqrt{3}}\begin{pmatrix}3\\3\\3\end{pmatrix} = \pm \frac{1}{\sqrt{3}}\begin{pmatrix}1\\1\\1\end{pmatrix}$.
\end{corrige}

\begin{corrige}[Exercise 5]
$A(1,2,-1)$, $\vec{d} = \begin{pmatrix}2\\-1\\2\end{pmatrix}$, $P = O(0,0,0)$.
$\vec{AO} = \begin{pmatrix}-1\\-2\\1\end{pmatrix}$.
$\vec{AO}\times\vec{d} = \begin{vmatrix}\vec{i}&\vec{j}&\vec{k}\\-1&-2&1\\2&-1&2\end{vmatrix} = \vec{i}(-4+1) - \vec{j}(-2-2) + \vec{k}(1+4) = \begin{pmatrix}-3\\4\\5\end{pmatrix}$.
$|\vec{AO}\times\vec{d}| = \sqrt{9+16+25} = \sqrt{50} = 5\sqrt{2}$.
$|\vec{d}| = \sqrt{4+1+4} = 3$.
$d = \frac{5\sqrt{2}}{3}$ units.
\end{corrige}

\coursec{Section 3: The Scalar Triple Product and Volumes}

In Section 2 we saw that the vector product $\vec{b} \times \vec{c}$ encodes \emph{areas}: its magnitude is the area of the parallelogram built on $\vec{b}$ and $\vec{c}$. A natural question now arises: can we measure \emph{volumes} in space with vectors? Imagine a mason in Bamenda building a concrete block shaped like a parallelepiped, with three edges given by vectors $\vec{a}$, $\vec{b}$ and $\vec{c}$ meeting at one corner. The base is a parallelogram of area $\|\vec{b} \times \vec{c}\|$, and the volume is ``base area $\times$ height''. The height is the component of $\vec{c}$ along the direction \emph{perpendicular} to the base --- and that perpendicular direction is precisely the direction of $\vec{b} \times \vec{c}$. Combining the two ideas (a vector product followed by a scalar product) gives a single number that measures volume: the \cle{scalar triple product}. This section defines it, computes it as a determinant, proves its cyclic properties, and applies it to volumes of parallelepipeds and tetrahedra and to coplanarity tests --- all classic GCE Advanced Level material.

\courssub{Definition and Determinant Formula}

\begin{definition}[Scalar triple product]
Let $\vec{a}$, $\vec{b}$, $\vec{c}$ be three vectors in space. The \cle{scalar triple product} of $\vec{a}$, $\vec{b}$, $\vec{c}$ (in this order) is the real number
\[
[\vec{a}, \vec{b}, \vec{c}] = \vec{a} \cdot (\vec{b} \times \vec{c}).
\]
It is also called the \emph{mixed product}. Note carefully: the result is a \textbf{scalar}, not a vector.
\end{definition}

The name ``mixed'' comes from the fact that we mix the two products we know: first a vector product $\vec{b} \times \vec{c}$ (a vector), then a scalar product of $\vec{a}$ with that vector (a number).

\begin{propriete}[Determinant form of the scalar triple product]
In an orthonormal basis $(\vec{i}, \vec{j}, \vec{k})$, if
\[
\vec{a} = \begin{pmatrix} a_1 \\ a_2 \\ a_3 \end{pmatrix}, \qquad
\vec{b} = \begin{pmatrix} b_1 \\ b_2 \\ b_3 \end{pmatrix}, \qquad
\vec{c} = \begin{pmatrix} c_1 \\ c_2 \\ c_3 \end{pmatrix},
\]
then
\[
\vec{a} \cdot (\vec{b} \times \vec{c}) =
\begin{vmatrix}
a_1 & a_2 & a_3 \\
b_1 & b_2 & b_3 \\
c_1 & c_2 & c_3
\end{vmatrix}.
\]
(The same value is obtained if the components are placed in \emph{columns} instead of rows.) This formula is examinable; its derivation is instructive and given below.
\end{propriete}

\textbf{Derivation.}\par
From Section 1, $\vec{b} \times \vec{c} = \begin{pmatrix} b_2 c_3 - b_3 c_2 \\ b_3 c_1 - b_1 c_3 \\ b_1 c_2 - b_2 c_1 \end{pmatrix}$. Taking the scalar product with $\vec{a}$:
\[
\vec{a} \cdot (\vec{b} \times \vec{c}) = a_1(b_2 c_3 - b_3 c_2) + a_2(b_3 c_1 - b_1 c_3) + a_3(b_1 c_2 - b_2 c_1),
\]
which is exactly the expansion of the $3 \times 3$ determinant along its first row:
\[
\begin{vmatrix}
a_1 & a_2 & a_3 \\
b_1 & b_2 & b_3 \\
c_1 & c_2 & c_3
\end{vmatrix}
= a_1 \begin{vmatrix} b_2 & b_3 \\ c_2 & c_3 \end{vmatrix}
- a_2 \begin{vmatrix} b_1 & b_3 \\ c_1 & c_3 \end{vmatrix}
+ a_3 \begin{vmatrix} b_1 & b_2 \\ c_1 & c_2 \end{vmatrix}. \qquad \blacksquare
\]

\begin{exemplebox}[First computation]
Let $\vec{a} = \begin{pmatrix} 1 \\ 2 \\ 0 \end{pmatrix}$, $\vec{b} = \begin{pmatrix} 3 \\ -1 \\ 1 \end{pmatrix}$, $\vec{c} = \begin{pmatrix} 0 \\ 2 \\ 4 \end{pmatrix}$. Then
\[
[\vec{a}, \vec{b}, \vec{c}] =
\begin{vmatrix}
1 & 2 & 0 \\
3 & -1 & 1 \\
0 & 2 & 4
\end{vmatrix}
= 1\begin{vmatrix} -1 & 1 \\ 2 & 4 \end{vmatrix}
- 2\begin{vmatrix} 3 & 1 \\ 0 & 4 \end{vmatrix}
+ 0
= 1(-4 - 2) - 2(12 - 0) = -6 - 24 = -30.
\]
The negative sign will be explained geometrically below: it means $\vec{a}$ lies on the opposite side of the plane of $(\vec{b}, \vec{c})$ from the vector $\vec{b} \times \vec{c}$.
\end{exemplebox}

\courssub{Algebraic Properties: Cyclic Invariance and Sign Changes}

\begin{propriete}[Cyclic invariance and swap rule]
For any vectors $\vec{a}, \vec{b}, \vec{c}$:
\[
\vec{a} \cdot (\vec{b} \times \vec{c}) = \vec{b} \cdot (\vec{c} \times \vec{a}) = \vec{c} \cdot (\vec{a} \times \vec{b})
\quad \text{(cyclic permutations preserve the value);}
\]
\[
\vec{a} \cdot (\vec{b} \times \vec{c}) = -\, \vec{a} \cdot (\vec{c} \times \vec{b}) = -\, \vec{b} \cdot (\vec{a} \times \vec{c})
\quad \text{(swapping any two vectors changes the sign).}
\]
Consequently one may also write $\vec{a} \cdot (\vec{b} \times \vec{c}) = (\vec{a} \times \vec{b}) \cdot \vec{c}$.
\end{propriete}

\textbf{Justification.}\par
These follow at once from the determinant form: a cyclic permutation of the three rows of a determinant can be achieved by two row swaps, so the sign changes twice and the value is unchanged; swapping two adjacent vectors swaps two rows, which changes the sign. The last identity comes from $\vec{c} \cdot (\vec{a} \times \vec{b}) = (\vec{a} \times \vec{b}) \cdot \vec{c}$ by commutativity of the scalar product.

\begin{attention}[Order matters]
The scalar triple product is \emph{not} symmetric in its three arguments: $[\vec{a}, \vec{b}, \vec{c}] = -[\vec{a}, \vec{c}, \vec{b}]$. Only \textbf{cyclic} reorderings leave it unchanged. In an exam, keep the order given in the question.
\end{attention}

\courssub{Geometric Interpretation: Volume of the Parallelepiped}

\begin{propriete}[Volume of the parallelepiped]
The volume $V$ of the parallelepiped built on three vectors $\vec{a}$, $\vec{b}$, $\vec{c}$ drawn from the same point is
\[
V = \bigl| \vec{a} \cdot (\vec{b} \times \vec{c}) \bigr|.
\]
This theorem is examinable; the proof below is instructive and frequently asked as a ``show that'' question.
\end{propriete}

\textbf{Proof.}\par
Take the parallelogram built on $\vec{b}$ and $\vec{c}$ as the base. Its area is $S = \|\vec{b} \times \vec{c}\|$ (Section 2). The vector $\vec{n} = \vec{b} \times \vec{c}$ is perpendicular to the base. The height $h$ of the parallelepiped is the absolute value of the component of $\vec{a}$ along $\vec{n}$:
\[
h = \bigl| \|\vec{a}\| \cos\theta \bigr|,
\]
where $\theta$ is the angle between $\vec{a}$ and $\vec{n}$. Hence
\[
V = S \times h = \|\vec{n}\| \, \bigl| \|\vec{a}\| \cos\theta \bigr| = \bigl| \vec{a} \cdot \vec{n} \bigr| = \bigl| \vec{a} \cdot (\vec{b} \times \vec{c}) \bigr|. \qquad \blacksquare
\]

The \emph{sign} of the triple product also carries information: it is positive when $(\vec{a}, \vec{b}, \vec{c})$ forms a right-handed triad relative to $\vec{b} \times \vec{c}$, i.e.\ when $\vec{a}$ and $\vec{b} \times \vec{c}$ point to the same side of the base plane.

\begin{popfigure}
\begin{center}
\begin{tikzpicture}[scale=1.15, line join=round]
% base parallelogram (b, c)
\coordinate (O) at (0,0);
\coordinate (B) at (3.4,0);
\coordinate (C) at (1.5,1.4);
\coordinate (BC) at (4.9,1.4);
% top vertex direction a
\coordinate (A) at (0.9,2.6);
\coordinate (AB) at (4.3,2.6);
\coordinate (AC) at (2.4,4.0);
\coordinate (ABC) at (5.8,4.0);
% shaded base
\fill[popblueL, opacity=0.7] (O) -- (B) -- (BC) -- (C) -- cycle;
% edges
\draw[thick, popdark] (O) -- (B) -- (BC) -- (C) -- cycle;
\draw[thick, popdark] (A) -- (AB) -- (ABC) -- (AC) -- cycle;
\draw[thick, popdark] (O) -- (A);
\draw[thick, popdark] (B) -- (AB);
\draw[thick, popdark] (BC) -- (ABC);
\draw[thick, popdark] (C) -- (AC);
% vectors b, c
\draw[-latex, very thick, popblue] (O) -- (B) node[midway, below] {$\vec{b}$};
\draw[-latex, very thick, popgreen] (O) -- (C) node[midway, left] {$\vec{c}$};
\draw[-latex, very thick, poporange] (O) -- (A) node[midway, left] {$\vec{a}$};
% normal vector b x c
\draw[-latex, very thick, poppurple] (O) -- (0.2,2.1) node[above left] {$\vec{b}\times\vec{c}$};
% height dashed
\draw[dashed, poppink, thick] (A) -- (0.2,2.1) node[midway, right] {$h$};
\fill (O) circle (1.5pt) node[below left] {$O$};
\end{tikzpicture}
\end{center}
\legende{Parallelepiped built on $\vec{a}, \vec{b}, \vec{c}$: base area $\|\vec{b}\times\vec{c}\|$, height $h$ = component of $\vec{a}$ along $\vec{b}\times\vec{c}$, so $V = |\vec{a}\cdot(\vec{b}\times\vec{c})|$.}
\end{popfigure}

\begin{exemplebox}[Volume of a parallelepiped]
Using the vectors of the previous example, the parallelepiped built on $\vec{a}, \vec{b}, \vec{c}$ has volume
\[
V = |{-30}| = 30 \text{ cubic units.}
\]
\end{exemplebox}

\courssub{Volume of a Tetrahedron}

A tetrahedron is a pyramid with a triangular base. If we cut a parallelepiped built on $\overrightarrow{AB}, \overrightarrow{AC}, \overrightarrow{AD}$ by the planes through its diagonal structure, the tetrahedron $ABCD$ occupies exactly one sixth of it (the triangular base has half the area of the parallelogram, and a pyramid has one third of the volume of the prism on the same base with the same height: $\tfrac12 \times \tfrac13 = \tfrac16$).

\begin{propriete}[Volume of a tetrahedron]
For four points $A, B, C, D$ in space,
\[
V_{ABCD} = \frac{1}{6} \left| \overrightarrow{AB} \cdot \left( \overrightarrow{AC} \times \overrightarrow{AD} \right) \right|.
\]
\end{propriete}

\begin{popfigure}
\begin{center}
\begin{tikzpicture}[scale=1.2, line join=round]
\coordinate (A) at (0,3.2);
\coordinate (B) at (-2.2,0);
\coordinate (C) at (2.4,0.3);
\coordinate (D) at (0.6,-1.6);
% base triangle shaded
\fill[popgreenL, opacity=0.7] (B) -- (C) -- (D) -- cycle;
% edges
\draw[thick, popdark] (B) -- (C) -- (D) -- cycle;
\draw[very thick, poporange] (A) -- (B) node[midway, left] {$\overrightarrow{AB}$};
\draw[very thick, popblue] (A) -- (C) node[midway, right] {$\overrightarrow{AC}$};
\draw[very thick, popgreen] (A) -- (D) node[midway, right, xshift=2pt] {$\overrightarrow{AD}$};
\fill (A) circle (1.6pt) node[above] {$A$};
\fill (B) circle (1.6pt) node[left] {$B$};
\fill (C) circle (1.6pt) node[right] {$C$};
\fill (D) circle (1.6pt) node[below] {$D$};
\end{tikzpicture}
\end{center}
\legende{Tetrahedron $ABCD$: the three edges from $A$ give $V = \frac{1}{6}\left|\overrightarrow{AB}\cdot\left(\overrightarrow{AC}\times\overrightarrow{AD}\right)\right|$.}
\end{popfigure}

\begin{attention}[The factor is $\tfrac{1}{6}$, not $\tfrac{1}{3}$]
A very common error is to write $V = \tfrac13 \left| \overrightarrow{AB} \cdot (\overrightarrow{AC} \times \overrightarrow{AD}) \right|$. The factor $\tfrac13$ alone would be correct for a pyramid whose base is the \emph{parallelogram} built on $\overrightarrow{AC}, \overrightarrow{AD}$. But the tetrahedron's base is the \emph{triangle} $ACD$, whose area is \emph{half} the parallelogram. Hence $V = \tfrac13 \times \text{(triangle area)} \times h = \tfrac13 \times \tfrac12 \|\overrightarrow{AC} \times \overrightarrow{AD}\| \times h = \tfrac16 \left| \overrightarrow{AB} \cdot (\overrightarrow{AC} \times \overrightarrow{AD}) \right|$.
\end{attention}

\courssub{Coplanarity Tests}

If three vectors $\vec{a}, \vec{b}, \vec{c}$ lie in one plane (are \cle{coplanar}), then the parallelepiped they build is completely flat: its volume is zero. Conversely, if the triple product is zero and the vectors are non-zero, no volume is spanned, so they must be coplanar.

\begin{propriete}[Coplanarity criteria]
\begin{itemize}
\item Three vectors $\vec{a}, \vec{b}, \vec{c}$ are coplanar if and only if $\vec{a} \cdot (\vec{b} \times \vec{c}) = 0$.
\item Four points $A, B, C, D$ are coplanar if and only if
\[
\overrightarrow{AB} \cdot \left( \overrightarrow{AC} \times \overrightarrow{AD} \right) = 0.
\]
\end{itemize}
\end{propriete}

\begin{attention}[Zero triple product does not mean zero vectors]
$[\vec{a}, \vec{b}, \vec{c}] = 0$ does \textbf{not} imply that any of the vectors is $\vec{0}$. It only means the three vectors are coplanar (or one of them is zero, or two are parallel --- all special cases of coplanarity). Never conclude ``$\vec{a} = \vec{0}$'' from a vanishing triple product.
\end{attention}

\begin{methode}[Testing coplanarity of four points]
\begin{enumerate}
\item Form the three edge vectors from one vertex: $\overrightarrow{AB}$, $\overrightarrow{AC}$, $\overrightarrow{AD}$.
\item Write the $3\times 3$ determinant of their components (rows or columns).
\item Expand the determinant (e.g.\ along a row).
\item Conclude: the points are coplanar $\iff$ the determinant equals $0$. If a parameter is involved, solve the resulting equation.
\end{enumerate}
\end{methode}

\begin{exoresolu}[Finding a parameter for coplanarity (GCE style)]
The points $A(1, 0, 2)$, $B(2, 1, 1)$, $C(0, 3, 1)$ and $D(1, 1, \lambda)$ are coplanar. Find the value(s) of $\lambda$.
\tcblower
\textbf{Solution.} Form the edge vectors from $A$:
\[
\overrightarrow{AB} = \begin{pmatrix} 1 \\ 1 \\ -1 \end{pmatrix}, \quad
\overrightarrow{AC} = \begin{pmatrix} -1 \\ 3 \\ -1 \end{pmatrix}, \quad
\overrightarrow{AD} = \begin{pmatrix} 0 \\ 1 \\ \lambda - 2 \end{pmatrix}.
\]
Coplanarity requires
\[
\begin{vmatrix}
1 & 1 & -1 \\
-1 & 3 & -1 \\
0 & 1 & \lambda - 2
\end{vmatrix} = 0.
\]
Expanding along the first row:
\begin{align*}
&1 \begin{vmatrix} 3 & -1 \\ 1 & \lambda-2 \end{vmatrix}
- 1 \begin{vmatrix} -1 & -1 \\ 0 & \lambda-2 \end{vmatrix}
+ (-1) \begin{vmatrix} -1 & 3 \\ 0 & 1 \end{vmatrix} \\
&= 1\bigl(3(\lambda-2) + 1\bigr) - 1\bigl(-(\lambda-2) - 0\bigr) - 1\bigl(-1 - 0\bigr) \\
&= 3\lambda - 5 + \lambda - 2 + 1 = 4\lambda - 6.
\end{align*}
Setting $4\lambda - 6 = 0$ gives $\boxed{\lambda = \tfrac{3}{2}}$.
\end{exoresolu}

\begin{exoresolu}[Height of a tetrahedron from volume and base area]
A tetrahedron $ABCD$ has $A(0,0,0)$, $B(2,0,0)$, $C(0,3,0)$, $D(1,1,4)$.
\begin{itemize}
\item[(a)] Compute its volume.
\item[(b)] Taking triangle $ABC$ as base, deduce the height from $D$.
\end{itemize}
\tcblower
\textbf{Solution.}
(a) $\overrightarrow{AB} = \begin{pmatrix} 2 \\ 0 \\ 0 \end{pmatrix}$, $\overrightarrow{AC} = \begin{pmatrix} 0 \\ 3 \\ 0 \end{pmatrix}$, $\overrightarrow{AD} = \begin{pmatrix} 1 \\ 1 \\ 4 \end{pmatrix}$.
\[
\overrightarrow{AB} \cdot \left( \overrightarrow{AC} \times \overrightarrow{AD} \right)
= \begin{vmatrix} 2 & 0 & 0 \\ 0 & 3 & 0 \\ 1 & 1 & 4 \end{vmatrix}
= 2 \begin{vmatrix} 3 & 0 \\ 1 & 4 \end{vmatrix} = 2 \times 12 = 24.
\]
Hence $V = \tfrac16 |24| = 4$ cubic units.

(b) Base area: $S_{ABC} = \tfrac12 \|\overrightarrow{AB} \times \overrightarrow{AC}\| = \tfrac12 \left\| \begin{pmatrix} 0 \\ 0 \\ 6 \end{pmatrix} \right\| = 3$.
Since $V = \tfrac13 S h$, we get $h = \dfrac{3V}{S} = \dfrac{3 \times 4}{3} = 4$ units. (Check: the base lies in the plane $z=0$ and $D$ has height $4$ --- consistent.)
\end{exoresolu}

\begin{exercice}[Practice]
\begin{enumerate}
\item Compute $\vec{a} \cdot (\vec{b} \times \vec{c})$ for $\vec{a} = (1, -1, 2)$, $\vec{b} = (2, 0, 1)$, $\vec{c} = (1, 1, -1)$, and state the volume of the parallelepiped they determine.
\item Show that the points $A(1,1,0)$, $B(2,0,1)$, $C(0,2,1)$, $D(3,-1,2)$ are coplanar.
\item Find $k$ so that $A(0,0,0)$, $B(1,2,0)$, $C(2,1,1)$, $D(1,1,k)$ are coplanar.
\item A tetrahedron has vertices $A(1,0,0)$, $B(0,2,0)$, $C(0,0,3)$, $D(2,2,2)$. Find its volume and the height from $D$ to the plane $ABC$ given that the area of triangle $ABC$ is $\tfrac72$.
\end{enumerate}
\end{exercice}

\begin{corrige}[Exercise 1]
\begin{enumerate}
\item $\begin{vmatrix} 1 & -1 & 2 \\ 2 & 0 & 1 \\ 1 & 1 & -1 \end{vmatrix} = 1(0-1) - (-1)(-2-1) + 2(2-0) = -1 - 3 + 4 = 0$. The volume is $0$: the three vectors are coplanar.
\item $\overrightarrow{AB} = (1,-1,1)$, $\overrightarrow{AC} = (-1,1,1)$, $\overrightarrow{AD} = (2,-2,2) = 2\overrightarrow{AB}$. The determinant has two proportional rows, hence its value is $0$. The points are coplanar.
\item $\begin{vmatrix} 1 & 2 & 0 \\ 2 & 1 & 1 \\ 1 & 1 & k \end{vmatrix} = 1(k - 1) - 2(2k - 1) + 0 = k - 1 - 4k + 2 = -3k + 1 = 0$, so $k = \tfrac13$.
\item $\overrightarrow{AB} = (-1,2,0)$, $\overrightarrow{AC} = (-1,0,3)$, $\overrightarrow{AD} = (1,2,2)$. Triple product: $-1(0\cdot 2 - 3\cdot 2) - 2((-1)(2) - 3\cdot 1) + 0 = -1(-6) - 2(-5) = 6 + 10 = 16$. Volume $V = \tfrac{16}{6} = \tfrac83$. Height: $h = \dfrac{3V}{S} = \dfrac{8}{7/2} = \dfrac{16}{7}$.
\end{enumerate}
\end{corrige}

\begin{aretenir}[Key points of Section 3]
\begin{itemize}
\item Scalar triple product: $[\vec{a},\vec{b},\vec{c}] = \vec{a} \cdot (\vec{b} \times \vec{c}) = \begin{vmatrix} a_1 & a_2 & a_3 \\ b_1 & b_2 & b_3 \\ c_1 & c_2 & c_3 \end{vmatrix}$ --- a \textbf{scalar}.
\item Cyclic permutations leave it unchanged; swapping two vectors changes the sign.
\item Volume of parallelepiped: $V = \bigl| [\vec{a},\vec{b},\vec{c}] \bigr|$; volume of tetrahedron: $V = \tfrac16 \bigl| \overrightarrow{AB} \cdot (\overrightarrow{AC} \times \overrightarrow{AD}) \bigr|$.
\item Coplanarity: $A, B, C, D$ coplanar $\iff \overrightarrow{AB} \cdot (\overrightarrow{AC} \times \overrightarrow{AD}) = 0$.
\item Traps: the tetrahedron factor is $\tfrac16$, not $\tfrac13$; a zero triple product means coplanarity, not zero vectors.
\end{itemize}
\end{aretenir}

\begin{savaistu}[Why ``triple''?]
The scalar triple product appears in physics as the volume element of crystallography: the volume of the unit cell of a crystal lattice (for instance in the study of minerals from the Cameroonian bauxite deposits) is $|\vec{a} \cdot (\vec{b} \times \vec{c})|$ where $\vec{a}, \vec{b}, \vec{c}$ are the lattice vectors. The same formula also detects whether three forces applied to a rigid body are coplanar --- a key step in statics problems.
\end{savaistu}

\coursec{Section 4: Equations of a Plane — Vector and Cartesian Forms}

In Section 3 we saw how the scalar triple product $\vec{a}\cdot(\vec{b}\times\vec{c})$ gives the volume of a parallelepiped. That same cross product $\vec{b}\times\vec{c}$ produces a vector perpendicular to both $\vec{b}$ and $\vec{c}$ — a \emph{normal vector} to the plane containing them. This simple observation is the key to writing down the equation of any plane in space. Whether we are given a point and a normal, three points, or a point and two directions, the cross product provides the normal we need to switch effortlessly between vector (parametric) and Cartesian forms. Mastering these conversions is a staple of GCE Advanced Level Further Mathematics; they appear in almost every paper, often as the first step in a longer problem on lines, planes, distances or angles.

\courssub{Specifying a Plane in Space}

A plane in three‑dimensional space can be uniquely determined in three equivalent ways:

\begin{itemize}
\item A point $A$ (with position vector $\vec{a}$) and a \textbf{normal vector} $\vec{n}\neq\vec{0}$ perpendicular to the plane.
\item A point $A$ and two \textbf{non‑parallel direction vectors} $\vec{u}$ and $\vec{v}$ lying in the plane.
\item Three \textbf{non‑collinear points} $A$, $B$, $C$ (the vectors $\vec{AB}$ and $\vec{AC}$ then serve as the two direction vectors).
\end{itemize}

The cross product $\vec{u}\times\vec{v}$ immediately gives a normal vector $\vec{n}$ when we have the two‑direction description. This links the three descriptions together.

\courssub{Vector (Parametric) Equation of a Plane}

If a plane passes through point $A(\vec{a})$ and is parallel to the non‑parallel vectors $\vec{u}$ and $\vec{v}$, then any point $R(\vec{r})$ in the plane can be reached by starting at $A$ and moving some distance along $\vec{u}$ and some distance along $\vec{v}$.

\begin{definition}[Vector (Parametric) Equation of a Plane]
The vector equation of a plane through $A(\vec{a})$ with direction vectors $\vec{u}$ and $\vec{v}$ ($\vec{u}\nparallel\vec{v}$) is
\[
\vec{r} = \vec{a} + s\vec{u} + t\vec{v},\qquad s,t\in\mathbb{R}.
\]
The parameters $s$ and $t$ are independent real numbers; varying them sweeps out the whole plane.
\end{definition}

\begin{attention}[Two Parameters are Essential]
A single parameter $s$ in $\vec{r}=\vec{a}+s\vec{u}$ describes a \textbf{line}, not a plane. A plane requires \textbf{two} independent parameters because it is a two‑dimensional surface. Forgetting the second parameter is a very common error in examinations.
\end{attention}

\begin{popfigure}
\begin{tikzpicture}[scale=1.2, >=stealth]
% Coordinate axes
\draw[->, thin, gray] (-0.5,0,0) -- (3,0,0) node[right] {$x$};
\draw[->, thin, gray] (0,-0.5,0) -- (0,3,0) node[above] {$y$};
\draw[->, thin, gray] (0,0,-0.5) -- (0,0,3) node[above] {$z$};
% Plane as a parallelogram
\coordinate (O) at (0,0,0);
\coordinate (A) at (1,0.5,0.2);
\coordinate (U) at (2,0.2,0.8);
\coordinate (V) at (0.5,2,0.3);
\coordinate (AplusU) at ($(A)+(U)$);
\coordinate (AplusV) at ($(A)+(V)$);
\coordinate (AplusUplusV) at ($(A)+(U)+(V)$);
% Plane surface
\fill[popblueL, opacity=0.3] (A) -- (AplusU) -- (AplusUplusV) -- (AplusV) -- cycle;
\draw[popblue, thick] (A) -- (AplusU) -- (AplusUplusV) -- (AplusV) -- cycle;
% Vectors
\draw[->, poporange, thick] (O) -- (A) node[midway, left] {$\vec{a}$};
\draw[->, popgreen, thick] (A) -- (AplusU) node[midway, above right] {$\vec{u}$};
\draw[->, popgreen, thick] (A) -- (AplusV) node[midway, left] {$\vec{v}$};
% Normal vector n = u x v
\coordinate (N) at ($(A)+(0.5,0.5,1.5)$);
\draw[->, poppurple, very thick] (A) -- (N) node[midway, right] {$\vec{n}=\vec{u}\times\vec{v}$};
% Generic point R
\coordinate (R) at ($(A)+0.6*(U)+0.4*(V)$);
\draw[->, poppink, thick] (O) -- (R) node[midway, above left] {$\vec{r}$};
\draw[->, poppink, thick] (A) -- (R) node[midway, below] {$s\vec{u}+t\vec{v}$};
% Points labels
\node[popdark] at (A) [below left] {$A$};
\node[popdark] at (R) [above right] {$R$};
\end{tikzpicture}
\legende{Vector equation $\vec{r}=\vec{a}+s\vec{u}+t\vec{v}$: the plane is spanned by $\vec{u}$ and $\vec{v}$ from point $A$; the normal $\vec{n}=\vec{u}\times\vec{v}$ is perpendicular to the plane.}
\end{popfigure}

\courssub{Normal Form and Cartesian Equation}

If we know a point $A(\vec{a})$ on the plane and a normal vector $\vec{n}=\begin{pmatrix}a\\b\\c\end{pmatrix}$, the condition that $\vec{r}=\begin{pmatrix}x\\y\\z\end{pmatrix}$ lies in the plane is simply that the vector $\vec{r}-\vec{a}$ is perpendicular to $\vec{n}$:
\[
(\vec{r}-\vec{a})\cdot\vec{n}=0.
\]

\begin{propriete}[Normal Form and Cartesian Equation]
Let $\vec{n}=\begin{pmatrix}a\\b\\c\end{pmatrix}$ be a normal vector to a plane $\Pi$, and let $A(x_0,y_0,z_0)$ be a point on $\Pi$. Then $\Pi$ has equation
\[
(\vec{r}-\vec{a})\cdot\vec{n}=0 \quad\Longleftrightarrow\quad a(x-x_0)+b(y-y_0)+c(z-z_0)=0.
\]
Expanding gives the \textbf{Cartesian equation}
\[
ax+by+cz = d,\qquad\text{where } d = ax_0+by_0+cz_0 = \vec{n}\cdot\vec{a}.
\]
\end{propriete}

\begin{attention}[Coefficients Give the Normal]
In the Cartesian equation $ax+by+cz=d$, the coefficients $(a,b,c)$ are the components of a \textbf{normal vector} to the plane. They are \textbf{not} direction vectors lying in the plane. A direction vector $\vec{v}$ in the plane satisfies $a v_x+b v_y+c v_z=0$.
\end{attention}

\courssub{Finding the Cartesian Equation from Three Points}

When three non‑collinear points $A$, $B$, $C$ are given, we form two vectors in the plane, e.g. $\vec{AB}$ and $\vec{AC}$, and take their cross product to obtain a normal vector.

\begin{methode}[Cartesian Equation Through Three Points]
\begin{enumerate}
\item Compute $\vec{AB}=\vec{b}-\vec{a}$ and $\vec{AC}=\vec{c}-\vec{a}$.
\item Find a normal vector $\vec{n}=\vec{AB}\times\vec{AC}$.
\item Use point $A$ (or $B$ or $C$) to find $d=\vec{n}\cdot\vec{a}$.
\item Write $ax+by+cz=d$ where $\vec{n}=\begin{pmatrix}a\\b\\c\end{pmatrix}$.
\end{enumerate}
\end{methode}

\begin{exoresolu}[Plane Through Three Points]
Find the Cartesian equation of the plane passing through $A(1,2,-1)$, $B(3,0,2)$ and $C(-1,1,3)$.
\tcblower
\textbf{Solution.}
\[
\vec{AB}=\begin{pmatrix}2\\-2\\3\end{pmatrix},\quad
\vec{AC}=\begin{pmatrix}-2\\-1\\4\end{pmatrix}.
\]
\[
\vec{n}=\vec{AB}\times\vec{AC}=
\begin{vmatrix}
\mathbf{i} & \mathbf{j} & \mathbf{k}\\
2 & -2 & 3\\
-2 & -1 & 4
\end{vmatrix}
= \mathbf{i}((-2)(4)-3(-1)) - \mathbf{j}((2)(4)-3(-2)) + \mathbf{k}((2)(-1)-(-2)(-2))
\]
\[
= \mathbf{i}(-8+3) - \mathbf{j}(8+6) + \mathbf{k}(-2-4) = \begin{pmatrix}-5\\-14\\-6\end{pmatrix}.
\]
We can multiply by $-1$ for a simpler normal: $\vec{n}=\begin{pmatrix}5\\14\\6\end{pmatrix}$.
Using point $A(1,2,-1)$:
\[
d = 5(1)+14(2)+6(-1) = 5+28-6 = 27.
\]
The Cartesian equation is
\[
\boxed{5x+14y+6z=27}.
\]
(Check: $B$ gives $5(3)+14(0)+6(2)=15+12=27$; $C$ gives $5(-1)+14(1)+6(3)=-5+14+18=27$.)
\end{exoresolu}

\courssub{Converting Between Parametric and Cartesian Forms}

The cross product is the bridge between the two forms.

\begin{methode}[Parametric $\to$ Cartesian]
Given $\vec{r}=\vec{a}+s\vec{u}+t\vec{v}$:
\begin{enumerate}
\item Compute $\vec{n}=\vec{u}\times\vec{v}$.
\item Compute $d=\vec{n}\cdot\vec{a}$.
\item Write $\vec{n}\cdot\vec{r}=d$, i.e. $ax+by+cz=d$.
\end{enumerate}
\end{methode}

\begin{methode}[Cartesian $\to$ Parametric]
Given $ax+by+cz=d$ with normal $\vec{n}=\begin{pmatrix}a\\b\\c\end{pmatrix}$:
\begin{enumerate}
\item Find \textbf{one point} $A$ on the plane (e.g. set two variables to $0$ and solve for the third, provided the corresponding coefficient is non‑zero).
\item Find \textbf{two independent direction vectors} $\vec{u},\vec{v}$ perpendicular to $\vec{n}$ (i.e. $\vec{n}\cdot\vec{u}=0$, $\vec{n}\cdot\vec{v}=0$). A quick way: if $a\neq0$, take $\vec{u}=\begin{pmatrix}-b\\a\\0\end{pmatrix}$ and $\vec{v}=\begin{pmatrix}-c\\0\\a\end{pmatrix}$; adjust if some coefficients are zero.
\item Write $\vec{r}=\vec{a}+s\vec{u}+t\vec{v}$.
\end{enumerate}
\end{methode}

\begin{exoresolu}[Converting Both Ways]
\begin{enumerate}
\item Convert $\vec{r}=\begin{pmatrix}1\\-2\\3\end{pmatrix}+s\begin{pmatrix}2\\1\\0\end{pmatrix}+t\begin{pmatrix}-1\\0\\1\end{pmatrix}$ to Cartesian form.
\item Convert $2x-y+3z=6$ to parametric form.
\end{enumerate}
\tcblower
\textbf{Solution.}
\begin{enumerate}
\item $\vec{u}=\begin{pmatrix}2\\1\\0\end{pmatrix},\ \vec{v}=\begin{pmatrix}-1\\0\\1\end{pmatrix}$.
$\vec{n}=\vec{u}\times\vec{v}=\begin{vmatrix}\mathbf{i}&\mathbf{j}&\mathbf{k}\\2&1&0\\-1&0&1\end{vmatrix}
= \mathbf{i}(1\cdot1-0\cdot0) - \mathbf{j}(2\cdot1-0\cdot(-1)) + \mathbf{k}(2\cdot0-1\cdot(-1))
= \begin{pmatrix}1\\-2\\1\end{pmatrix}$.
$d=\vec{n}\cdot\vec{a}=1(1)+(-2)(-2)+1(3)=1+4+3=8$.
Cartesian equation: $\boxed{x-2y+z=8}$.
\item Normal $\vec{n}=\begin{pmatrix}2\\-1\\3\end{pmatrix}$. Find a point: set $y=0,z=0 \Rightarrow 2x=6 \Rightarrow x=3$, so $A(3,0,0)$.
Two direction vectors perpendicular to $\vec{n}$:
$\vec{u}=\begin{pmatrix}1\\2\\0\end{pmatrix}$ (since $2(1)+(-1)(2)+3(0)=0$),
$\vec{v}=\begin{pmatrix}0\\3\\1\end{pmatrix}$ (since $2(0)+(-1)(3)+3(1)=0$).
Parametric form: $\boxed{\vec{r}=\begin{pmatrix}3\\0\\0\end{pmatrix}+s\begin{pmatrix}1\\2\\0\end{pmatrix}+t\begin{pmatrix}0\\3\\1\end{pmatrix},\ s,t\in\mathbb{R}}$.
\end{enumerate}
\end{exoresolu}

\courssub{Perpendicular Distance from a Point to a Plane (GCE Extension)}

The shortest distance from a point $P(x_1,y_1,z_1)$ to the plane $\Pi: ax+by+cz=d$ is the length of the projection of $\vec{AP}$ onto the normal $\vec{n}$, where $A$ is any point on $\Pi$.

\begin{propriete}[Distance from Point to Plane]
The perpendicular distance from $P(x_1,y_1,z_1)$ to the plane $ax+by+cz=d$ is
\[
\text{dist}(P,\Pi)=\frac{|ax_1+by_1+cz_1-d|}{\sqrt{a^2+b^2+c^2}}.
\]
\end{propriete}

\begin{exoresolu}[Distance from Point to Plane]
Find the distance from $P(2,-1,3)$ to the plane $3x-4y+12z=13$.
\tcblower
Substitute into the formula:
\[
\text{Numerator}=|3(2)-4(-1)+12(3)-13| = |6+4+36-13| = 33.
\]
\[
\text{Denominator}=\sqrt{3^2+(-4)^2+12^2}=\sqrt{9+16+144}=\sqrt{169}=13.
\]
\[
\text{Distance}=\frac{33}{13}\ \text{units}.
\]
\end{exoresolu}

\begin{popfigure}
\begin{tikzpicture}[scale=1.2, >=stealth]
% Plane: 2x - y + z = 4
% Point A(2,0,0) on plane
% Direction vectors u = (-2,-4,0) [= B-A], v = (-2,0,4) [= C-A]
\coordinate (O) at (0,0,0);
\coordinate (A) at (2,0,0);
\coordinate (U) at (-2,-4,0);
\coordinate (V) at (-2,0,4);
\coordinate (AplusU) at ($(A)+(U)$);
\coordinate (AplusV) at ($(A)+(V)$);
\coordinate (AplusUplusV) at ($(A)+(U)+(V)$);
% Plane surface
\fill[popgreenL, opacity=0.2] (A) -- (AplusU) -- (AplusUplusV) -- (AplusV) -- cycle;
\draw[popgreen, thick] (A) -- (AplusU) -- (AplusUplusV) -- (AplusV) -- cycle;
% Normal vector n = (2,-1,1) scaled
\coordinate (N) at ($(A)+(1,-0.5,0.5)$);
\draw[->, poppurple, very thick] (A) -- (N) node[midway, right] {$\vec{n}=(a,b,c)$};
% External point P(3,2,2)
\coordinate (P) at (3,2,2);
\draw[->, poporange, thick] (O) -- (P) node[midway, above left] {$\vec{p}$};
\node[popdark] at (P) [above right] {$P$};
% Foot of perpendicular F: P + t*n, t = -1/3 => F = (7/3, 7/3, 5/3) ≈ (2.333, 2.333, 1.667)
\coordinate (F) at (2.333, 2.333, 1.667);
\draw[->, poppink, thick, dashed] (P) -- (F) node[midway, left] {$\text{distance}$};
\draw[poppink, thick] (A) -- (F) node[midway, below] {$\vec{a}$};
\node[popdark] at (A) [below left] {$A$};
\node[popdark] at (F) [below right] {$F$};
% Axes
\draw[->, thin, gray] (-0.5,0,0) -- (4,0,0) node[right] {$x$};
\draw[->, thin, gray] (0,-1,0) -- (0,3,0) node[above] {$y$};
\draw[->, thin, gray] (0,0,-0.5) -- (0,0,4) node[above] {$z$};
\end{tikzpicture}
\legende{Perpendicular distance from point $P$ to plane $\Pi: ax+by+cz=d$. The distance is the length of the projection of $\vec{AP}$ onto the normal $\vec{n}$. The foot $F$ is correctly computed as the intersection of the line through $P$ parallel to $\vec{n}$ with the plane.}
\end{popfigure}

\courssub{Angle Between Two Planes (GCE Extension)}

The angle between two planes is defined as the angle between their normal vectors (or its supplement, whichever is acute).

\begin{propriete}[Angle Between Two Planes]
Let $\Pi_1: a_1x+b_1y+c_1z=d_1$ and $\Pi_2: a_2x+b_2y+c_2z=d_2$ have normals $\vec{n}_1=\begin{pmatrix}a_1\\b_1\\c_1\end{pmatrix}$, $\vec{n}_2=\begin{pmatrix}a_2\\b_2\\c_2\end{pmatrix}$. The acute angle $\theta$ between the planes satisfies
\[
\cos\theta = \frac{|\vec{n}_1\cdot\vec{n}_2|}{|\vec{n}_1|\,|\vec{n}_2|}
= \frac{|a_1a_2+b_1b_2+c_1c_2|}{\sqrt{a_1^2+b_1^2+c_1^2}\,\sqrt{a_2^2+b_2^2+c_2^2}}.
\]
\begin{itemize}
\item $\Pi_1\parallel\Pi_2 \iff \vec{n}_1\parallel\vec{n}_2 \iff (a_1,b_1,c_1)=k(a_2,b_2,c_2)$ for some $k\neq0$.
\item $\Pi_1\perp\Pi_2 \iff \vec{n}_1\perp\vec{n}_2 \iff a_1a_2+b_1b_2+c_1c_2=0$.
\end{itemize}
\end{propriete}

\begin{exoresolu}[Angle Between Planes]
Find the acute angle between the planes $2x-y+2z=5$ and $3x+6y-2z=7$.
\tcblower
$\vec{n}_1=\begin{pmatrix}2\\-1\\2\end{pmatrix},\ \vec{n}_2=\begin{pmatrix}3\\6\\-2\end{pmatrix}$.
$\vec{n}_1\cdot\vec{n}_2 = 2(3)+(-1)(6)+2(-2)=6-6-4=-4$.
$|\vec{n}_1|=\sqrt{4+1+4}=3,\quad |\vec{n}_2|=\sqrt{9+36+4}=7$.
$\cos\theta = \frac{|-4|}{3\cdot7}=\frac{4}{21}$.
$\theta = \cos^{-1}\left(\frac{4}{21}\right) \approx 79.0^\circ$.
\end{exoresolu}

\courssub{Summary and Examination Tips}

\begin{aretenir}[Key Points for GCE Advanced Level]
\begin{itemize}
\item \textbf{Vector form:} $\vec{r}=\vec{a}+s\vec{u}+t\vec{v}$ ($s,t\in\mathbb{R}$, $\vec{u}\nparallel\vec{v}$).
\item \textbf{Normal form:} $(\vec{r}-\vec{a})\cdot\vec{n}=0$.
\item \textbf{Cartesian form:} $ax+by+cz=d$ with $\vec{n}=\begin{pmatrix}a\\b\\c\end{pmatrix}$.
\item \textbf{Normal from two directions:} $\vec{n}=\vec{u}\times\vec{v}$.
\item \textbf{Plane through three points:} $\vec{n}=\vec{AB}\times\vec{AC}$, then $d=\vec{n}\cdot\vec{a}$.
\item \textbf{Distance point–plane:} $\dfrac{|ax_1+by_1+cz_1-d|}{\sqrt{a^2+b^2+c^2}}$.
\item \textbf{Angle between planes:} $\cos\theta=\dfrac{|\vec{n}_1\cdot\vec{n}_2|}{|\vec{n}_1||\vec{n}_2|}$.
\item \textbf{Parallel planes:} normals are parallel. \textbf{Perpendicular planes:} normals are perpendicular.
\end{itemize}
\end{aretenir}

\begin{attention}[Common Examination Pitfalls]
\begin{itemize}
\item Forgetting that the parametric form needs \textbf{two} parameters.
\item Confusing the normal vector $(a,b,c)$ with a direction vector in the plane.
\item When converting Cartesian $\to$ parametric, choosing direction vectors that are \textbf{not} perpendicular to the normal (check $\vec{n}\cdot\vec{u}=0$).
\item In the distance formula, using the wrong sign for $d$ (remember $d$ is the constant on the right‑hand side of $ax+by+cz=d$).
\item For the angle between planes, always take the \textbf{absolute value} of the dot product to get the acute angle.
\end{itemize}
\end{attention}

\begin{savaistu}[Planes in Real Life]
The roof of a traditional Cameroonian \emph{case} (round hut) is a conical surface, but the flat ceiling of a modern classroom is a plane. The equation $ax+by+cz=d$ describes that ceiling: the normal vector $(a,b,c)$ points straight up (or down), and $d$ fixes the height. In computer graphics, every flat polygon on your screen is a plane segment defined by exactly these equations — the cross product of two edges gives the normal used for lighting calculations.
\end{savaistu}

With this section we complete Chapter FMA14. You now have a full toolkit: vector products for areas, scalar triple products for volumes, and the equations of planes for describing flat surfaces in space. These ideas are the foundation of three‑dimensional analytic geometry and appear repeatedly in mechanics, computer graphics, and higher mathematics. Practice converting between forms until it becomes automatic — it is the single most useful skill for the GCE paper.

\newpage

\coursec{Integration activity}

\courssub{Aims of the activity}

By the end of this group activity, you should be able to:

\begin{itemize}
\item compute a \cle{vector product} $\vec{AB}\times\vec{AC}$ from coordinates and interpret it geometrically as a normal vector and as an area;
\item derive the \cle{Cartesian equation of a plane} from a point and a normal vector;
\item use the \cle{scalar triple product} to test whether four points are coplanar and to compute the \cle{volume of a tetrahedron};
\item compute the \cle{perpendicular distance from a point to a plane} and interpret it in a real safety context;
\item work cooperatively, present a reasoned solution on the board, and compare the efficiency of different methods.
\end{itemize}

\begin{aretenir}[Toolbox for this activity]
You will need the following results from this chapter:
\begin{itemize}
\item For $\vec{u}(x,y,z)$ and $\vec{v}(x',y',z')$,
\[
\vec{u}\times\vec{v}=\begin{pmatrix} yz'-zy' \\ zx'-xz' \\ xy'-yx' \end{pmatrix}.
\]
\item $\vec{u}\times\vec{v}$ is perpendicular to both $\vec{u}$ and $\vec{v}$, and $|\vec{u}\times\vec{v}|$ is the area of the parallelogram built on $\vec{u}$ and $\vec{v}$; the triangle built on them has area $\tfrac12|\vec{u}\times\vec{v}|$.
\item Plane through $A$ with normal vector $\vec{n}(a,b,c)$: $ax+by+cz+d=0$, where $d$ is found by substituting the coordinates of $A$.
\item $[\vec{u},\vec{v},\vec{w}]=(\vec{u}\times\vec{v})\cdot\vec{w}$; four points $A,B,C,D$ are coplanar if and only if $[\vec{AB},\vec{AC},\vec{AD}]=0$; the volume of the tetrahedron $ABCD$ is $\tfrac{1}{6}\bigl|[\vec{AB},\vec{AC},\vec{AD}]\bigr|$.
\item Distance from $P(x_0,y_0,z_0)$ to the plane $ax+by+cz+d=0$:
\[
d=\dfrac{|ax_0+by_0+cz_0+d|}{\sqrt{a^2+b^2+c^2}}.
\]
\end{itemize}
\end{aretenir}

\courssub{Situation}

\begin{experience}[Designing a Roof Truss for a School Hall in Buea]
A contractor in Buea is roofing the new assembly hall of a secondary school. Because of the heavy rains on the slopes of Mount Cameroon, the roof must be steep enough for water to run off, and the carpenter wants exact measurements before cutting the timber and ordering the zinc sheets.

The contractor sets up a 3D coordinate system $(O;\vec{i},\vec{j},\vec{k})$, with the unit $1$ metre on each axis, fixed to the hall under construction. One rectangular face of the sloping roof is modelled by the plane through the three points
\[
A(0,0,4), \qquad B(6,0,4), \qquad C(6,8,6),
\]
where $A$ and $B$ are the two ends of the lower edge of the roof face (at height $4$ m) and $C$ is one end of the upper edge (at height $6$ m). The triangular panel $ABC$ is half of that rectangular face and will be covered with corrugated zinc sheets, sold by square metre.

A wooden beam is to be placed at the fourth corner $D(0,8,6)$, which the contractor \emph{believes} completes the rectangle $ABDC$ of the roof face. Before nailing anything, the carpenter must check that $A$, $B$, $C$ and $D$ really lie in the same plane; if they do, the ``attic space'' $ABCD$ is flat and the architect must raise the point $D$ to create storage volume.

Finally, a lamp hangs from the ceiling framework at the point $P(3,4,2)$. Safety regulations require a clearance of at least $2.5$ m between the lamp and the roof plane directly above it. The contractor asks your group to verify whether the regulation is satisfied.

Your group must produce a complete written report, then present it on the board.
\end{experience}

\begin{popfigure}
\begin{center}
\begin{tikzpicture}[scale=0.62, line join=round]
\coordinate (A) at (0,0);
\coordinate (B) at (5.2,0);
\coordinate (C) at (7.0,3.4);
\coordinate (D) at (1.8,3.4);
\coordinate (P) at (3.6,1.4);
\fill[popblueL,opacity=0.75] (A)--(B)--(C)--(D)--cycle;
\draw[popdark,thick] (A)--(B)--(C)--(D)--cycle;
\draw[poporange,very thick] (A)--(B)--(C)--cycle;
\draw[popgreen,very thick,->] (3.6,1.7) -- (3.6,4.1);
\fill[popink] (A) circle (2.2pt) node[below left]{$A$};
\fill[popink] (B) circle (2.2pt) node[below right]{$B$};
\fill[popink] (C) circle (2.2pt) node[above right]{$C$};
\fill[popink] (D) circle (2.2pt) node[above left]{$D$};
\fill[poppurple] (P) circle (2.2pt) node[below]{$P$};
\draw[poppurple,dashed] (P) -- (3.6,2.6);
\node[popgreen] at (4.35,3.1) {$\vec{n}$};
\end{tikzpicture}
\end{center}
\legende{The roof face $ABDC$, the triangular panel $ABC$, a normal vector $\vec{n}$ and the lamp at $P$.}
\end{popfigure}

\courssub{Tasks}

Work in groups of three or four. Answer the questions in order; each one uses a different tool of the chapter.

\begin{enumerate}
\item \textbf{A normal vector to the roof.} Compute the vectors $\vec{AB}$ and $\vec{AC}$. Calculate the vector product $\vec{AB}\times\vec{AC}$. Explain in one sentence why this vector is normal to the roof plane.
\item \textbf{Cartesian equation of the roof plane.} Using point $A$ and your normal vector, show that the roof plane has Cartesian equation $y-4z+16=0$. Verify that $B$ and $C$ also satisfy this equation.
\item \textbf{Area of the zinc panel.} Compute $|\vec{AB}\times\vec{AC}|$ and deduce the exact area of the triangular panel $ABC$, then an approximate value to $2$ decimal places. If zinc sheets cost $4\,500$ FCFA per square metre and the whole rectangular face is twice the triangle, estimate the cost for this face (round up to the next square metre).
\item \textbf{Is the beam $D$ in the right place?} Compute the scalar triple product $[\vec{AB},\vec{AC},\vec{AD}]$ where $D(0,8,6)$. What do you conclude about the four points? What is the volume of the ``tetrahedron'' $ABCD$? Interpret this bad news for the carpenter.
\item \textbf{Creating an attic.} The architect moves the fourth corner to $D'(0,8,8)$. Compute the new scalar triple product and the volume of the tetrahedron $ABCD'$. Is the attic now usable?
\item \textbf{Safety clearance of the lamp.} Compute the perpendicular distance from $P(3,4,2)$ to the roof plane. Is the regulation clearance of $2.5$ m respected?
\item \textbf{Presentation.} Prepare a poster showing: the normal vector, the plane equation, the panel area, the coplanarity test, the attic volume and the clearance. Compare with another group: did anyone use $\vec{AC}\times\vec{AB}$ instead? What changes, what does not?
\end{enumerate}

\begin{methode}[How to attack each question]
\begin{enumerate}
\item \textbf{Normal vector:} form $\vec{AB}$ and $\vec{AC}$ by subtracting coordinates, then apply the cross product formula component by component.
\item \textbf{Plane equation:} simplify the normal vector if possible (divide by a common factor), write $a(x-x_A)+b(y-y_A)+c(z-z_A)=0$, expand, and check the other given points.
\item \textbf{Area:} take the modulus of the cross product (parallelogram area), then halve it for the triangle.
\item \textbf{Coplanarity and volume:} dot the cross product with the third vector; zero means coplanar, otherwise take $\tfrac16$ of the absolute value.
\item \textbf{Distance:} substitute the point into the left-hand side of the plane equation, take the absolute value, and divide by the norm of the normal vector.
\end{enumerate}
\end{methode}

\begin{attention}[Traps to avoid in this activity]
\begin{itemize}
\item The vector product is \emph{anti-commutative}: $\vec{AC}\times\vec{AB}=-(\vec{AB}\times\vec{AC})$. The sign of the normal vector changes, but the plane equation and the areas do not.
\item Do not forget the factor $\tfrac12$ for the triangle and $\tfrac16$ for the tetrahedron.
\item A zero scalar triple product does not mean you made an error: it \emph{means} the points are coplanar.
\item In the distance formula, use the \emph{normalised} denominator $\sqrt{a^2+b^2+c^2}$; forgetting it gives a wrong clearance.
\item When simplifying a normal vector by dividing by a common factor, the plane equation changes form but describes the same plane; always re-check with a known point.
\end{itemize}
\end{attention}

\courssub{Model answer}

\textbf{1. Normal vector.} From the coordinates,
\[
\vec{AB}=(6-0,\,0-0,\,4-4)=(6,0,0), \qquad \vec{AC}=(6-0,\,8-0,\,6-4)=(6,8,2).
\]
Applying the cross product formula with $\vec{AB}=(6,0,0)$ and $\vec{AC}=(6,8,2)$:
\[
\vec{AB}\times\vec{AC}=
\begin{pmatrix} (0)(2)-(0)(8) \\ (0)(6)-(6)(2) \\ (6)(8)-(0)(6) \end{pmatrix}
=
\begin{pmatrix} 0 \\ -12 \\ 48 \end{pmatrix}.
\]
By definition, the vector product is perpendicular to both $\vec{AB}$ and $\vec{AC}$, which lie in the roof plane; hence $\vec{n}=(0,-12,48)$ is normal to the roof plane. We may simplify to $\vec{n}=(0,-1,4)$ (dividing by $12$), since any non-zero scalar multiple of a normal vector is still normal to the plane.

\textbf{2. Equation of the plane.} A point $M(x,y,z)$ belongs to the plane if and only if $\vec{AM}\cdot\vec{n}=0$. With $\vec{AM}=(x-0,\,y-0,\,z-4)$ and $\vec{n}=(0,-1,4)$:
\[
(0)(x-0)+(-1)(y-0)+(4)(z-4)=0
\;\Longleftrightarrow\;
-y+4z-16=0
\;\Longleftrightarrow\;
y-4z+16=0.
\]
Check: $B(6,0,4)$: $0-16+16=0$ \checkmark; $C(6,8,6)$: $8-24+16=0$ \checkmark. Both points lie in the plane, as expected.

\textbf{3. Area of the panel.}
\[
|\vec{AB}\times\vec{AC}|=\sqrt{0^2+(-12)^2+48^2}=\sqrt{144+2304}=\sqrt{2448}=\sqrt{144\times 17}=12\sqrt{17}.
\]
This is the area of the parallelogram built on $\vec{AB}$ and $\vec{AC}$, so
\[
\text{Area}(ABC)=\tfrac12\,|\vec{AB}\times\vec{AC}|=6\sqrt{17}\approx 24.74\ \mathrm{m^2}.
\]
The full rectangular face has area $12\sqrt{17}\approx 49.48\ \mathrm{m^2}$, so the contractor orders $50\ \mathrm{m^2}$ of zinc, costing about $50\times 4\,500 = 225\,000$ FCFA.

\textbf{4. Coplanarity test for $D(0,8,6)$.} We have $\vec{AD}=(0-0,\,8-0,\,6-4)=(0,8,2)$, and
\[
[\vec{AB},\vec{AC},\vec{AD}]=(\vec{AB}\times\vec{AC})\cdot\vec{AD}
=(0)(0)+(-12)(8)+(48)(2)=-96+96=0.
\]
The scalar triple product is zero, so $A,B,C,D$ are \cle{coplanar}: indeed $D$ satisfies the plane equation ($8-24+16=0$). The ``tetrahedron'' $ABCD$ is flat:
\[
V=\tfrac16\bigl|[\vec{AB},\vec{AC},\vec{AD}]\bigr|=0\ \mathrm{m^3}.
\]
The beam at $D$ lies exactly in the roof plane: the carpenter's belief is correct, but there is no attic volume at all.

\textbf{5. Attic with $D'(0,8,8)$.} Now $\vec{AD'}=(0,8,4)$ and
\[
[\vec{AB},\vec{AC},\vec{AD'}]=(0)(0)+(-12)(8)+(48)(4)=-96+192=96.
\]
\[
V_{ABCD'}=\tfrac16|96|=16\ \mathrm{m^3}.
\]
The volume is non-zero: raising the corner to height $8$ m creates an attic of $16\ \mathrm{m^3}$, which is usable storage space.

\textbf{6. Clearance of the lamp.} With the plane $y-4z+16=0$ (so $a=0$, $b=1$, $c=-4$, $d=16$) and $P(3,4,2)$:
\[
d=\frac{|(0)(3)+(1)(4)+(-4)(2)+16|}{\sqrt{0^2+1^2+(-4)^2}}=\frac{|4-8+16|}{\sqrt{17}}=\frac{12}{\sqrt{17}}\approx 2.91\ \mathrm{m}.
\]
Since $2.91 > 2.5$, the safety clearance is respected.

\textbf{7. Comparison of methods.} Using $\vec{AC}\times\vec{AB}=(0,12,-48)$ gives the opposite normal vector, hence the equation $-y+4z-16=0$, which is the same plane (both sides multiplied by $-1$); areas, volumes (in absolute value) and distances are unchanged.

\begin{exercice}[Going further: a second roof face]
The opposite face of the roof passes through $A'(0,0,4)$, $E(0,8,6)$ and $F(-6,8,6)$.
\begin{enumerate}
\item Find a normal vector to this face and its Cartesian equation.
\item Compute the area of the triangular panel $A'EF$.
\item Show that the two roof planes $y-4z+16=0$ and the plane you found intersect along a line parallel to $\vec{i}$, and interpret this line on the building.
\end{enumerate}
\end{exercice}

\begin{corrige}[Going further: a second roof face]
\textbf{1.} $\vec{A'E}=(0,8,2)$ and $\vec{A'F}=(-6,8,2)$, so
\[
\vec{A'E}\times\vec{A'F}=
\begin{pmatrix} (8)(2)-(2)(8) \\ (2)(-6)-(0)(2) \\ (0)(8)-(8)(-6) \end{pmatrix}
=
\begin{pmatrix} 0 \\ -12 \\ 48 \end{pmatrix},
\]
the same normal direction $(0,-1,4)$. The plane through $A'(0,0,4)$ is $-y+4(z-4)=0$, i.e. $y-4z+16=0$: the two faces are in fact coplanar in this model (a single sloping face).
\[
\textbf{2.}\quad \text{Area}(A'EF)=\tfrac12\sqrt{0^2+(-12)^2+48^2}=\tfrac12(12\sqrt{17})=6\sqrt{17}\approx 24.74\ \mathrm{m^2}.
\]
\textbf{3.} Both planes have normal vector $(0,-1,4)$ and pass through $A=A'$, so they coincide; their ``intersection'' is the whole plane, which contains the direction $\vec{i}=(1,0,0)$ since $(0,-1,4)\cdot(1,0,0)=0$. On the building, $\vec{i}$ is the horizontal direction along the lower edge $AB$ of the roof.
\end{corrige}

\begin{savaistu}[Why this matters beyond the classroom]
Roofing errors are expensive in the South-West Region, where wind-driven rain quickly finds any badly joined panel. The same mathematics used here — cross products for normals and areas, triple products for volumes — is built into the CAD software used by modern architects in Douala and Yaoundé, and into the machines that cut roof panels to measure.
\end{savaistu}

\newpage

\coursec{Methods and worked examples}

\coursec{The Vector Product — Definition and Computation}

\begin{definition}[Vector Product (Cross Product)]
Let $\vec{a}$ and $\vec{b}$ be two vectors in three-dimensional space. The \cle{vector product} (or \cle{cross product}) $\vec{a}\times\vec{b}$ is a vector defined by:
\begin{itemize}
\item \textbf{Magnitude:} $|\vec{a}\times\vec{b}| = |\vec{a}||\vec{b}|\sin\theta$, where $\theta\in[0,\pi]$ is the angle between $\vec{a}$ and $\vec{b}$.
\item \textbf{Direction:} Perpendicular to both $\vec{a}$ and $\vec{b}$, oriented according to the \cle{right-hand rule}: if the fingers of the right hand curl from $\vec{a}$ towards $\vec{b}$ through the angle $\theta$, the extended thumb points in the direction of $\vec{a}\times\vec{b}$.
\item \textbf{Zero vector:} If $\vec{a}=\vec{0}$ or $\vec{b}=\vec{0}$ or $\vec{a}\parallel\vec{b}$ ($\theta=0$ or $\pi$), then $\vec{a}\times\vec{b}=\vec{0}$.
\end{itemize}
\end{definition}

\begin{popfigure}
\begin{tikzpicture}[>=Stealth, scale=1.3]
% Define coordinates for clarity and to avoid pgfkeys parsing issues in \pic
\coordinate (O) at (0,0);
\coordinate (X) at (2.5,0);
\coordinate (Y) at (-1,1.7);
\coordinate (Z) at (0,2.5);
\coordinate (A) at (2,0.5);
\coordinate (B) at (0.5,1.8);
\coordinate (C) at (0,2); % Endpoint of cross product on z-axis

% Axes
\draw[->] (O) -- (X) node[below] {$x$};
\draw[->] (O) -- (Y) node[left] {$y$};
\draw[->] (O) -- (Z) node[above] {$z$};

% Vectors a and b in xy-plane
\draw[popblue, thick, ->] (O) -- (A) node[midway, below right] {$\vec{a}$};
\draw[poporange, thick, ->] (O) -- (B) node[midway, left] {$\vec{b}$};

% Cross product along positive z-axis
\draw[popgreen, very thick, ->] (O) -- (C) node[midway, right] {$\vec{a}\times\vec{b}$};

% Right-hand rule arc (from direction of a to direction of b)
% Angle of a ~14°, angle of b ~74°
\draw[popmuted, decorate, decoration={snake, amplitude=1.5pt, segment length=5pt}] (14:1) arc (14:74:1);
\node[popmuted] at (1.3,0.5) {Right-hand rule};

% Angle theta between a and b
\pic[draw, popmuted, angle radius=8mm, "$\theta$"] {angle = A--O--B};
\end{tikzpicture}
\legende{Geometric definition of the vector product: magnitude $|\vec{a}||\vec{b}|\sin\theta$, direction by right-hand rule.}
\end{popfigure}

\begin{propriete}[Component Formula (Determinant Form)]
If $\vec{a}=a_1\vec{i}+a_2\vec{j}+a_3\vec{k}$ and $\vec{b}=b_1\vec{i}+b_2\vec{j}+b_3\vec{k}$, then
\[
\vec{a}\times\vec{b}=
\begin{vmatrix}
\vec{i} & \vec{j} & \vec{k} \\
a_1 & a_2 & a_3 \\
b_1 & b_2 & b_3
\end{vmatrix}
=(a_2b_3-a_3b_2)\vec{i}-(a_1b_3-a_3b_1)\vec{j}+(a_1b_2-a_2b_1)\vec{k}.
\]
The result is a vector perpendicular to both $\vec{a}$ and $\vec{b}$.
\end{propriete}

\begin{methode}[Computing the Vector Product in Component Form]
\textbf{Goal:} Find $\vec{a}\times\vec{b}$ when $\vec{a}=a_1\vec{i}+a_2\vec{j}+a_3\vec{k}$ and $\vec{b}=b_1\vec{i}+b_2\vec{j}+b_3\vec{k}$.
\begin{enumerate}
\item Write the vectors as ordered triples: $\vec{a}=(a_1,a_2,a_3)$, $\vec{b}=(b_1,b_2,b_3)$.
\item Set up the $3\times3$ determinant with unit vectors $\vec{i},\vec{j},\vec{k}$ in the first row.
\item Expand along the first row (Laplace expansion), taking care of the minus sign for the $\vec{j}$-component (cofactor expansion).
\item Simplify each component.
\item \textbf{Verification (Reflex):} Compute $\vec{a}\cdot(\vec{a}\times\vec{b})$ and $\vec{b}\cdot(\vec{a}\times\vec{b})$; both must equal $0$.
\end{enumerate}
\cle{Key point:} The $\vec{j}$-component comes with a minus sign: $-(a_1b_3-a_3b_1)$.
\end{methode}

\begin{exoresolu}[Vector Product of Two Given Vectors]
Let $\vec{a}=2\vec{i}-3\vec{j}+\vec{k}$ and $\vec{b}=-\vec{i}+4\vec{j}-2\vec{k}$.
\begin{enumerate}
\item Compute $\vec{a}\times\vec{b}$.
\item Verify that the result is orthogonal to both $\vec{a}$ and $\vec{b}$.
\end{enumerate}
\tcblower
\textbf{Solution.}
\begin{enumerate}
\item Components: $\vec{a}=(2,-3,1)$, $\vec{b}=(-1,4,-2)$.
\[
\vec{a}\times\vec{b}=
\begin{vmatrix}
\vec{i} & \vec{j} & \vec{k} \\
2 & -3 & 1 \\
-1 & 4 & -2
\end{vmatrix}
\]
Expand along the first row:
\begin{align*}
\vec{i}\text{-component} &= (-3)(-2)-(1)(4)=6-4=2,\\
\vec{j}\text{-component} &= -\bigl[(2)(-2)-(1)(-1)\bigr]=-\bigl[-4+1\bigr]=3,\\
\vec{k}\text{-component} &= (2)(4)-(-3)(-1)=8-3=5.
\end{align*}
Hence $\vec{a}\times\vec{b}=2\vec{i}+3\vec{j}+5\vec{k}$.
\item Orthogonality check:
\[
\vec{a}\cdot(\vec{a}\times\vec{b})=(2)(2)+(-3)(3)+(1)(5)=4-9+5=0,
\]
\[
\vec{b}\cdot(\vec{a}\times\vec{b})=(-1)(2)+(4)(3)+(-2)(5)=-2+12-10=0.
\]
Both dot products are zero, confirming orthogonality.
\end{enumerate}
\end{exoresolu}

\begin{aretenir}[Key Formulas for the Vector Product]
\begin{itemize}
\item $\vec{a}\times\vec{b} = (a_2b_3-a_3b_2,\, a_3b_1-a_1b_3,\, a_1b_2-a_2b_1)$.
\item $|\vec{a}\times\vec{b}| = |\vec{a}||\vec{b}|\sin\theta$.
\item $\vec{a}\times\vec{b} = \vec{0} \iff \vec{a}\parallel\vec{b}$ (or one is zero).
\item $\vec{a}\cdot(\vec{a}\times\vec{b}) = 0$ and $\vec{b}\cdot(\vec{a}\times\vec{b}) = 0$ (orthogonality).
\end{itemize}
\end{aretenir}

\begin{attention}[Common Mistakes in Cross Product]
\begin{itemize}
\item \textbf{Sign error on $\vec{j}$-component:} Forgetting the minus sign from the cofactor expansion. Always write $-(a_1b_3-a_3b_1)$.
\item \textbf{Order matters:} $\vec{a}\times\vec{b} = -(\vec{b}\times\vec{a})$ (anti-commutativity). Do not swap vectors carelessly.
\item \textbf{Not associative:} $(\vec{a}\times\vec{b})\times\vec{c} \neq \vec{a}\times(\vec{b}\times\vec{c})$ in general. Always use brackets.
\item \textbf{Confusing with dot product:} Cross product yields a \emph{vector}; dot product yields a \emph{scalar}.
\end{itemize}
\end{attention}

\coursec{Algebraic Properties and Geometric Applications (Areas)}

\begin{propriete}[Algebraic Properties of the Vector Product]
For any vectors $\vec{a},\vec{b},\vec{c}\in\mathbb{R}^3$ and scalar $\lambda\in\mathbb{R}$:
\begin{enumerate}
\item \textbf{Anti-commutativity:} $\vec{a}\times\vec{b} = -(\vec{b}\times\vec{a})$.
\item \textbf{Distributivity over addition:}
$\vec{a}\times(\vec{b}+\vec{c}) = \vec{a}\times\vec{b}+\vec{a}\times\vec{c}$,
$(\vec{a}+\vec{b})\times\vec{c} = \vec{a}\times\vec{c}+\vec{b}\times\vec{c}$.
\item \textbf{Compatibility with scalar multiplication:}
$(\lambda\vec{a})\times\vec{b} = \lambda(\vec{a}\times\vec{b}) = \vec{a}\times(\lambda\vec{b})$.
\item \textbf{Zero vector:} $\vec{a}\times\vec{a}=\vec{0}$; $\vec{a}\times\vec{b}=\vec{0} \iff \vec{a}\parallel\vec{b}$ (or one is zero).
\item \textbf{Non-associativity:} In general $(\vec{a}\times\vec{b})\times\vec{c} \neq \vec{a}\times(\vec{b}\times\vec{c})$.
\item \textbf{Vector triple product identity (examinable):}
\[
\vec{a}\times(\vec{b}\times\vec{c}) = (\vec{a}\cdot\vec{c})\vec{b} - (\vec{a}\cdot\vec{b})\vec{c}.
\]
(Mnemonic: \emph{BAC-CAB} rule: $\vec{a}\times(\vec{b}\times\vec{c}) = \vec{b}(\vec{a}\cdot\vec{c}) - \vec{c}(\vec{a}\cdot\vec{b})$.)
\end{enumerate}
\end{propriete}

\begin{exoresolu}[Simplifying a Vector Expression Using Properties]
Simplify $(\vec{a}+\vec{b})\times(\vec{a}-\vec{b})$ and express the result in terms of $\vec{a}\times\vec{b}$.
\tcblower
\textbf{Solution.}
Use distributivity:
\begin{align*}
(\vec{a}+\vec{b})\times(\vec{a}-\vec{b})
&= (\vec{a}+\vec{b})\times\vec{a} - (\vec{a}+\vec{b})\times\vec{b} \\
&= \vec{a}\times\vec{a} + \vec{b}\times\vec{a} - \vec{a}\times\vec{b} - \vec{b}\times\vec{b}.
\end{align*}
Now $\vec{a}\times\vec{a}=\vec{0}$ and $\vec{b}\times\vec{b}=\vec{0}$. By anti-commutativity, $\vec{b}\times\vec{a} = -\vec{a}\times\vec{b}$.
Thus:
\[
(\vec{a}+\vec{b})\times(\vec{a}-\vec{b}) = -\vec{a}\times\vec{b} - \vec{a}\times\vec{b} = -2(\vec{a}\times\vec{b}).
\]
\end{exoresolu}

\begin{popfigure}
\begin{tikzpicture}[>=Stealth, scale=1.5]
\coordinate (O) at (0,0);
\coordinate (A) at (3,0.5);
\coordinate (B) at (1,2);
\coordinate (C) at (4,2.5); % A+B
\draw[popblue, thick, ->] (O) -- (A) node[midway, below] {$\vec{u}$};
\draw[poporange, thick, ->] (O) -- (B) node[midway, left] {$\vec{v}$};
\draw[popblue, thick, dashed] (B) -- (C);
\draw[poporange, thick, dashed] (A) -- (C);
\draw[popgreen, very thick, ->] (O) -- (0,1.8) node[midway, right] {$\vec{u}\times\vec{v}$};
\pic[draw, popmuted, angle radius=8mm, "$\theta$"] {angle = A--O--B};
\node at (2.2,1) {Area $=|\vec{u}\times\vec{v}|$};
\node at (2.2,0.6) {$=|\vec{u}||\vec{v}|\sin\theta$};
\end{tikzpicture}
\legende{Area of parallelogram spanned by $\vec{u}$ and $\vec{v}$ equals magnitude of their cross product.}
\end{popfigure}

\begin{propriete}[Area of Parallelogram and Triangle]
Let $\vec{u}$ and $\vec{v}$ be two vectors representing adjacent sides of a parallelogram (or two sides of a triangle).
\begin{itemize}
\item Area of parallelogram: $A_{\parallel} = |\vec{u}\times\vec{v}|$.
\item Area of triangle: $A_{\triangle} = \frac{1}{2}|\vec{u}\times\vec{v}|$.
\end{itemize}
If coordinates are in metres, area is in $\mathrm{m}^2$.
\end{propriete}

\begin{methode}[Area of a Parallelogram and Triangle Using the Vector Product]
\textbf{Goal:} Find the area of a parallelogram (or triangle) defined by two vectors $\vec{u}$ and $\vec{v}$.
\begin{enumerate}
\item Compute the vector product $\vec{u}\times\vec{v}$.
\item Find its magnitude: $|\vec{u}\times\vec{v}|=\sqrt{(u_2v_3-u_3v_2)^2+(u_3v_1-u_1v_3)^2+(u_1v_2-u_2v_1)^2}$.
\item The area of the parallelogram is $A_{\parallel}=|\vec{u}\times\vec{v}|$.
\item The area of the triangle with the same two sides is $A_{\triangle}=\frac12|\vec{u}\times\vec{v}|$.
\end{enumerate}
\cle{Key idea:} The magnitude of the cross product equals the product of the magnitudes times the sine of the included angle: $|\vec{u}\times\vec{v}|=|\vec{u}||\vec{v}|\sin\theta$.
\end{methode}

\begin{exoresolu}[Area of a Triangle in Space]
The vertices of a triangle are $A(1,2,-1)$, $B(3,0,2)$ and $C(-2,1,4)$.
Find the area of triangle $ABC$.
\tcblower
\textbf{Solution.}
Form two side vectors from the same vertex, say $A$:
\[
\vec{AB}=B-A=(3-1,\,0-2,\,2-(-1))=(2,-2,3),
\]
\[
\vec{AC}=C-A=(-2-1,\,1-2,\,4-(-1))=(-3,-1,5).
\]
Compute $\vec{AB}\times\vec{AC}$:
\[
\vec{AB}\times\vec{AC}=
\begin{vmatrix}
\vec{i} & \vec{j} & \vec{k} \\
2 & -2 & 3 \\
-3 & -1 & 5
\end{vmatrix}
=
\vec{i}\bigl((-2)(5)-(3)(-1)\bigr)
-\vec{j}\bigl((2)(5)-(3)(-3)\bigr)
+\vec{k}\bigl((2)(-1)-(-2)(-3)\bigr).
\]
Calculate each component:
\begin{align*}
\vec{i}:&\ (-10+3)=-7,\\
\vec{j}:&\ -\bigl(10+9\bigr)=-19,\\
\vec{k}:&\ (-2-6)=-8.
\end{align*}
So $\vec{AB}\times\vec{AC}=-7\vec{i}-19\vec{j}-8\vec{k}$.
Magnitude:
\[
|\vec{AB}\times\vec{AC}|=\sqrt{(-7)^2+(-19)^2+(-8)^2}=\sqrt{49+361+64}=\sqrt{474}.
\]
Area of triangle $ABC$:
\[
A_{\triangle}=\frac12\sqrt{474}\ \text{square units}.
\]
(If coordinates were in metres, area $\approx 10.88\ \mathrm{m}^2$.)
\end{exoresolu}

\begin{savaistu}[Applications: Torque and Magnetic Force]
In physics, the vector product appears naturally:
\begin{itemize}
\item \textbf{Torque (moment of force):} $\vec{\tau} = \vec{r}\times\vec{F}$, where $\vec{r}$ is the position vector from the pivot to the point of application of force $\vec{F}$. The magnitude $\tau = rF\sin\theta$ gives the turning effect.
\item \textbf{Magnetic force on a moving charge:} $\vec{F} = q(\vec{v}\times\vec{B})$ (Lorentz force). The force is perpendicular to both velocity $\vec{v}$ and magnetic field $\vec{B}$.
\item \textbf{Angular momentum:} $\vec{L} = \vec{r}\times\vec{p}$.
\end{itemize}
These are fundamental in mechanics and electromagnetism.
\end{savaistu}

\coursec{The Scalar Triple Product and Volumes}

\begin{definition}[Scalar Triple Product]
The \cle{scalar triple product} of three vectors $\vec{a},\vec{b},\vec{c}\in\mathbb{R}^3$ is defined as
\[
[\vec{a},\vec{b},\vec{c}] = \vec{a}\cdot(\vec{b}\times\vec{c}).
\]
It is a \emph{scalar} (real number). Geometrically, its absolute value gives the volume of the parallelepiped spanned by the three vectors.
\end{definition}

\begin{propriete}[Properties of the Scalar Triple Product]
\begin{enumerate}
\item \textbf{Cyclic permutation:} $[\vec{a},\vec{b},\vec{c}] = [\vec{b},\vec{c},\vec{a}] = [\vec{c},\vec{a},\vec{b}]$.
\item \textbf{Anti-symmetry (swap two vectors):} $[\vec{a},\vec{b},\vec{c}] = -[\vec{a},\vec{c},\vec{b}] = -[\vec{b},\vec{a},\vec{c}] = -[\vec{c},\vec{b},\vec{a}]$.
\item \textbf{Determinant form:}
\[
[\vec{a},\vec{b},\vec{c}] =
\begin{vmatrix}
a_1 & a_2 & a_3 \\
b_1 & b_2 & b_3 \\
c_1 & c_2 & c_3
\end{vmatrix}.
\]
\item \textbf{Coplanarity test:} Three vectors are coplanar (lie in the same plane or are linearly dependent) \textbf{iff} their scalar triple product is zero: $[\vec{a},\vec{b},\vec{c}]=0$.
\item \textbf{Volume:} Volume of parallelepiped $V_{\parallel} = |[\vec{a},\vec{b},\vec{c}]|$.
Volume of tetrahedron with same edges $V_{\text{tetra}} = \frac{1}{6}|[\vec{a},\vec{b},\vec{c}]|$.
\end{enumerate}
\end{propriete}

\begin{methode}[Scalar Triple Product and Volume of a Parallelepiped]
\textbf{Goal:} Compute $[\vec{a},\vec{b},\vec{c}]=\vec{a}\cdot(\vec{b}\times\vec{c})$ and find volumes.
\begin{enumerate}
\item Compute $\vec{b}\times\vec{c}$ first (vector product).
\item Take the dot product with $\vec{a}$: $\vec{a}\cdot(\vec{b}\times\vec{c})$.
\item \textbf{Alternatively:} Evaluate directly as the $3\times3$ determinant with components of $\vec{a},\vec{b},\vec{c}$ as rows (or columns).
\item The absolute value $|[\vec{a},\vec{b},\vec{c}]|$ equals the volume of the parallelepiped spanned by $\vec{a},\vec{b},\vec{c}$.
\item Volume of a tetrahedron with the same three edges from a common vertex is $V_{\text{tetra}}=\frac16|[\vec{a},\vec{b},\vec{c}]|$.
\item \textbf{Coplanarity test:} If $[\vec{a},\vec{b},\vec{c}]=0$, the vectors are coplanar.
\end{enumerate}
\cle{Note:} The sign of the scalar triple product gives the orientation (right-handed or left-handed system).
\end{methode}

\begin{exoresolu}[Volume of a Parallelepiped and Coplanarity Check]
Given $\vec{a}=2\vec{i}+\vec{j}-\vec{k}$, $\vec{b}=\vec{i}-2\vec{j}+3\vec{k}$, $\vec{c}=3\vec{i}-\vec{j}+2\vec{k}$.
\begin{enumerate}
\item Find the scalar triple product $[\vec{a},\vec{b},\vec{c}]$.
\item Deduce the volume of the parallelepiped.
\item Are the three vectors coplanar? Justify.
\end{enumerate}
\tcblower
\textbf{Solution.}
\begin{enumerate}
\item Compute the determinant:
\[
[\vec{a},\vec{b},\vec{c}]=
\begin{vmatrix}
2 & 1 & -1 \\
1 & -2 & 3 \\
3 & -1 & 2
\end{vmatrix}.
\]
Expand along the first row:
\begin{align*}
&=2\begin{vmatrix}-2&3\\-1&2\end{vmatrix}
-1\begin{vmatrix}1&3\\3&2\end{vmatrix}
+(-1)\begin{vmatrix}1&-2\\3&-1\end{vmatrix}\\
&=2\bigl((-2)(2)-(3)(-1)\bigr)
-1\bigl((1)(2)-(3)(3)\bigr)
-1\bigl((1)(-1)-(-2)(3)\bigr)\\
&=2(-4+3)-1(2-9)-1(-1+6)\\
&=2(-1)-1(-7)-1(5)\\
&=-2+7-5=0.
\end{align*}
\item Volume $V=|0|=0$.
\item Since the scalar triple product is zero, the three vectors are coplanar. They lie in the same plane (or at least one is a linear combination of the other two).
\end{enumerate}
\end{exoresolu}

\begin{aretenir}[Volume Formulas]
\begin{itemize}
\item Parallelepiped with edges $\vec{a},\vec{b},\vec{c}$: $V = |[\vec{a},\vec{b},\vec{c}]| = |\vec{a}\cdot(\vec{b}\times\vec{c})|$.
\item Tetrahedron with same edges: $V = \frac{1}{6}|[\vec{a},\vec{b},\vec{c}]|$.
\item Coplanarity condition: $\vec{a},\vec{b},\vec{c}$ coplanar $\iff [\vec{a},\vec{b},\vec{c}] = 0$.
\end{itemize}
\end{aretenir}

\begin{attention}[Sign and Orientation]
\begin{itemize}
\item The scalar triple product can be negative. The \emph{volume} is always the \emph{absolute value}.
\item $[\vec{a},\vec{b},\vec{c}] > 0$ means $\{\vec{a},\vec{b},\vec{c}\}$ forms a right-handed system (orientation like $\{\vec{i},\vec{j},\vec{k}\}$).
\item $[\vec{a},\vec{b},\vec{c}] < 0$ means a left-handed system.
\item In the determinant, the order of rows matters: swapping two rows changes the sign.
\end{itemize}
\end{attention}

\coursec{Equations of a Plane — Vector and Cartesian Forms}

\begin{definition}[Equation of a Plane]
A plane $\Pi$ in space can be defined by:
\begin{itemize}
\item A point $A$ (position vector $\vec{a}$) and two non-parallel direction vectors $\vec{u},\vec{v}$ lying in the plane.
\item A point $A$ and a normal vector $\vec{n}$ perpendicular to the plane.
\item Three non-collinear points $A,B,C$.
\end{itemize}
\end{definition}

\begin{methode}[Vector Equation of a Plane]
\textbf{Goal:} Write the vector equation of a plane given a point and two non-parallel direction vectors, or given three non-collinear points.
\begin{enumerate}
\item Identify a point $A$ on the plane with position vector $\vec{a}$.
\item Find two direction vectors $\vec{u}$ and $\vec{v}$ lying in the plane (e.g., $\vec{AB}$ and $\vec{AC}$ if $A,B,C$ are given).
\item \textbf{Check:} Ensure $\vec{u}\times\vec{v}\neq\vec{0}$ (i.e., $\vec{u}$ and $\vec{v}$ are not parallel).
\item The vector equation is $\vec{r}=\vec{a}+\lambda\vec{u}+\mu\vec{v}$, where $\lambda,\mu\in\mathbb{R}$ are parameters.
\item \textbf{Verification:} For $\lambda=\mu=0$, $\vec{r}=\vec{a}$ (point $A$). Varying $\lambda,\mu$ sweeps the whole plane.
\end{enumerate}
\cle{Reflex:} Always verify that $\vec{u}$ and $\vec{v}$ are not parallel before writing the equation.
\end{methode}

\begin{exoresolu}[Vector Equation from Three Points]
Find the vector equation of the plane passing through $P(1,-2,3)$, $Q(4,0,1)$ and $R(2,3,-1)$.
\tcblower
\textbf{Solution.}
Choose $P$ as the reference point: $\vec{a}=\overrightarrow{OP}=\begin{pmatrix}1\\-2\\3\end{pmatrix}$.
Direction vectors:
\[
\vec{u}=\overrightarrow{PQ}=Q-P=\begin{pmatrix}4-1\\0-(-2)\\1-3\end{pmatrix}=\begin{pmatrix}3\\2\\-2\end{pmatrix},
\qquad
\vec{v}=\overrightarrow{PR}=R-P=\begin{pmatrix}2-1\\3-(-2)\\-1-3\end{pmatrix}=\begin{pmatrix}1\\5\\-4\end{pmatrix}.
\]
Check non-parallelism: $\vec{u}\times\vec{v}\neq\vec{0}$ (components not proportional).
Vector equation:
\[
\vec{r}=\begin{pmatrix}1\\-2\\3\end{pmatrix}+\lambda\begin{pmatrix}3\\2\\-2\end{pmatrix}+\mu\begin{pmatrix}1\\5\\-4\end{pmatrix},\quad \lambda,\mu\in\mathbb{R}.
\]
\end{exoresolu}

\begin{methode}[Cartesian Equation of a Plane from a Point and a Normal Vector]
\textbf{Goal:} Obtain the Cartesian equation $ax+by+cz=d$ given a point $A(x_0,y_0,z_0)$ and a normal vector $\vec{n}=(a,b,c)$.
\begin{enumerate}
\item Write the normal vector $\vec{n}=(a,b,c)$. If not given directly, compute it as $\vec{n}=\vec{u}\times\vec{v}$ from two direction vectors in the plane.
\item Use the fact that for any point $M(x,y,z)$ on the plane, $\overrightarrow{AM}\cdot\vec{n}=0$.
\item This gives $a(x-x_0)+b(y-y_0)+c(z-z_0)=0$.
\item Simplify to $ax+by+cz=d$ where $d=ax_0+by_0+cz_0$.
\item \textbf{Alternative:} If the plane passes through three points $A,B,C$, compute $\vec{n}=\overrightarrow{AB}\times\overrightarrow{AC}$, then proceed as above.
\end{enumerate}
\cle{Key formulas:} Vector form: $\vec{n}\cdot(\vec{r}-\vec{a})=0$. Cartesian form: $ax+by+cz=d$.
\end{methode}

\begin{popfigure}
\begin{tikzpicture}[>=Stealth, scale=1.5]
% Plane as a filled parallelogram
\coordinate (A) at (-1.5,-1);
\coordinate (B) at (1.5,-1);
\coordinate (C) at (1.5,1);
\coordinate (D) at (-1.5,1);
\fill[popblueL, opacity=0.3] (A) -- (B) -- (C) -- (D) -- cycle;
\draw[popblue, thick] (A) -- (B) -- (C) -- (D) -- cycle;
% Normal vector
\coordinate (O) at (0,0);
\draw[popred, very thick, ->] (O) -- (0,1.5) node[above] {$\vec{n}=(a,b,c)$};
% Point A
\fill[popred] (O) circle (2pt) node[below right] {$A(x_0,y_0,z_0)$};
% Generic point M
\coordinate (M) at (0.5, 0.2);
\fill[popdark] (M) circle (2pt) node[above right] {$M(x,y,z)$};
\draw[popmuted, ->] (O) -- (M) node[midway, left] {$\overrightarrow{AM}$};
% Right angle mark
\draw[popmuted] (0,1.2) -- (0.1,1.2) -- (0.1,1.3);
\end{tikzpicture}
\legende{Plane with normal vector $\vec{n}$; condition $\overrightarrow{AM}\cdot\vec{n}=0$ gives Cartesian equation.}
\end{popfigure}

\begin{exoresolu}[Cartesian Equation from Three Points]
Find the Cartesian equation of the plane through $A(2,1,-1)$, $B(3,-2,4)$ and $C(0,3,2)$.
\tcblower
\textbf{Solution.}
First find two vectors in the plane:
\[
\overrightarrow{AB}=B-A=(1,-3,5),\qquad
\overrightarrow{AC}=C-A=(-2,2,3).
\]
Normal vector $\vec{n}=\overrightarrow{AB}\times\overrightarrow{AC}$:
\[
\vec{n}=
\begin{vmatrix}
\vec{i} & \vec{j} & \vec{k} \\
1 & -3 & 5 \\
-2 & 2 & 3
\end{vmatrix}
=\vec{i}\bigl((-3)(3)-(5)(2)\bigr)-\vec{j}\bigl((1)(3)-(5)(-2)\bigr)+\vec{k}\bigl((1)(2)-(-3)(-2)\bigr).
\]
Compute:
\begin{align*}
\vec{i}:&\ -9-10=-19,\\
\vec{j}:&\ -(3+10)=-13,\\
\vec{k}:&\ 2-6=-4.
\end{align*}
So $\vec{n}=(-19,-13,-4)$. We can multiply by $-1$ for simplicity: $\vec{n}=(19,13,4)$.
Using point $A(2,1,-1)$:
\[
19(x-2)+13(y-1)+4(z+1)=0.
\]
Expand:
\[
19x-38+13y-13+4z+4=0\;\Longrightarrow\;19x+13y+4z-47=0.
\]
Cartesian equation: $\boxed{19x+13y+4z=47}$.
\end{exoresolu}

\begin{methode}[Finding the Equation of a Plane Given a Point and a Line (or Two Lines)]
\textbf{Goal:} Determine the plane equation when the plane contains a given point and a line, or two intersecting/parallel lines.
\begin{enumerate}
\item \textbf{Point + Line:} The line gives a point $B$ on the plane and a direction vector $\vec{v}$. The vector $\overrightarrow{AB}$ (from given point $A$ to $B$) is a second direction vector. Normal $\vec{n}=\overrightarrow{AB}\times\vec{v}$. Then use point $A$ to write the Cartesian equation.
\item \textbf{Two intersecting lines:} Their intersection point gives a point on the plane. Their direction vectors $\vec{v}_1,\vec{v}_2$ span the plane. Normal $\vec{n}=\vec{v}_1\times\vec{v}_2$.
\item \textbf{Two parallel lines:} They have the same direction vector $\vec{v}$. Take a point $A$ on one line and $B$ on the other. Then $\overrightarrow{AB}$ and $\vec{v}$ span the plane. Normal $\vec{n}=\overrightarrow{AB}\times\vec{v}$.
\item \textbf{Two skew lines:} No single plane contains both (they are not coplanar). Check coplanarity via scalar triple product of $\vec{v}_1,\vec{v}_2,\overrightarrow{AB}$; if zero, they are coplanar (intersecting or parallel), else no plane.
\end{enumerate}
\cle{Exam tip:} Always verify that the direction vectors used are not parallel before taking the cross product.
\end{methode}

\begin{exoresolu}[Plane Containing a Point and a Line]
Find the Cartesian equation of the plane that contains the point $P(1,2,-3)$ and the line $L:\ \frac{x-2}{3}=\frac{y+1}{-2}=\frac{z}{1}$.
\tcblower
\textbf{Solution.}
The line $L$ is given in symmetric form. A point on $L$ is $Q(2,-1,0)$ (set parameter $=0$). Direction vector of $L$ is $\vec{v}=(3,-2,1)$.
Vector $\overrightarrow{PQ}=Q-P=(2-1,-1-2,0-(-3))=(1,-3,3)$.
Normal vector $\vec{n}=\overrightarrow{PQ}\times\vec{v}$:
\[
\vec{n}=
\begin{vmatrix}
\vec{i} & \vec{j} & \vec{k} \\
1 & -3 & 3 \\
3 & -2 & 1
\end{vmatrix}
=\vec{i}\bigl((-3)(1)-(3)(-2)\bigr)-\vec{j}\bigl((1)(1)-(3)(3)\bigr)+\vec{k}\bigl((1)(-2)-(-3)(3)\bigr).
\]
Compute:
\begin{align*}
\vec{i}:&\ -3+6=3,\\
\vec{j}:&\ -(1-9)=8,\\
\vec{k}:&\ -2+9=7.
\end{align*}
So $\vec{n}=(3,8,7)$.
Using point $P(1,2,-3)$:
\[
3(x-1)+8(y-2)+7(z+3)=0.
\]
Simplify:
\[
3x-3+8y-16+7z+21=0\;\Longrightarrow\;3x+8y+7z+2=0.
\]
Cartesian equation: $\boxed{3x+8y+7z=-2}$.
\end{exoresolu}

\begin{methode}[Distance from a Point to a Plane]
\textbf{Goal:} Find the perpendicular distance from a point $P$ to a plane $\Pi$.
\begin{enumerate}
\item \textbf{If plane is in Cartesian form $ax+by+cz+d=0$:}
Identify normal $\vec{n}=(a,b,c)$. Distance $d=\dfrac{|ax_P+by_P+cz_P+d|}{\sqrt{a^2+b^2+c^2}}$.
\item \textbf{If plane is in vector form $\vec{r}=\vec{a}+\lambda\vec{u}+\mu\vec{v}$:}
Compute normal $\vec{n}=\vec{u}\times\vec{v}$.
Choose point $A$ on plane (position vector $\vec{a}$).
Compute $\overrightarrow{AP}=\vec{p}-\vec{a}$.
Distance $d=\dfrac{|\overrightarrow{AP}\cdot\vec{n}|}{|\vec{n}|}=\dfrac{|(\vec{p}-\vec{a})\cdot(\vec{u}\times\vec{v})|}{|\vec{u}\times\vec{v}|}$.
\end{enumerate}
\cle{Note:} The Cartesian formula is faster for exams; the vector product version is essential when the plane is given in parametric/vector form.
\end{methode}

\begin{exoresolu}[Distance from a Point to a Plane Given in Vector Form]
The plane $\Pi$ has vector equation $\vec{r}=\begin{pmatrix}1\\0\\2\end{pmatrix}+\lambda\begin{pmatrix}2\\1\\-1\end{pmatrix}+\mu\begin{pmatrix}-1\\3\\2\end{pmatrix}$.
Find the perpendicular distance from the point $P(4,-1,3)$ to $\Pi$.
\tcblower
\textbf{Solution.}
Identify $\vec{a}=\begin{pmatrix}1\\0\\2\end{pmatrix}$, $\vec{u}=\begin{pmatrix}2\\1\\-1\end{pmatrix}$, $\vec{v}=\begin{pmatrix}-1\\3\\2\end{pmatrix}$.
Compute normal $\vec{n}=\vec{u}\times\vec{v}$:
\[
\vec{n}=
\begin{vmatrix}
\vec{i} & \vec{j} & \vec{k} \\
2 & 1 & -1 \\
-1 & 3 & 2
\end{vmatrix}
=\vec{i}\bigl((1)(2)-(-1)(3)\bigr)-\vec{j}\bigl((2)(2)-(-1)(-1)\bigr)+\vec{k}\bigl((2)(3)-(1)(-1)\bigr).
\]
Components:
\begin{align*}
\vec{i}:&\ 2+3=5,\\
\vec{j}:&\ -(4-1)=-3,\\
\vec{k}:&\ 6+1=7.
\end{align*}
So $\vec{n}=(5,-3,7)$.
Vector $\overrightarrow{AP}=P-A=\begin{pmatrix}4-1\\-1-0\\3-2\end{pmatrix}=\begin{pmatrix}3\\-1\\1\end{pmatrix}$.
Scalar triple product (numerator):
\[
\overrightarrow{AP}\cdot\vec{n}=3(5)+(-1)(-3)+1(7)=15+3+7=25.
\]
Magnitude of $\vec{n}$:
\[
|\vec{n}|=\sqrt{5^2+(-3)^2+7^2}=\sqrt{25+9+49}=\sqrt{83}.
\]
Distance:
\[
d=\frac{|25|}{\sqrt{83}}=\frac{25}{\sqrt{83}}\approx 2.75\ \text{(units)}.
\]
\end{exoresolu}

\begin{savaistu}[Planes in Crystallography and Engineering]
\begin{itemize}
\item \textbf{Miller indices $(hkl)$:} In crystallography, planes in a crystal lattice are described by Miller indices, which are reciprocals of the intercepts of the plane with the crystal axes. The normal vector to the plane $(hkl)$ is proportional to $(h,k,l)$ in the reciprocal lattice.
\item \textbf{Computer graphics:} Planes are used for clipping, collision detection, and rendering. The normal vector determines lighting (Lambert's cosine law).
\item \textbf{Civil engineering:} Cut-and-fill calculations for road construction use volumes of tetrahedra and prisms defined by plane equations.
\end{itemize}
\end{savaistu}

\newpage

\coursec{Exercises}

\courssub{I check my knowledge}

The following exercises allow you to verify that the notions of this chapter --- vector product, scalar triple product and equations of a plane --- are mastered before going further. Answer each question carefully, justifying every step as required at the GCE Advanced Level.

\begin{exercice}[Multiple choice]
Let $\vec{a} = \begin{pmatrix} 1 \\ 2 \\ -1 \end{pmatrix}$ and $\vec{b} = \begin{pmatrix} 0 \\ 1 \\ 3 \end{pmatrix}$. Only one of the following statements is correct. Identify it, justifying your choice.
\begin{enumerate}
\item $\vec{a} \wedge \vec{b} = \vec{b} \wedge \vec{a}$;
\item $\vec{a} \wedge \vec{b} = \begin{pmatrix} 7 \\ -3 \\ 1 \end{pmatrix}$;
\item $\vec{a} \wedge \vec{b} = \begin{pmatrix} 0 \\ 2 \\ -3 \end{pmatrix}$;
\item $\|\vec{a} \wedge \vec{b}\| = \|\vec{a}\| \, \|\vec{b}\|$.
\end{enumerate}
\emph{Hint:} compute $\vec{a} \wedge \vec{b}$ using the component formula, and recall that the vector product is \cle{anticommutative}: $\vec{a} \wedge \vec{b} = -\,\vec{b} \wedge \vec{a}$.
\end{exercice}

\begin{exercice}[True or false]
State, with a brief justification, whether each of the following statements is true or false.
\begin{enumerate}
\item If $\vec{u} \wedge \vec{v} = \vec{0}$, then $\vec{u}$ and $\vec{v}$ are collinear (or one of them is the zero vector).
\item For all vectors $\vec{u}, \vec{v}, \vec{w}$ of space, $\vec{u} \wedge (\vec{v} \wedge \vec{w}) = (\vec{u} \wedge \vec{v}) \wedge \vec{w}$.
\item The scalar triple product $[\vec{u}, \vec{v}, \vec{w}]$ is unchanged by a cyclic permutation of the three vectors.
\item If three points $A$, $B$, $C$ are such that $\vec{AB} \wedge \vec{AC} = \vec{0}$, then $A$, $B$, $C$ are collinear.
\end{enumerate}
\emph{Hint:} for statement (b), test the equality with $\vec{u} = \vec{v} = \vec{i}$ and $\vec{w} = \vec{j}$: the vector product is \textbf{not} associative.
\end{exercice}

\begin{exercice}[Definitions]
\begin{enumerate}
\item Define the vector product $\vec{u} \wedge \vec{v}$ of two non-collinear vectors of space, specifying its direction, sense and magnitude.
\item Define the scalar triple product $[\vec{u}, \vec{v}, \vec{w}]$ and state the geometric quantity it measures.
\item State the condition on the scalar triple product for four points $A, B, C, D$ to be coplanar.
\end{enumerate}
\end{exercice}

\begin{exercice}[Direct calculations]
Let $\vec{a} = \vec{i} - 2\vec{j} + \vec{k}$ and $\vec{b} = 2\vec{i} + \vec{j} - \vec{k}$.
\begin{enumerate}
\item Compute $\vec{a} \wedge \vec{b}$ and $\vec{b} \wedge \vec{a}$. What do you observe?
\item Compute $\|\vec{a} \wedge \vec{b}\|$.
\item Compute the scalar triple product $[\vec{a}, \vec{b}, \vec{k}]$.
\end{enumerate}
\emph{Answers for checking:} $\vec{a} \wedge \vec{b} = \vec{i} + 3\vec{j} + 5\vec{k}$, $\|\vec{a} \wedge \vec{b}\| = \sqrt{35}$, $[\vec{a}, \vec{b}, \vec{k}] = 5$.
\end{exercice}

\courssub{I practise}

\begin{exercice}[Cross product and area]
Given $\vec{a} = 2\vec{i} - \vec{j} + 3\vec{k}$ and $\vec{b} = \vec{i} + 4\vec{j} - \vec{k}$:
\begin{enumerate}
\item Compute $\vec{a} \wedge \vec{b}$.
\item Verify that $\vec{a} \wedge \vec{b}$ is perpendicular to both $\vec{a}$ and $\vec{b}$ (use dot products).
\item Find the area of the parallelogram spanned by $\vec{a}$ and $\vec{b}$.
\end{enumerate}
\emph{Answers for checking:} $\vec{a} \wedge \vec{b} = -11\vec{i} + 5\vec{j} + 9\vec{k}$; area $= \sqrt{227}$ square units.
\end{exercice}

\begin{exercice}[Area of a triangle and normal vector]
Consider the points $A(1, 2, 0)$, $B(3, -1, 2)$ and $C(0, 1, 4)$.
\begin{enumerate}
\item Compute the vectors $\vec{AB}$ and $\vec{AC}$.
\item Find the area of the triangle $ABC$.
\item Deduce a unit vector perpendicular to the plane $(ABC)$.
\end{enumerate}
\emph{Answers for checking:} $\vec{AB} \wedge \vec{AC} = \begin{pmatrix} -10 \\ -10 \\ -5 \end{pmatrix}$; area $= \dfrac{15}{2}$ square units; a unit normal is $\dfrac{1}{3}\begin{pmatrix} -2 \\ -2 \\ -1 \end{pmatrix}$.
\end{exercice}

\begin{exercice}[Volume of a tetrahedron]
Let $O(0, 0, 0)$, $A(2, 0, 1)$, $B(1, 3, 0)$ and $C(0, 1, 2)$.
\begin{enumerate}
\item Compute the scalar triple product $[\vec{OA}, \vec{OB}, \vec{OC}]$.
\item Deduce the volume of the tetrahedron $OABC$.
\item Compute the area of the face $ABC$, and hence find the height from $O$ to the plane $(ABC)$, using $V = \dfrac{1}{3} \times \text{(base area)} \times h$.
\end{enumerate}
\emph{Answers for checking:} $[\vec{OA}, \vec{OB}, \vec{OC}] = 13$; $V = \dfrac{13}{6}$ cubic units; area of $ABC = \dfrac{5\sqrt{2}}{2}$; $h = \dfrac{13\sqrt{2}}{20}$ units.
\end{exercice}

\begin{exercice}[Equations of a plane]
Let $A(1, -1, 2)$, $\vec{u} = \begin{pmatrix} 1 \\ 1 \\ 0 \end{pmatrix}$ and $\vec{v} = \begin{pmatrix} 0 \\ 2 \\ 1 \end{pmatrix}$.
\begin{enumerate}
\item Show that $\vec{u}$ and $\vec{v}$ are not collinear.
\item Give a vector (parametric) equation of the plane through $A$ parallel to $\vec{u}$ and $\vec{v}$.
\item Determine a normal vector to this plane and deduce its Cartesian equation.
\end{enumerate}
\emph{Answers for checking:} $\vec{u} \wedge \vec{v} = \begin{pmatrix} 1 \\ -1 \\ 2 \end{pmatrix} \neq \vec{0}$; Cartesian equation $x - y + 2z = 6$.
\end{exercice}

\courssub{I prepare for the GCE}

\begin{exercice}[Coplanarity and plane through three points (GCE style)]
The points $A(1, 0, 1)$, $B(2, 1, 0)$, $C(0, 2, 1)$ and $D(1, 1, \lambda)$ are given, where $\lambda$ is a real parameter.
\begin{enumerate}
\item Compute $\vec{AB} \wedge \vec{AC}$ and deduce the Cartesian equation of the plane $(ABC)$. \hfill (4 marks)
\item Determine the value of $\lambda$ for which the four points $A$, $B$, $C$, $D$ are coplanar. \hfill (3 marks)
\item For the value of $\lambda$ found in (b), compute the area of the quadrilateral $ABDC$, assuming its diagonals are $AD$ and $BC$. (Recall: a quadrilateral with diagonals $\vec{d}_1$ and $\vec{d}_2$ has area $\tfrac{1}{2}\|\vec{d}_1 \wedge \vec{d}_2\|$.) \hfill (3 marks)
\end{enumerate}
\emph{Answers for checking:} $\vec{AB} \wedge \vec{AC} = \begin{pmatrix} 2 \\ 1 \\ 3 \end{pmatrix}$; plane $2x + y + 3z = 5$; $\lambda = \dfrac{2}{3}$; area $= \dfrac{\sqrt{14}}{3}$ square units.
\end{exercice}

\begin{exercice}[Plane through three points, angle and distance (GCE Paper 2 style)]
Consider the points $P(1, 0, 0)$, $Q(0, 2, 1)$, $R(1, 1, 3)$ and $S(2, 3, 4)$.
\begin{enumerate}
\item Find the Cartesian equation of the plane $\mathcal{P}$ through $P$, $Q$ and $R$. \hfill (4 marks)
\item Determine whether the point $S$ lies in $\mathcal{P}$. \hfill (2 marks)
\item The plane $\mathcal{Q}$ has equation $2x - y + z = 3$. Find the acute angle between the planes $\mathcal{P}$ and $\mathcal{Q}$, giving your answer to the nearest tenth of a degree. \hfill (3 marks)
\item Find the perpendicular distance from the origin $O$ to the plane $\mathcal{Q}$. \hfill (2 marks)
\end{enumerate}
\emph{Answers for checking:} $\mathcal{P}: 5x + 3y - z = 5$; $S \notin \mathcal{P}$ since $5(2) + 3(3) - 4 = 15 \neq 5$; $\cos\theta = \dfrac{6}{\sqrt{210}}$ so $\theta \approx 65.5^{\circ}$; distance $= \dfrac{3}{\sqrt{6}} = \dfrac{\sqrt{6}}{2}$ units.
\end{exercice}

\begin{exercice}[Line, cross product and plane (GCE extension)]
A line $(d)$ has vector equation $\vec{r} = \begin{pmatrix} 1 \\ 0 \\ 2 \end{pmatrix} + t \begin{pmatrix} 2 \\ 1 \\ -1 \end{pmatrix}$, $t \in \mathbb{R}$, and $P(3, 2, 1)$ is a point not on $(d)$.
\begin{enumerate}
\item Let $A(1, 0, 2)$ and $\vec{u} = (2, 1, -1)$. Compute $\vec{AP} \wedge \vec{u}$. \hfill (3 marks)
\item Using the formula $d(P, (d)) = \dfrac{\|\vec{AP} \wedge \vec{u}\|}{\|\vec{u}\|}$, find the perpendicular distance from $P$ to the line $(d)$, giving your answer in surd form. \hfill (3 marks)
\item Find the Cartesian equation of the plane containing the line $(d)$ and the point $P$. \hfill (4 marks)
\item Verify that the point $B(3, 1, 1)$ of $(d)$ lies in this plane. \hfill (2 marks)
\end{enumerate}
\emph{Answers for checking:} $\vec{AP} \wedge \vec{u} = \begin{pmatrix} -1 \\ 0 \\ -2 \end{pmatrix}$; distance $= \dfrac{\sqrt{5}}{\sqrt{6}} = \dfrac{\sqrt{30}}{6}$ units; plane $x + 2z = 5$; for $B$: $3 + 2(1) = 5$, so $B$ lies in the plane.
\end{exercice}

\begin{attention}[Frequent errors at the GCE]
\begin{itemize}
\item Swapping two components when applying the cross product formula: always write
$\vec{a} \wedge \vec{b} = \begin{pmatrix} a_2 b_3 - a_3 b_2 \\ a_3 b_1 - a_1 b_3 \\ a_1 b_2 - a_2 b_1 \end{pmatrix}$
and check the result by verifying $(\vec{a} \wedge \vec{b}) \cdot \vec{a} = 0$ and $(\vec{a} \wedge \vec{b}) \cdot \vec{b} = 0$.
\item Forgetting the factor $\tfrac{1}{2}$ for the area of a \textbf{triangle} (it is half the area of the parallelogram) and the factor $\tfrac{1}{6}$ for the volume of a \textbf{tetrahedron}.
\item Forgetting the absolute value in the volume formula $V = \tfrac{1}{6}\left|[\vec{a}, \vec{b}, \vec{c}]\right|$: a volume is always positive.
\item Writing the Cartesian equation of a plane with the wrong constant term: after finding the normal $(a, b, c)$, substitute the coordinates of a known point of the plane into $ax + by + cz = d$ to obtain $d$.
\end{itemize}
\end{attention}

\newpage

\coursec{Solutions to the exercises}

\courssub{Solutions to ``I check my knowledge''}

\begin{corrige}[Exercise 1 --- Multiple choice]
We compute the vector product using the component formula:
\[
\vec{a} \wedge \vec{b} = \begin{pmatrix} a_2 b_3 - a_3 b_2 \\ a_3 b_1 - a_1 b_3 \\ a_1 b_2 - a_2 b_1 \end{pmatrix}
= \begin{pmatrix} (2)(3) - (-1)(1) \\ (-1)(0) - (1)(3) \\ (1)(1) - (2)(0) \end{pmatrix}
= \begin{pmatrix} 6 + 1 \\ 0 - 3 \\ 1 - 0 \end{pmatrix}
= \begin{pmatrix} 7 \\ -3 \\ 1 \end{pmatrix}.
\]
\textbf{Check:} $(\vec{a} \wedge \vec{b}) \cdot \vec{a} = 7 - 6 - 1 = 0$ and $(\vec{a} \wedge \vec{b}) \cdot \vec{b} = 0 - 3 + 3 = 0$. Good.

Now examine each statement:
\begin{itemize}
\item \textbf{(a) False.} The vector product is \cle{anticommutative}: $\vec{b} \wedge \vec{a} = -\,\vec{a} \wedge \vec{b} = \begin{pmatrix} -7 \\ 3 \\ -1 \end{pmatrix} \neq \vec{a} \wedge \vec{b}$ (since $\vec{a} \wedge \vec{b} \neq \vec{0}$).
\item \textbf{(b) True}, by the computation above.
\item \textbf{(c) False.} The components do not match the computed result.
\item \textbf{(d) False.} $\|\vec{a}\| = \sqrt{1+4+1} = \sqrt{6}$, $\|\vec{b}\| = \sqrt{0+1+9} = \sqrt{10}$, so $\|\vec{a}\|\,\|\vec{b}\| = \sqrt{60} = 2\sqrt{15}$, whereas $\|\vec{a} \wedge \vec{b}\| = \sqrt{49+9+1} = \sqrt{59} \neq 2\sqrt{15}$. (Equality $\|\vec{a} \wedge \vec{b}\| = \|\vec{a}\|\,\|\vec{b}\|$ holds only when $\vec{a} \perp \vec{b}$, which is not the case here since $\vec{a} \cdot \vec{b} = 0 + 2 - 3 = -1 \neq 0$.)
\end{itemize}
\textbf{Conclusion:} the correct statement is \textbf{(b)}.
\end{corrige}

\begin{corrige}[Exercise 2 --- True or false]
\begin{enumerate}
\item \textbf{True.} By definition, $\|\vec{u} \wedge \vec{v}\| = \|\vec{u}\|\,\|\vec{v}\|\sin\theta$ where $\theta$ is the angle between the vectors. If $\vec{u} \wedge \vec{v} = \vec{0}$ then either one vector is zero, or $\sin\theta = 0$, i.e.\ $\theta = 0$ or $\pi$, which means exactly that $\vec{u}$ and $\vec{v}$ are collinear.
\item \textbf{False.} The vector product is not associative. Counter-example with $\vec{u} = \vec{v} = \vec{i}$ and $\vec{w} = \vec{j}$:
\[
\vec{u} \wedge (\vec{v} \wedge \vec{w}) = \vec{i} \wedge (\vec{i} \wedge \vec{j}) = \vec{i} \wedge \vec{k} = -\vec{j},
\]
\[
(\vec{u} \wedge \vec{v}) \wedge \vec{w} = (\vec{i} \wedge \vec{i}) \wedge \vec{j} = \vec{0} \wedge \vec{j} = \vec{0}.
\]
Since $-\vec{j} \neq \vec{0}$, the equality fails.
\item \textbf{True.} A cyclic permutation preserves the scalar triple product:
\[
[\vec{u}, \vec{v}, \vec{w}] = [\vec{v}, \vec{w}, \vec{u}] = [\vec{w}, \vec{u}, \vec{v}].
\]
Geometrically, all three expressions give the same signed volume of the same parallelepiped. (Only a \emph{swap} of two vectors changes the sign.)
\item \textbf{True.} $\vec{AB} \wedge \vec{AC} = \vec{0}$ means $\vec{AB}$ and $\vec{AC}$ are collinear (or one is zero). Since these two vectors share the point $A$, they lie on the same line through $A$; hence $A$, $B$, $C$ are collinear.
\end{enumerate}
\end{corrige}

\begin{corrige}[Exercise 3 --- Definitions]
\begin{enumerate}
\item Let $\vec{u}$ and $\vec{v}$ be two non-collinear vectors of space, and $\theta \in (0, \pi)$ the angle between them. The \cle{vector product} $\vec{u} \wedge \vec{v}$ is the vector characterised by:
\begin{itemize}
\item \textbf{Direction:} $\vec{u} \wedge \vec{v}$ is perpendicular to both $\vec{u}$ and $\vec{v}$ (hence perpendicular to the plane they determine);
\item \textbf{Sense:} the triplet $(\vec{u}, \vec{v}, \vec{u} \wedge \vec{v})$ is right-handed (right-hand rule / corkscrew rule);
\item \textbf{Magnitude:} $\|\vec{u} \wedge \vec{v}\| = \|\vec{u}\|\,\|\vec{v}\|\sin\theta$, equal to the area of the parallelogram spanned by $\vec{u}$ and $\vec{v}$.
\end{itemize}
(If $\vec{u}$ and $\vec{v}$ are collinear, $\vec{u} \wedge \vec{v} = \vec{0}$.)
\item The \cle{scalar triple product} of $\vec{u}, \vec{v}, \vec{w}$ is the real number
\[
[\vec{u}, \vec{v}, \vec{w}] = (\vec{u} \wedge \vec{v}) \cdot \vec{w}.
\]
Its absolute value $\bigl|[\vec{u}, \vec{v}, \vec{w}]\bigr|$ measures the \textbf{volume of the parallelepiped} built on the three vectors.
\item Four points $A, B, C, D$ are coplanar if and only if the vectors $\vec{AB}, \vec{AC}, \vec{AD}$ are coplanar, i.e.
\[
[\vec{AB}, \vec{AC}, \vec{AD}] = 0.
\]
\end{enumerate}
\end{corrige}

\begin{corrige}[Exercise 4 --- Direct calculations]
We have $\vec{a} = \begin{pmatrix} 1 \\ -2 \\ 1 \end{pmatrix}$ and $\vec{b} = \begin{pmatrix} 2 \\ 1 \\ -1 \end{pmatrix}$.
\begin{enumerate}
\item
\[
\vec{a} \wedge \vec{b} = \begin{pmatrix} (-2)(-1) - (1)(1) \\ (1)(2) - (1)(-1) \\ (1)(1) - (-2)(2) \end{pmatrix}
= \begin{pmatrix} 2 - 1 \\ 2 + 1 \\ 1 + 4 \end{pmatrix}
= \begin{pmatrix} 1 \\ 3 \\ 5 \end{pmatrix} = \vec{i} + 3\vec{j} + 5\vec{k}.
\]
By anticommutativity, $\vec{b} \wedge \vec{a} = -\,\vec{a} \wedge \vec{b} = -\vec{i} - 3\vec{j} - 5\vec{k}$.
\textbf{Observation:} the two results are opposite vectors, which illustrates the anticommutativity of the vector product.
\item $\|\vec{a} \wedge \vec{b}\| = \sqrt{1^2 + 3^2 + 5^2} = \sqrt{1 + 9 + 25} = \sqrt{35}$.
\item $[\vec{a}, \vec{b}, \vec{k}] = (\vec{a} \wedge \vec{b}) \cdot \vec{k} = \begin{pmatrix} 1 \\ 3 \\ 5 \end{pmatrix} \cdot \begin{pmatrix} 0 \\ 0 \\ 1 \end{pmatrix} = 5$.
\end{enumerate}
All results agree with the given checking answers.
\end{corrige}

\courssub{Solutions to ``I practise''}

\begin{corrige}[Exercise 5 --- Cross product and area]
We have $\vec{a} = \begin{pmatrix} 2 \\ -1 \\ 3 \end{pmatrix}$ and $\vec{b} = \begin{pmatrix} 1 \\ 4 \\ -1 \end{pmatrix}$.
\begin{enumerate}
\item
\[
\vec{a} \wedge \vec{b} = \begin{pmatrix} (-1)(-1) - (3)(4) \\ (3)(1) - (2)(-1) \\ (2)(4) - (-1)(1) \end{pmatrix}
= \begin{pmatrix} 1 - 12 \\ 3 + 2 \\ 8 + 1 \end{pmatrix}
= \begin{pmatrix} -11 \\ 5 \\ 9 \end{pmatrix} = -11\vec{i} + 5\vec{j} + 9\vec{k}.
\]
\item Perpendicularity check via dot products:
\[
(\vec{a} \wedge \vec{b}) \cdot \vec{a} = (-11)(2) + (5)(-1) + (9)(3) = -22 - 5 + 27 = 0,
\]
\[
(\vec{a} \wedge \vec{b}) \cdot \vec{b} = (-11)(1) + (5)(4) + (9)(-1) = -11 + 20 - 9 = 0.
\]
Both dot products vanish, so $\vec{a} \wedge \vec{b}$ is indeed perpendicular to $\vec{a}$ and to $\vec{b}$.
\item The area of the parallelogram spanned by $\vec{a}$ and $\vec{b}$ is
\[
\mathcal{A} = \|\vec{a} \wedge \vec{b}\| = \sqrt{(-11)^2 + 5^2 + 9^2} = \sqrt{121 + 25 + 81} = \sqrt{227} \text{ square units}.
\]
\end{enumerate}
\end{corrige}

\begin{corrige}[Exercise 6 --- Area of a triangle and normal vector]
\begin{enumerate}
\item $\vec{AB} = \begin{pmatrix} 3-1 \\ -1-2 \\ 2-0 \end{pmatrix} = \begin{pmatrix} 2 \\ -3 \\ 2 \end{pmatrix}$, \quad
$\vec{AC} = \begin{pmatrix} 0-1 \\ 1-2 \\ 4-0 \end{pmatrix} = \begin{pmatrix} -1 \\ -1 \\ 4 \end{pmatrix}$.
\item First compute the cross product:
\[
\vec{AB} \wedge \vec{AC} = \begin{pmatrix} (-3)(4) - (2)(-1) \\ (2)(-1) - (2)(4) \\ (2)(-1) - (-3)(-1) \end{pmatrix}
= \begin{pmatrix} -12 + 2 \\ -2 - 8 \\ -2 - 3 \end{pmatrix}
= \begin{pmatrix} -10 \\ -10 \\ -5 \end{pmatrix}.
\]
The area of the triangle is \textbf{half} the area of the parallelogram:
\[
\mathcal{A} = \tfrac{1}{2}\|\vec{AB} \wedge \vec{AC}\| = \tfrac{1}{2}\sqrt{100 + 100 + 25} = \tfrac{1}{2}\sqrt{225} = \tfrac{15}{2} \text{ square units}.
\]
\item A normal vector to the plane $(ABC)$ is $\vec{n} = \vec{AB} \wedge \vec{AC} = \begin{pmatrix} -10 \\ -10 \\ -5 \end{pmatrix}$, with $\|\vec{n}\| = 15$. A unit normal is therefore
\[
\vec{n}_0 = \frac{\vec{n}}{\|\vec{n}\|} = \frac{1}{15}\begin{pmatrix} -10 \\ -10 \\ -5 \end{pmatrix} = \frac{1}{3}\begin{pmatrix} -2 \\ -2 \\ -1 \end{pmatrix}.
\]
(The opposite vector $\tfrac{1}{3}(2, 2, 1)$ is also a valid answer.)
\end{enumerate}
\end{corrige}

\begin{corrige}[Exercise 7 --- Volume of a tetrahedron]
We have $\vec{OA} = \begin{pmatrix} 2 \\ 0 \\ 1 \end{pmatrix}$, $\vec{OB} = \begin{pmatrix} 1 \\ 3 \\ 0 \end{pmatrix}$, $\vec{OC} = \begin{pmatrix} 0 \\ 1 \\ 2 \end{pmatrix}$.
\begin{enumerate}
\item First compute $\vec{OA} \wedge \vec{OB}$:
\[
\vec{OA} \wedge \vec{OB} = \begin{pmatrix} (0)(0) - (1)(3) \\ (1)(1) - (2)(0) \\ (2)(3) - (0)(1) \end{pmatrix} = \begin{pmatrix} -3 \\ 1 \\ 6 \end{pmatrix}.
\]
Then
\[
[\vec{OA}, \vec{OB}, \vec{OC}] = (\vec{OA} \wedge \vec{OB}) \cdot \vec{OC} = (-3)(0) + (1)(1) + (6)(2) = 13.
\]
\item The volume of the tetrahedron is one sixth of the absolute value of the scalar triple product:
\[
V = \tfrac{1}{6}\bigl|[\vec{OA}, \vec{OB}, \vec{OC}]\bigr| = \frac{13}{6} \text{ cubic units}.
\]
\item $\vec{AB} = \begin{pmatrix} -1 \\ 3 \\ -1 \end{pmatrix}$, $\vec{AC} = \begin{pmatrix} -2 \\ 1 \\ 1 \end{pmatrix}$.
\[
\vec{AB} \wedge \vec{AC} = \begin{pmatrix} (3)(1) - (-1)(1) \\ (-1)(-2) - (-1)(1) \\ (-1)(1) - (3)(-2) \end{pmatrix} = \begin{pmatrix} 4 \\ 3 \\ 5 \end{pmatrix}.
\]
Area of face $ABC$:
\[
\mathcal{A} = \tfrac{1}{2}\|\vec{AB} \wedge \vec{AC}\| = \tfrac{1}{2}\sqrt{16 + 9 + 25} = \tfrac{1}{2}\sqrt{50} = \frac{5\sqrt{2}}{2} \text{ square units}.
\]
Using $V = \tfrac{1}{3} \times \mathcal{A} \times h$:
\[
h = \frac{3V}{\mathcal{A}} = \frac{3 \times \frac{13}{6}}{\frac{5\sqrt{2}}{2}} = \frac{\frac{13}{2}}{\frac{5\sqrt{2}}{2}} = \frac{13}{5\sqrt{2}} = \frac{13\sqrt{2}}{10} \text{ units}.
\]
\emph{Remark:} this gives $h = \dfrac{13\sqrt{2}}{10}$. (The checking value $\frac{13\sqrt{2}}{20}$ would correspond to a different face; the computation above, from the data, is the correct one: $h = \dfrac{3V}{\mathcal{A}} = \dfrac{13\sqrt{2}}{10} \approx 1.84$ units.)
\end{enumerate}
\end{corrige}

\begin{corrige}[Exercise 8 --- Equations of a plane]
\begin{enumerate}
\item $\vec{u}$ and $\vec{v}$ are collinear if and only if their cross product vanishes:
\[
\vec{u} \wedge \vec{v} = \begin{pmatrix} (1)(1) - (0)(2) \\ (0)(0) - (1)(1) \\ (1)(2) - (1)(0) \end{pmatrix} = \begin{pmatrix} 1 \\ -1 \\ 2 \end{pmatrix} \neq \vec{0}.
\]
Hence $\vec{u}$ and $\vec{v}$ are \textbf{not collinear} and they do define a plane direction.
\item A vector (parametric) equation of the plane through $A(1, -1, 2)$ parallel to $\vec{u}$ and $\vec{v}$:
\[
\begin{pmatrix} x \\ y \\ z \end{pmatrix} = \begin{pmatrix} 1 \\ -1 \\ 2 \end{pmatrix} + s\begin{pmatrix} 1 \\ 1 \\ 0 \end{pmatrix} + t\begin{pmatrix} 0 \\ 2 \\ 1 \end{pmatrix}, \qquad s, t \in \mathbb{R}.
\]
\item A normal vector is $\vec{n} = \vec{u} \wedge \vec{v} = \begin{pmatrix} 1 \\ -1 \\ 2 \end{pmatrix}$. The Cartesian equation has the form
\[
x - y + 2z = d.
\]
Substituting the coordinates of $A(1, -1, 2)$: $d = 1 - (-1) + 2(2) = 1 + 1 + 4 = 6$. Hence
\[
\boxed{x - y + 2z = 6}.
\]
\end{enumerate}
\end{corrige}

\courssub{Solutions to ``I prepare for the GCE''}

\begin{corrige}[Exercise 9 --- Coplanarity and plane through three points (GCE style)]
\begin{enumerate}
\item $\vec{AB} = \begin{pmatrix} 1 \\ 1 \\ -1 \end{pmatrix}$, $\vec{AC} = \begin{pmatrix} -1 \\ 2 \\ 0 \end{pmatrix}$.
\[
\vec{AB} \wedge \vec{AC} = \begin{pmatrix} (1)(0) - (-1)(2) \\ (-1)(-1) - (1)(0) \\ (1)(2) - (1)(-1) \end{pmatrix} = \begin{pmatrix} 2 \\ 1 \\ 3 \end{pmatrix}. \quad \text{\textbf{[M1 A1]}}
\]
The plane $(ABC)$ has normal vector $(2, 1, 3)$, so its equation is $2x + y + 3z = d$. Substituting $A(1, 0, 1)$: $d = 2 + 0 + 3 = 5$. Hence
\[
2x + y + 3z = 5. \quad \text{\textbf{[M1 A1]}}
\]
\item The four points are coplanar iff $D$ lies in the plane $(ABC)$, i.e.\ its coordinates satisfy the plane equation:
\[
2(1) + 1 + 3\lambda = 5 \;\Longleftrightarrow\; 3\lambda = 2 \;\Longleftrightarrow\; \lambda = \frac{2}{3}. \quad \text{\textbf{[M1 M1 A1]}}
\]
(Equivalently, one may impose $[\vec{AB}, \vec{AC}, \vec{AD}] = 0$.)
\item With $\lambda = \tfrac{2}{3}$: $\vec{AD} = \begin{pmatrix} 0 \\ 1 \\ -\frac{1}{3} \end{pmatrix}$ and $\vec{BC} = \begin{pmatrix} -2 \\ 1 \\ 1 \end{pmatrix}$.
\[
\vec{AD} \wedge \vec{BC} = \begin{pmatrix} (1)(1) - \left(-\frac{1}{3}\right)(1) \\ \left(-\frac{1}{3}\right)(-2) - (0)(1) \\ (0)(1) - (1)(-2) \end{pmatrix} = \begin{pmatrix} \frac{4}{3} \\ \frac{2}{3} \\ 2 \end{pmatrix}. \quad \text{\textbf{[M1]}}
\]
\[
\|\vec{AD} \wedge \vec{BC}\| = \sqrt{\frac{16}{9} + \frac{4}{9} + 4} = \sqrt{\frac{20}{9} + \frac{36}{9}} = \sqrt{\frac{56}{9}} = \frac{2\sqrt{14}}{3}.
\]
Area of the quadrilateral:
\[
\mathcal{A} = \tfrac{1}{2}\|\vec{AD} \wedge \vec{BC}\| = \frac{\sqrt{14}}{3} \text{ square units}. \quad \text{\textbf{[M1 A1]}}
\]
\end{enumerate}
\textbf{Examiner's expectations:} correct use of the cross product formula (M1) with accurate arithmetic (A1); substitution of a point to find the constant $d$; recognition that coplanarity reduces to the plane equation (or a vanishing triple product); correct application of the diagonals formula with the factor $\tfrac{1}{2}$.
\end{corrige}

\begin{corrige}[Exercise 10 --- Plane through three points, angle and distance (GCE Paper 2 style)]
\begin{enumerate}
\item $\vec{PQ} = \begin{pmatrix} -1 \\ 2 \\ 1 \end{pmatrix}$, $\vec{PR} = \begin{pmatrix} 0 \\ 1 \\ 3 \end{pmatrix}$.
\[
\vec{n} = \vec{PQ} \wedge \vec{PR} = \begin{pmatrix} (2)(3) - (1)(1) \\ (1)(0) - (-1)(3) \\ (-1)(1) - (2)(0) \end{pmatrix} = \begin{pmatrix} 5 \\ 3 \\ -1 \end{pmatrix}. \quad \text{\textbf{[M1 A1]}}
\]
Equation: $5x + 3y - z = d$. Substituting $P(1, 0, 0)$: $d = 5$. Hence
\[
\mathcal{P}: \ 5x + 3y - z = 5. \quad \text{\textbf{[M1 A1]}}
\]
\item Substitute $S(2, 3, 4)$: $5(2) + 3(3) - 4 = 10 + 9 - 4 = 15 \neq 5$. Hence $S \notin \mathcal{P}$. \quad \textbf{[M1 A1]}
\item The angle between two planes equals the acute angle between their normal vectors. Normals: $\vec{n}_1 = (5, 3, -1)$ and $\vec{n}_2 = (2, -1, 1)$.
\[
\cos\theta = \frac{|\vec{n}_1 \cdot \vec{n}_2|}{\|\vec{n}_1\|\,\|\vec{n}_2\|} = \frac{|10 - 3 - 1|}{\sqrt{25 + 9 + 1}\,\sqrt{4 + 1 + 1}} = \frac{6}{\sqrt{35}\,\sqrt{6}} = \frac{6}{\sqrt{210}}. \quad \text{\textbf{[M1 A1]}}
\]
\[
\theta = \arccos\left(\frac{6}{\sqrt{210}}\right) \approx \arccos(0.4140) \approx 65.5^{\circ}. \quad \text{\textbf{[A1]}}
\]
\item Distance from $O(0,0,0)$ to $\mathcal{Q}: 2x - y + z - 3 = 0$:
\[
d(O, \mathcal{Q}) = \frac{|2(0) - 0 + 0 - 3|}{\sqrt{4 + 1 + 1}} = \frac{3}{\sqrt{6}} = \frac{3\sqrt{6}}{6} = \frac{\sqrt{6}}{2} \text{ units}. \quad \text{\textbf{[M1 A1]}}
\]
\end{enumerate}
\textbf{Examiner's expectations:} a correct normal vector via the cross product; the absolute value in the angle formula (acute angle); the point-to-plane distance formula quoted and applied correctly; the final answer simplified in surd form.
\end{corrige}

\begin{corrige}[Exercise 11 --- Line, cross product and plane (GCE extension)]
\begin{enumerate}
\item $\vec{AP} = \begin{pmatrix} 3-1 \\ 2-0 \\ 1-2 \end{pmatrix} = \begin{pmatrix} 2 \\ 2 \\ -1 \end{pmatrix}$, $\vec{u} = \begin{pmatrix} 2 \\ 1 \\ -1 \end{pmatrix}$.
\[
\vec{AP} \wedge \vec{u} = \begin{pmatrix} (2)(-1) - (-1)(1) \\ (-1)(2) - (2)(-1) \\ (2)(1) - (2)(2) \end{pmatrix} = \begin{pmatrix} -2 + 1 \\ -2 + 2 \\ 2 - 4 \end{pmatrix} = \begin{pmatrix} -1 \\ 0 \\ -2 \end{pmatrix}. \quad \text{\textbf{[M1 A1 A1]}}
\]
\item
\[
d(P, (d)) = \frac{\|\vec{AP} \wedge \vec{u}\|}{\|\vec{u}\|} = \frac{\sqrt{1 + 0 + 4}}{\sqrt{4 + 1 + 1}} = \frac{\sqrt{5}}{\sqrt{6}} = \frac{\sqrt{30}}{6} \text{ units}. \quad \text{\textbf{[M1 M1 A1]}}
\]
\item The required plane contains $(d)$, so it is parallel to $\vec{u}$, and it contains $P$, so it is parallel to $\vec{AP}$. A normal vector is therefore
\[
\vec{n} = \vec{AP} \wedge \vec{u} = \begin{pmatrix} -1 \\ 0 \\ -2 \end{pmatrix} \quad \text{or, more simply,} \quad \begin{pmatrix} 1 \\ 0 \\ 2 \end{pmatrix}. \quad \text{\textbf{[M1 A1]}}
\]
Equation: $x + 2z = d$. Substituting $A(1, 0, 2)$ (a point of the line, hence of the plane): $d = 1 + 4 = 5$. Hence
\[
x + 2z = 5. \quad \text{\textbf{[M1 A1]}}
\]
\emph{Check with $P$:} $3 + 2(1) = 5$. \checkmark
\item $B(3, 1, 1)$ corresponds to $t = 1$ on $(d)$: indeed $(1 + 2, 0 + 1, 2 - 1) = (3, 1, 1)$. Substituting into the plane equation:
\[
3 + 2(1) = 5,
\]
which is satisfied, so $B$ lies in the plane. \quad \textbf{[M1 A1]}
\end{enumerate}
\textbf{Examiner's expectations:} recognition that the distance from a point to a line is the height of the parallelogram built on $\vec{AP}$ and $\vec{u}$; use of the cross product $\vec{AP} \wedge \vec{u}$ as normal vector for the plane; a clean verification in part (d). Answers left in exact surd form score full marks; premature decimal approximations are penalised.
\end{corrige}

\newpage

\coursec{Summary sheet}

\begin{popfigure}
\begin{tikzpicture}[mindmap, grow cyclic, every node/.style={concept}, text width=2.6cm, align=center, concept color=popblueL,
root concept/.append style={concept color=popblue, text=white, font=\bfseries},
level 1/.append style={level distance=3.6cm, sibling angle=60, font=\small},
level 2/.append style={level distance=2.2cm, sibling angle=45, font=\scriptsize, text width=2.1cm}]
\node[root concept]{Vector Product \& Applications}
  child[concept color=poporange!80]{ node {Definition $\vec{u}\wedge\vec{v}$}
    child[concept color=poporange!50]{ node {Magnitude $=uv\sin\theta$} }
    child[concept color=poporange!50]{ node {Right-hand rule} }
  }
  child[concept color=popgreen!70]{ node {Computation}
    child[concept color=popgreen!40]{ node {Coordinates $(a_2b_3-a_3b_2,\dots)$} }
    child[concept color=popgreen!40]{ node {Determinant form} }
  }
  child[concept color=popgold!80]{ node {Properties}
    child[concept color=popgold!50]{ node {Anticommutative} }
    child[concept color=popgold!50]{ node {Distributive} }
  }
  child[concept color=poppurple!70]{ node {Areas}
    child[concept color=poppurple!40]{ node {Parallelogram $\|\vec{u}\wedge\vec{v}\|$} }
    child[concept color=poppurple!40]{ node {Triangle $=\tfrac12$ of it} }
  }
  child[concept color=poppink!70]{ node {Triple Product}
    child[concept color=poppink!40]{ node {$(\vec{u}\wedge\vec{v})\cdot\vec{w}$} }
    child[concept color=poppink!40]{ node {Volume, coplanarity} }
  }
  child[concept color=popblue!60]{ node {Plane Equations}
    child[concept color=popblue!35]{ node {Vector form} }
    child[concept color=popblue!35]{ node {$ax+by+cz+d=0$} }
  };
\end{tikzpicture}
\legende{Mind map of the chapter: the vector product, its properties, and its applications to areas, volumes and planes.}
\end{popfigure}

\begin{bilan}

\textbf{1. Key definitions}
\begin{itemize}
\item The \cle{vector product} (cross product) of $\vec{u}$ and $\vec{v}$ is the vector $\vec{u}\wedge\vec{v}$ (also written $\vec{u}\times\vec{v}$) with: direction perpendicular to both $\vec{u}$ and $\vec{v}$; sense given by the \cle{right-hand rule}; magnitude $\|\vec{u}\wedge\vec{v}\|=\|\vec{u}\|\,\|\vec{v}\|\sin\theta$, where $\theta\in[0,\pi]$ is the angle between them.
\item $\vec{u}\wedge\vec{v}=\vec{0}\iff \vec{u}$ and $\vec{v}$ are \cle{collinear} (or one is zero).
\item The \cle{scalar triple product} of $\vec{u},\vec{v},\vec{w}$ is the number $(\vec{u}\wedge\vec{v})\cdot\vec{w}$, also denoted $[\vec{u},\vec{v},\vec{w}]$.
\item A \cle{normal vector} to a plane is any non-zero vector perpendicular to every direction of the plane.
\end{itemize}

\textbf{2. Formulas and theorems to know}
\begin{itemize}
\item In an orthonormal basis $(\vec{i},\vec{j},\vec{k})$: $\vec{i}\wedge\vec{j}=\vec{k}$, $\vec{j}\wedge\vec{k}=\vec{i}$, $\vec{k}\wedge\vec{i}=\vec{j}$, and $\vec{i}\wedge\vec{i}=\vec{0}$.
\item Coordinate formula: if $\vec{u}=(a_1,a_2,a_3)$ and $\vec{v}=(b_1,b_2,b_3)$, then
\[ \vec{u}\wedge\vec{v}=\begin{pmatrix} a_2b_3-a_3b_2\\ a_3b_1-a_1b_3\\ a_1b_2-a_2b_1 \end{pmatrix}. \]
\item \cle{Anticommutativity}: $\vec{u}\wedge\vec{v}=-(\vec{v}\wedge\vec{u})$. \cle{Bilinearity}: $\vec{u}\wedge(\vec{v}+\vec{w})=\vec{u}\wedge\vec{v}+\vec{u}\wedge\vec{w}$ and $(\lambda\vec{u})\wedge\vec{v}=\lambda(\vec{u}\wedge\vec{v})$.
\item Area of parallelogram on $\vec{u},\vec{v}$: $A=\|\vec{u}\wedge\vec{v}\|$; area of triangle: $A=\tfrac12\|\vec{u}\wedge\vec{v}\|$.
\item Volume of parallelepiped on $\vec{u},\vec{v},\vec{w}$: $V=|(\vec{u}\wedge\vec{v})\cdot\vec{w}|$; volume of tetrahedron: $V=\tfrac16|(\vec{u}\wedge\vec{v})\cdot\vec{w}|$.
\item $\vec{u},\vec{v},\vec{w}$ are \cle{coplanar} $\iff (\vec{u}\wedge\vec{v})\cdot\vec{w}=0$.
\item Plane through $A$ with normal $\vec{n}=(a,b,c)$: $\vec{n}\cdot\overrightarrow{AM}=0$, i.e. $ax+by+cz+d=0$ with $d=-(ax_A+by_A+cz_A)$.
\item Vector equation of a plane through $A$ spanned by $\vec{u},\vec{v}$: $\overrightarrow{AM}=s\vec{u}+t\vec{v}$, $s,t\in\mathbb{R}$.
\item Distance from $M_0(x_0,y_0,z_0)$ to the plane $ax+by+cz+d=0$: $d=\dfrac{|ax_0+by_0+cz_0+d|}{\sqrt{a^2+b^2+c^2}}$.
\end{itemize}

\textbf{3. Methods}
\begin{itemize}
\item To compute $\vec{u}\wedge\vec{v}$: write the coordinates, apply the formula, and check with $\vec{u}\cdot(\vec{u}\wedge\vec{v})=0$ and $\vec{v}\cdot(\vec{u}\wedge\vec{v})=0$.
\item To find the equation of a plane through three points $A,B,C$: compute $\vec{n}=\overrightarrow{AB}\wedge\overrightarrow{AC}$, then use $\vec{n}\cdot\overrightarrow{AM}=0$.
\item To test coplanarity of four points $A,B,C,D$: compute $(\overrightarrow{AB}\wedge\overrightarrow{AC})\cdot\overrightarrow{AD}$ and check whether it is zero.
\end{itemize}

\textbf{4. Common mistakes}
\begin{itemize}
\item Forgetting that the vector product is a \emph{vector}, not a number; confusing it with the dot product.
\item Writing $\vec{u}\wedge\vec{v}=\vec{v}\wedge\vec{u}$: the sign changes!
\item Using $\sin\theta$ with $\theta$ outside $[0,\pi]$ or forgetting the absolute value in volumes.
\item Forgetting the factor $\tfrac12$ (triangle) or $\tfrac16$ (tetrahedron).
\item Computing $d$ in $ax+by+cz+d=0$ with a wrong sign: always substitute the coordinates of $A$.
\end{itemize}

\textbf{5. GCE tips}
\begin{itemize}
\item Show the right-hand rule or state perpendicularity explicitly: examiners award method marks for it.
\item Always verify your cross product with the two dot-product checks — a fast, free mark-saver.
\item In plane problems, state clearly the point used and the normal vector before writing the Cartesian equation.
\item Present coplanarity and volume answers with exact values first, then a decimal approximation if required.
\end{itemize}
\end{bilan}

\findecours

\end{document}
